\documentclass[11pt,oneside]{report}

\usepackage[a4paper,margin=1in]{geometry}
\usepackage[T1]{fontenc}
\usepackage[utf8]{inputenc}
\usepackage{lmodern}
\usepackage{microtype}
\usepackage{parskip}
\usepackage{tocloft}
\usepackage{makeidx}
\usepackage{hyperref}

\makeindex

\setlength{\cftchapnumwidth}{3em}
\setlength{\cftsecnumwidth}{4.5em}

\hypersetup{
  colorlinks=true,
  linkcolor=blue,
  urlcolor=blue,
  pdftitle={Human Body Facts},
  pdfauthor={Chaman Singh Verma}
}

\title{\Huge Human Body Facts\\[0.5em]
       \large Answers to Interesting Questions}
\author{Chaman Singh Verma}
\date{}

\begin{document}
\pagenumbering{roman}

\maketitle
\tableofcontents

\chapter*{Preface}
\phantomsection
\addcontentsline{toc}{chapter}{Preface}

Every day we live inside one of the most intricate systems known to science,
yet most of us rarely pause to ask how it actually works. Why do our fingers
wrinkle in water, why do we shiver when cold, how do trillions of cells
coordinate without a central conductor, and why can a smell unlock a decades-old
memory?

This book grew out of a simple curiosity: collecting the questions people
genuinely ask about their own bodies and answering each one with clarity,
scientific honesty, and respect for the reader. Where physiology is well
established, the explanation goes straight to the mechanism. Where scientists
still debate competing hypotheses---especially in evolutionary biology and
consciousness research---that uncertainty is stated openly rather than smoothed
over.

Because many everyday questions touch on health, the text also separates
well-supported biology from popular myths and highlights warning signs that
deserve prompt medical attention. It is written for curious readers, students,
and lifelong learners alike, not as a clinical manual or a substitute for
professional medical care.

My hope is that these pages deepen your appreciation for the quiet, ceaseless
engineering that keeps us alive, aware, and adaptable from moment to moment.

\medskip
\noindent\textit{Chaman Singh Verma}

\clearpage
\pagenumbering{arabic}
\chapter{Biology, Evolution, and Human Variation}

\section{What is life?}\index{DNA}\index{genes}\index{homeostasis}\index{metabolism}\index{proteins}

There is no universally accepted one-sentence definition of life. A useful
scientific description is a bounded, self-maintaining chemical system that
exchanges matter and energy with its surroundings, regulates internal
conditions, responds to changes, and stores heritable information. Many living
systems grow, develop, repair themselves, and participate in lineages that can
reproduce and evolve. No single feature is sufficient in every case, so
biologists use a collection of interacting characteristics rather than one
strict test.

Every living human body is made primarily of cells, together with extracellular
materials (the supportive fluid and protein scaffolding outside cells). Each
cell is enclosed by a protective membrane, obtains and uses chemical energy
(largely via ATP, the cell's main energy-carrying molecule), makes or modifies
molecules, removes waste, responds to signals, and helps regulate local
conditions such as acidity, temperature, and water content. These activities are
aspects of metabolism and homeostasis. The body
is therefore not a static object but a coordinated process maintained by
trillions of interacting cells.

Biological information is stored mainly in DNA in cellular life, although RNA
can serve as genetic material in some viruses. DNA is not simply a list of
instructions: its protein-coding genes and regulatory switches work with cell
signals and developmental processes to influence what cells do. During
reproduction, heritable information is copied with occasional variation. If
inherited variation affects survival or reproduction, natural selection can
change a population over generations. An individual need not reproduce or evolve
for its species or lineage to do so.

Borderline cases show why the definition is useful but not absolute. A fire
uses energy and spreads, but it lacks a heritable biological system that
evolves by natural selection. Crystals can grow, but they do not metabolise or
use genetic information. Viruses contain genetic material and evolve, but they
cannot independently synthesise proteins or reproduce without a host cell; most
biologists therefore classify them as non-cellular infectious agents rather
than independent living organisms. Prions—misfolded proteins that induce normal
proteins to adopt the same abnormal shape—propagate without genes or
metabolism, making them another important contrast. Clinically, these
distinctions matter: antibiotics target living bacterial cells and therefore do
not treat viruses or prions, while medical definitions of death distinguish
irreversible loss of whole-brain or circulatory function from the slower,
staggered death of individual cells.

In short, life is an organised, self-maintaining, information-linked process,
not a special substance. Its boundaries are clearest for cellular organisms
and less clear for entities that possess only some life-like properties.

\section{Why do breasts remain developed when a person is not breastfeeding?}\index{breasts}\index{alveoli}\index{hormones}\index{nipples and areolas}\index{oestrogen}\index{oxytocin}

Breasts remain visible after puberty because puberty creates a lasting adult
structure rather than a temporary reservoir of stored milk. It includes ducts,
connective tissue, blood vessels, fat, and a variable amount of glandular
tissue. During puberty, hormones such as oestrogen stimulate duct growth and
change the surrounding tissue; the amount and distribution of fat and
connective tissue then contribute greatly to breast size and shape. This tissue
remains even when no milk is being made.

Pregnancy produces a further, temporary change. Oestrogen, progesterone,
prolactin, placental hormones, and other signals enlarge the ducts and develop
milk-producing lobules and alveoli. Before delivery, high oestrogen and
progesterone levels help prepare the breast but restrain full milk synthesis.
After delivery, their rapid fall allows prolactin to support milk production,
while oxytocin causes contraction of cells around the alveoli and helps eject
milk. When lactation ends, the milk-producing tissue usually becomes less active
and partly regresses, although some pregnancy-related changes may persist; the
breast does not normally return to a completely undeveloped state.

Breast size and prominence vary widely with genetics, age, body composition,
hormones, medicines, and life events. They can also change in people who do not
undergo typical female puberty, including through hormone treatment or
endocrine disorders. Breast size does not reliably indicate health, fertility,
or ability to produce milk.

From an evolutionary perspective, adult human breasts tend to remain more
prominent outside pregnancy than those of most other primates, though the
reasons remain uncertain; proposed explanations involving sexual selection or
developmental signalling are hypotheses rather than settled facts. Clinically,
a new persistent lump, skin dimpling, bloody discharge, unexplained one-sided
change, or new nipple inversion should be medically assessed.

\section{Why do males have nipples and areolas, and why are they often flatter?}\index{nipples and areolas}\index{ageing}\index{androgens}\index{breasts}\index{embryonic development}\index{hormones}

Nipples and areolas develop early in embryonic life—around the fifth to sixth
week of gestation—along the embryonic mammary ridges, long before puberty
produces major sex-related differences in breast growth. As a result,
most embryos develop these structures regardless of their genetic sex. Later
hormonal changes mainly modify the amount and organisation of the surrounding
breast tissue; they do not need to create the nipple and areola from the
beginning.

During a typical female puberty, oestrogen, progesterone, growth hormone, and
other signals stimulate substantial development of the milk ducts, lobules,
connective tissue, and fat. Typical male puberty produces much less glandular
and fatty breast development. With less surrounding tissue, male nipples and
areolas may appear flatter or less prominent, although projection is also
affected by smooth muscle, temperature, stimulation, and normal genetic
variation. Some males naturally have more prominent nipples, so ``flatter'' is a
general tendency rather than a rule.

Male breasts usually contain nipple tissue and a small system of milk ducts, but
few or no fully developed milk-producing lobules. \emph{Gynaecomastia} means
benign enlargement of male breast glandular tissue. It reflects increased
oestrogen activity relative to androgen activity in the breast and may result
from normal puberty or ageing, medicines, medical conditions, altered hormone
production, or increased tissue sensitivity; sometimes no specific cause is
found. A new firm lump, nipple discharge, skin change, or persistent one-sided
enlargement should be medically assessed, as although uncommon, male breast
cancer can occur.

\section{How do we remember, and how long can we remember?}\index{amygdala}\index{basal ganglia}\index{brain}\index{cerebellum}\index{cerebral cortex}\index{hippocampus}

\subsection*{How memories are formed and retrieved}

Remembering begins when the brain \emph{encodes} information. Attention,
meaning, emotion, and prior knowledge influence what is noticed and how well it
is learned. The hippocampus and nearby medial-temporal structures are especially
important for forming and initially organising many new memories of events and
facts. Over time, long-term retention depends on wider networks in the cerebral
cortex and, depending on the type of memory, on areas such as the basal
ganglia, cerebellum, and amygdala. Learning changes the strength and structure
of connections between neurons, a capacity called \emph{synaptic plasticity}
(commonly studied through mechanisms such as long-term potentiation, or LTP).

Memory is not one system. \emph{Declarative} (explicit) memory—such as
personal events and general knowledge—relies heavily on the hippocampus and
cortex, whereas \emph{procedural} (implicit) skills such as riding a bicycle,
along with emotional learning and working memory, rely on partly different
networks including the basal ganglia, cerebellum, and amygdala. This is why an
injury or disease can impair one kind of memory while leaving another relatively
intact, and why memory strength is not a single physical quantity stored in one
location.

New memories are fragile at first. Sleep, focused practice, repetition, and
meaningful associations can strengthen them through \emph{consolidation},
although consolidation does not make memories immune to later forgetting or
distortion. Later, a memory is \emph{retrieved} when a suitable cue reactivates
parts of the relevant network. Retrieval is a reconstruction influenced by the
original information and by the person's current context, not a literal replay.
Because retrieval can update a memory, recollections are useful but are not
perfect recordings of the past.

\subsection*{How long memories last and why we forget}

There is no single maximum human memory duration. Sensory traces usually last
only fractions of a second to a few seconds, although the exact duration depends
on the sensory system and task. Information held in working memory is generally
maintained for only about 15--30 seconds unless it is rehearsed or actively
refreshed, and the amount that can be held is limited to a few items or chunks
depending on the material and task.
Working memory usually includes manipulating information, whereas ``short-term
memory'' is often used more narrowly for temporary retention; the terms overlap
and are not used identically by every researcher. A well-learned long-term
memory may last for years or a lifetime, although its details can fade, become
difficult to retrieve, or change through later learning.

Forgetting is normal. Memories may weaken when they are not revisited, become
harder to retrieve because of interference from other learning, or be altered
when they are recalled. Age, stress, poor sleep, illness, and brain injury can
all affect memory, as can reversible medical factors such as delirium, vitamin
B12 deficiency, thyroid disease, and medication side effects; occasional lapses
are normal, whereas sudden or progressive memory loss, confusion, or a major
change in daily functioning should be medically assessed.

\section{What is the clinical significance of tonsils?}\index{tonsils}\index{adenoids}\index{antibiotics}\index{breathing}\index{dehydration}\index{fever}

The tonsils are clinically important mainly because they are a frequent site of
throat infections, can enlarge enough to obstruct breathing during sleep, and
are among the most commonly removed lymphoid organs.

\subsection*{Structure and immune role}

The visible tonsils on either side of the throat are the \emph{palatine tonsils};
the adenoids lie higher behind the nose and are not normally visible through the
mouth. Together with the lingual and tubal tonsils, they form \emph{Waldeyer's
ring}. Their location allows immune cells to sample material entering through
the mouth and nose, a surveillance role that is most prominent during
childhood.

\subsection*{Common clinical problems}

The most common clinical problem is \emph{tonsillitis}, inflammation of the
tonsils. It may cause a sore throat, fever, red or swollen tonsils, painful
swallowing, swollen neck glands, or white coating. Viruses cause most cases;
group A streptococcus is an important bacterial cause. The appearance of white
patches alone cannot identify the cause, so a clinician may use examination and,
when appropriate, a rapid strep test or throat culture. Antibiotics treat
susceptible bacterial infections, not viral tonsillitis.

\emph{Tonsil stones} are usually harmless collections of debris and bacteria
that form in tonsillar crypts. They can cause bad breath or irritation, but they
are different from the coating of acute tonsillitis and usually do not justify
tonsillectomy by themselves. Separately, persistently enlarged tonsils can
obstruct the upper airway and contribute to snoring, sleep-disordered
breathing, or obstructive sleep apnoea, particularly in children.

Rarely, infection spreads into the tissue beside a tonsil and forms a
\emph{peritonsillar abscess} (quinsy). Rapidly worsening one-sided throat pain
or swelling, a muffled voice, difficulty opening the mouth, drooling, or an
inability to swallow saliva needs urgent medical assessment. Difficulty
breathing is an emergency.

\subsection*{When tonsillectomy is considered}

\emph{Tonsillectomy} is surgery to remove the palatine tonsils. It is not needed
for every sore throat. Depending on the person's age, symptoms, examination,
and local clinical guidance, it may be considered for repeated, documented, and
disabling episodes of acute tonsillitis; obstructive sleep apnoea or other
significant breathing obstruction; recurrent peritonsillar abscesses; or a
concerning abnormality such as persistent one-sided enlargement. The decision
requires an individual assessment and discussion of alternatives. Surgery can
cause considerable pain and carries risks including dehydration, bleeding, and
anaesthetic complications. Removing the tonsils removes one local site of
immune surveillance, but most people live normally without them because other
lymphoid tissues continue to support immune defence. The expected benefit must
therefore be weighed against surgical risk for the individual patient.

\section{What are visual angle, visual field, and field of view?}\index{vision!visual field}\index{brain}\index{fovea}

\subsection*{Definitions}

The \emph{visual field} is the area detectable while the eyes are held in a
specified position, and clinical visual-field tests usually measure it under
controlled fixation. \emph{Field of view} is a related term whose meaning is
not completely standardised: in human vision it may include the effects of eye
or head movement, while in cameras and optics it usually describes the angular
extent captured by the device. The \emph{visual angle} is different: it is the
angle at the eye between lines drawn to opposite edges of an object. For an
object of width $s$ at distance $d$, its angular size is approximately $s/d$
when the angle is small, or exactly $2\arctan(s/(2d))$ for a centred object.
Thus a larger object or a nearer object usually has a larger visual angle.

\subsection*{Extent and sharpness of the visual field}

Approximate measurements for a healthy adult are about 150--160 degrees
horizontally and about 130 degrees vertically for one eye. The exact values
vary with facial anatomy, eyelids, gaze direction, and the testing method.
Because the two monocular fields overlap, the combined binocular field is
roughly 200 degrees wide and about 130--135 degrees high, though these are
approximate population values rather than fixed limits. The central overlapping
region provides the best binocular depth perception, or stereopsis; other depth
cues also come from motion, perspective, and the relative size of objects.

The whole field is not equally sharp. The fovea, especially its central
foveola, provides the finest detail and much of our highest-acuity colour vision
when we look directly at something. Reading requires this small central region,
so the eyes make quick movements called \emph{saccades} to place successive
parts of a scene on the fovea. Peripheral vision is less detailed but is useful
for detecting movement, changes in light, and objects outside the direct gaze.

\subsection*{Visual angle and clinical warning signs}

The phrase \emph{visual angle} therefore describes apparent angular size, not
the object's actual physical size. A nearby coin and a distant coin may have
the same visual angle if they appear the same size, even though their real sizes
and distances differ. Visual angle is one cue to perceived size, but the brain
also uses distance, context, perspective, and prior knowledge.

Clinically, loss of part of the visual field, a sudden shower of floaters or
flashes, a curtain-like shadow, or sudden loss or distortion of vision requires
prompt medical assessment because it may signal a serious eye problem.

\section{Why do we have hair on our bodies?}\index{hair}\index{androgens}\index{genes}\index{hormones}\index{puberty}\index{skin}

Humans retain body hair because follicles perform modest protective, sensory,
and sun-shielding roles while reflecting our evolutionary history, even though
most human skin is far less densely covered than the fur of other mammals.

\subsection*{Structure and growth cycle}

Body hair grows from tiny structures in the skin called \emph{hair follicles}.
The exposed hair shaft is made mainly of keratinised cells and has no blood
supply or sensation of its own, although movement of the shaft can stimulate
sensory receptors around the follicle. Living cells in the follicle's growth
region produce the shaft. Hair is not equally thick everywhere: fine, short
\emph{vellus} hairs cover much of the body, while thicker \emph{terminal} hairs
grow on the scalp, eyebrows, and eyelashes and, after puberty, often in areas
such as the armpits, face, chest, or pubic region.

Each follicle cycles through four phases: growth (\emph{anagen}, lasting years
on the scalp), transition (\emph{catagen}, a few weeks), rest (\emph{telogen},
a few months), and shedding (\emph{exogen}). Follicles
in different regions spend different lengths of time in each phase, so a hair's
maximum length depends largely on how long it grows before being shed, as well
as on its growth rate and breakage. Body hair is continually renewed, even
though individual hairs do not grow forever.

\subsection*{Functions and evolutionary context}

Hair has several useful functions. Eyelashes and eyebrows help divert sweat,
dust, and some particles away from the eyes. Hair at the openings of the nose
and ears can catch some larger particles, but the nasal lining, mucus, and
airway cilia do most of the deeper filtering and cleaning. Nerve endings around
hair follicles detect movement of the hair and can help the skin sense a light
touch or an insect before it is clearly felt. Scalp hair also provides some
protection from sunlight and can alter heat exchange at the scalp, although its
insulation is limited compared with fur.

In furry mammals, hair is important for insulation and warmth. Humans have much
less dense body hair than most similarly sized mammals. Scientists have proposed
several explanations for this evolutionary pattern, including improved heat loss
from a relatively hairless skin surface, changes in body shape and activity, and
other environmental pressures; no single explanation is established as the full
answer. Human body hair still provides a small amount of insulation, but sweating
and increased blood flow to the skin are more important for cooling. Its
thermoregulatory effect is modest compared with the dense fur of many mammals.

\subsection*{Hormonal regulation and variation}

The amount, thickness, and distribution of hair are influenced by genes, age,
ancestry, and hormones. Androgens, including testosterone and its more potent
derivative dihydrotestosterone, can transform fine vellus hairs into thicker
terminal hairs in some androgen-sensitive regions after puberty. Their effects
are region-specific: the same hormonal pathway can promote facial or body hair
while shortening the growth phase and contributing to scalp-hair
miniaturisation in genetically susceptible people. Normal body hair therefore
varies widely, and not every hair has a distinct, proven adaptive function.
Sudden marked hair growth or rapid hair loss, especially with other symptoms,
can sometimes signal a hormonal or medical problem and may deserve medical
advice.

\section{Why is pubic hair thicker and often curly?}\index{hair}\index{hair!pubic}\index{androgens}\index{genes}\index{hormones}\index{puberty}

Pubic hair changes at puberty—typically beginning in late childhood or early
adolescence (\emph{pubarche}, driven by adrenal and gonadal hormones and
clinically graded using Tanner stages)—because follicles in the pubic region are
responsive to androgens, a group of hormones that includes testosterone. These
hormones transform fine, lightly pigmented \emph{vellus} hairs into darker,
coarser, longer-growing \emph{terminal} hairs. The same general process affects
the armpits and, in many people, the face and chest, although different regions
respond differently to hormones.

Pubic hair is often curly because curl depends on several interacting features:
the shape and angle of the follicle, the cross-sectional shape of the shaft, and
the way the shaft grows and twists. A more circular shaft often grows straighter,
whereas a more oval or asymmetric shaft tends to bend or curl, but this is a
tendency rather than a rule. Genes strongly influence these features, so pubic
hair can naturally be straight, wavy, or tightly curly in different people.

Thicker, curled hair may reduce skin-on-skin friction or provide modest
protection from minor irritation, but the size and importance of these effects
in humans are uncertain. These possible functions do not make pubic hair a
measure of cleanliness or health. Removing it is a personal choice, not a
medical necessity, although shaving or waxing can irritate the skin or cause
ingrown hairs and inflamed follicles (\emph{folliculitis}). A sudden, marked
change in hair growth or loss that is not explained by grooming, especially with
acne, menstrual irregularity, voice changes, or other signs of hormone excess
(such as polycystic ovary syndrome or androgen-producing conditions), may
deserve medical advice.

\section{Why do we burp?}\index{belching (burping)}\index{breathing}\index{carbon dioxide}\index{diaphragm}\index{digestion}\index{muscles}

Burping, also called \emph{belching} or \emph{eructation}, is the release of
gas through the mouth. The gas is usually swallowed air taken in while eating,
drinking, talking, chewing gum, or breathing through the mouth. Fizzy drinks
add dissolved carbon dioxide, which comes out of solution in the stomach and can
make burping more frequent. A belch is different from intestinal gas, which
leaves through the anus, and is mainly a pressure-relief response rather than a
necessary part of digestion.

When gas collects in the stomach, stretching can trigger a brief relaxation of
the lower oesophageal sphincter, the muscular valve between the oesophagus and
stomach. The gas then moves up the oesophagus while the upper oesophageal
sphincter relaxes, allowing it to leave through the mouth. The diaphragm and
abdominal muscles may assist. A normal belch usually relieves pressure without
forcefully emptying stomach contents, unlike vomiting; swallowed air that does
not leave as a belch passes onward into the intestines, where it may eventually
leave as intestinal gas.

Some belches are \emph{supragastric}: air is drawn into the oesophagus and
released again without reaching the stomach. This pattern can occur repeatedly
and may be influenced by learned breathing or swallowing behaviours rather than
by excess stomach gas.

Occasional burping is normal. Frequent or troublesome belching may be
associated with rapid eating, habitual air swallowing (aerophagia),
gastro-oesophageal reflux, indigestion, fizzy drinks, or an individual response
to particular foods. Not all frequent belching is caused by excess gas in the
stomach; supragastric belching and reflux can produce similar symptoms.
Persistent symptoms should not automatically be attributed to food intolerance.

Belching accompanied by persistent abdominal pain, difficulty swallowing,
repeated vomiting, or unexplained weight loss should be medically assessed.
Belching with chest pain, faintness, or breathlessness requires urgent
assessment because the cause may be cardiac, respiratory, or otherwise
unrelated to digestion.

\section{Can we keep our eyes open while sneezing?}\index{sneezing}\index{brainstem}\index{muscles}\index{pain}

Yes, some people can keep their eyes partly or fully open during a sneeze,
although most close their eyelids automatically. Eyelid closure is part of a
coordinated reflex involving sensory pathways, the brainstem, and facial
muscles: irritation in the nasal lining is carried mainly by branches of the
trigeminal nerve to brainstem circuits, which coordinate the respiratory and
facial responses. The precise evolutionary reason for eyelid closure is not
fully established; it may help shield the eyes from droplets and particles,
but closing the eyes is not essential to the mechanics of sneezing.

The reflex is strong, yet it is not an unbreakable rule. In a healthy eye, the
eyeball is held securely in its socket by muscles, connective tissue, and
surrounding structures; an ordinary sneeze does not make it pop out. Pressure
in and around the eye may change briefly, but this is not normally enough to
displace or damage the eyeball. Trying to keep the eyes open is therefore not
generally dangerous, although pain, bleeding, or a sudden change in vision
after sneezing should be assessed.

Sneezing is a useful way for the nose to clear irritants. It is safer to let a
sneeze out into a tissue or the crook of the elbow rather than forcefully
blocking both the nose and mouth. Suppressing a sneeze in this way can sharply
increase pressure and, although uncommon, injure the throat, ears, or other
tissues. The risk comes from trapping the sneeze, not from leaving the eyes
open.

\section{How fast does blood flow?}\index{aorta}\index{arteries}\index{capillaries}\index{carbon dioxide}\index{cardiac output}\index{heart}

``Blood flow'' can refer either to the volume moved through a vessel per unit of
time or to the linear speed of blood cells. These are related but not identical,
and neither is constant throughout the circulation. At rest, the heart pumps
about 5 litres per minute through the circulation; this is the \emph{cardiac
output}. During exercise, cardiac output can rise severalfold so that working
muscles receive more oxygen and nutrients. All numerical values below are
approximate and vary with vessel, location, heartbeat phase, posture, health,
and activity.

Blood generally travels fastest in the large arteries close to the heart. In the
aorta, peak speed during part of a heartbeat can approach about 0.5 metres per
second, while the average is lower. As arteries branch into arterioles and
capillaries, the combined cross-sectional area of all the vessels becomes much
larger, so the average linear speed falls greatly. In capillaries, where oxygen,
carbon dioxide, nutrients, and waste products are exchanged with tissues, it may
be around 1 millimetre per second. In simplified form, volumetric flow equals
average velocity multiplied by total cross-sectional area, so the same overall
flow can move slowly through the enormous combined area of capillaries and then
speed up as vessels join into larger veins. Real flow is pulsatile and variable,
so this relationship describes averages rather than every individual cell.

At rest, a tracer or blood cell may take roughly a minute to pass through the
pulmonary and systemic circuits, but transit time differs among organs and
routes and changes with cardiac output and activity. Blood does not make isolated
trips; it circulates continuously, with different cells following different
paths and travelling at different speeds. The pressure wave felt as a pulse
travels through the arteries faster than the blood itself. Its speed also varies
with arterial elasticity, which is why a pulse at the wrist is detected soon
after the heart contracts.

\section{Why are we ticklish?}\index{ticklishness}\index{brain}\index{cerebellum}\index{mechanoreceptors}\index{skin}\index{touch}

Ticklishness is an interpretation created by the nervous system, thought to
have evolved either to help detect small threats such as parasites on vulnerable
skin or to support juvenile play and social bonding, though neither explanation
is settled. It begins when light, irregular, or unexpected touch activates
mechanoreceptors and other sensory endings in the skin, and the brain combines
those signals with attention, expectation, emotional state, and social context
to produce a crawling sensation, startle response, urge to move away, or
laughter.

Researchers often use two descriptive terms. \emph{Knismesis} refers to the
itchy or crawling sensation caused by very light touch, such as a feather moving
over the skin. \emph{Gargalesis} refers to the stronger response to repeated or
firm touch in sensitive areas and is more likely to produce familiar ticklish
laughter. These terms are useful, but they do not capture every individual form
of ticklishness. Sensitive regions often include the soles of the feet,
underarms, neck, and sides of the torso. Their sensitivity may reflect sensory
innervation, vulnerability to unexpected contact, and learned or social
associations rather than any single cause.

We usually cannot tickle ourselves in the same way. The cerebellum and other
brain systems predict the sensory consequences of our own movements. They use
that prediction to reduce the perceived intensity or surprise of self-produced
touch, a process called sensory attenuation. This prediction involves the
cerebellum and interacting sensory and motor networks rather than one isolated
brain centre. Self-touch may feel more ticklish when it is delayed, altered, or
produced by a device, while another person's touch is harder to predict.

The evolutionary value of ticklishness is not settled. Detecting small movements
on vulnerable skin could help an animal respond to parasites or unexpected
contact, while stronger tickling may have roles in juvenile play and social
bonding. These possibilities are hypotheses rather than established explanations
for every form of human ticklishness.

Ticklish laughter is not always a sign of amusement; people may laugh while
feeling surprised, uncomfortable, or unable to control the response. Depending
on the person and situation, ticklishness can be pleasant, neutral, or
distressing, so consent and social context matter. Sensitivity varies with
mood, attention, expectation, relationship, culture, and individual biology.

\section{How can sleeping position affect different people?}\index{breathing}\index{pain}\index{sleep!apnoea}\index{sleep}\index{stomach}

\subsection*{Adult sleeping positions}

There is no single best sleeping position for every adult. Side sleeping can
reduce snoring or positional obstructive sleep apnoea in some people because
gravity is less likely to let the tongue and soft tissues fall backward. It does
not work equally well for everyone and is not a substitute for assessment or
treatment of sleep apnoea. Keeping the neck in a comfortable, neutral position
and placing a pillow between the knees may reduce strain for some people with
back or hip pain. Either side is usually acceptable; some people with nighttime
reflux find the left side more comfortable because the stomach curves below the
oesophagus in a way that reduces acid pooling at the junction. These are comfort
and symptom management effects, not guarantees that a position will correct the
underlying condition.

Sleeping on the back can be comfortable and may support the spine when the
pillow and mattress keep it in a neutral position. In some people, however,
gravity allows the tongue and soft tissues to narrow the airway, worsening
snoring or obstructive sleep apnoea. For nighttime reflux, raising the head end
of the bed or using a suitable wedge may help more reliably than simply adding
pillows, which can bend the body and increase discomfort. Persistent loud
snoring, witnessed pauses in breathing, gasping during sleep, or severe daytime
sleepiness should be assessed rather than managed by position alone.

Sleeping on the stomach is not automatically dangerous for an adult, but it
often keeps the neck rotated and may increase strain on the neck or lower back.
It may therefore be uncomfortable for people with spinal or neck pain, while
some people find it comfortable. In general, the most suitable adult position
is the one that permits comfortable, uninterrupted sleep without pain or
breathing problems.

\subsection*{Pregnancy}

Pregnancy requires special advice. After 28 weeks, guidance commonly recommends
starting sleep on either side rather than deliberately settling flat on the
back, because the enlarged uterus can compress major blood vessels (notably the
inferior vena cava and aorta, reducing blood return to the heart and placenta). This does not
mean a person must remain in one position all night. Waking on the back is not a
reason to panic; the person can simply turn onto a side. This advice is separate
from adult comfort advice because it concerns pregnancy-related circulation and
safety.

\subsection*{Infants}

Infant sleep recommendations are different and stricter: babies should be
placed on their backs for every sleep on a firm, flat, non-inclined surface with
only a fitted sheet. Pillows, loose blankets, toys, and other soft objects
should be kept out of the sleep area. These recommendations are intended to
reduce sleep-related suffocation and sudden infant death syndrome (SIDS) risks
and should not be extrapolated to adult sleeping positions.

\section{Can stem-cell treatments extend healthy human lifespan?}\index{stem cells}\index{ageing}\index{cancer}\index{DNA}\index{heart}\index{kidneys}

\subsection*{Why stem cells are not a proven anti-ageing therapy}

No stem-cell intervention has convincingly been shown in well-controlled human
studies to extend maximum lifespan or broadly delay ageing in otherwise healthy
people, and routine stem-cell treatment is not an established anti-ageing
therapy. A stem cell can self-renew and can produce one or more specialised
cell types, but different stem cells have very different abilities. They are
not a universal supply of youth. Ageing affects the whole body: DNA and
proteins can become damaged, organs and immune function change, and cells can
lose some of their ability to repair tissue. Replacing one cell type would not
automatically reverse all of these processes.

\subsection*{Established medical uses and risks}

Stem-cell transplantation is already an important treatment in carefully
selected conditions. For example, transplantation of blood-forming
\emph{haematopoietic} stem cells can replace diseased or damaged marrow in some
blood cancers, marrow-failure syndromes, and serious blood or immune disorders.
Some transplants use the patient's own cells; others use cells from a donor.
This is intensive medical treatment, not a yearly wellness procedure. Infection
and graft failure are possible, while infertility and damage to organs such as
the liver, kidneys, lungs, or heart can result particularly from the high-dose
conditioning treatment. With donor cells, immune cells in the graft may attack
the patient's tissues, causing \emph{graft-versus-host disease}.

\subsection*{Research obstacles and unproven clinics}

Researchers are studying ways to repair ageing tissues, restore the function of
older stem cells, and replace or reprogramme damaged cells. Important obstacles
include poor survival or integration of transplanted cells, unwanted immune
reactions, abnormal tissue growth, tumour formation, and uncertain long-term
benefit. Encouraging results in cells or animals do not show that repeated
infusions extend human lifespan or healthspan.

Clinics offering routine anti-ageing stem-cell injections should therefore be
treated with great caution. For most such uses, safety and effectiveness have
not been established in rigorous clinical trials, and regulators have received
reports of serious harms from unapproved products, including infections,
blindness, and tumours. Regulatory approval differs by country; in the United
States, FDA-approved cellular and stem-cell products are limited to specific
medical indications; none are approved as general anti-ageing treatments. A treatment
being advertised as ``regenerative'' or listed in a clinical-trial database does
not by itself prove that it is approved, effective, or appropriate. Legitimate
experimental treatment should be provided through properly regulated research
with informed consent.

\section{Why do we sleep, and why does it occupy roughly one-third of life?}\index{sleep}\index{brain}\index{breathing}\index{circadian rhythm}\index{heart}\index{hormones}

\subsection*{Regulation and bodily maintenance}

Sleep is not a shutdown or a sign that the body has stopped working. It is an
active, highly organised state in which the brain and body perform tasks that
are difficult to coordinate during waking life. Sleep is regulated by at least
two interacting processes: a circadian clock—anchored by the brain's
suprachiasmatic nucleus and signalled partly by evening melatonin release—that
helps time sleep across the day and night, and a sleep pressure (driven in part
by accumulating adenosine) that builds the longer we remain awake. The brain
cycles roughly every 90 minutes through non-REM stages (from light N1 and N2
sleep to deep slow-wave N3 sleep) and REM sleep, each with partly different
patterns of brain activity, muscle tone, heart rate, and breathing.

The body maintains and repairs itself during waking life as well as during
sleep, but sleep alters internal conditions in ways that support growth-hormone
pulses, metabolism, and immune regulation without the demands of active
behaviour.

\subsection*{Brain functions during sleep}

The brain has its own reasons for needing sleep. During sleep, neural circuits
can strengthen useful memories, while other connections may be weakened or
reorganised. Non-REM sleep is especially associated with slow-wave activity and
physical recovery, while REM sleep supports changing patterns of brain activity
and may contribute to learning, emotion, and memory. Neither stage has one
single function, and their roles overlap. Brain-fluid movement may also
contribute to clearing some metabolic waste, but the size and importance of this
effect in living humans remain under study.

\subsection*{Why we spend about a third of life asleep}

Spending roughly one-third of life asleep is a broad human average, not a fixed
requirement. Sleep need changes with age and individual biology; most adults
are generally advised to obtain about 7--9 hours per night, while children and
adolescents usually need more. The body is not ``too inefficient'' to heal
continuously: some neural processing and maintenance are simply better
coordinated during sleep, while other repair continues while awake. Regularly
getting too little or poor-quality sleep can impair attention, mood, immune
function, metabolism, and physical recovery, though why sleep evolved and why
each stage is necessary remain active areas of research.

\section{How do children acquire language, and is there a universal grammar?}\index{language acquisition}\index{brain}\index{hearing}\index{learning}\index{memory}\index{speech}

\subsection*{How children acquire language}

Children acquire language through the interaction of a developing brain,
learning abilities, and a language-rich social environment. In the first months,
infants become sensitive to the sounds and rhythm of the languages around them
and begin taking turns in vocal interaction. Babbling, shared attention,
gesture, and responsive conversation help them learn which sounds, words, and
social signals matter. Typically, infants begin canonical babbling around 6--9
months, speak first meaningful words around 12 months, and combine two-word
phrases around 24 months, supported by developing language networks including
Broca's and Wernicke's areas. They gradually discover likely word boundaries,
connect words with people and events, and infer patterns from the language they
hear. They then use those patterns creatively, producing sentences they have
never heard before. Language development is therefore an active process
involving input, statistical learning, memory, attention, and social
interaction.

The same broad ability supports spoken and signed languages. Children do not
simply memorise a list of sentences; they learn flexible systems that can
generate new ones. A child exposed consistently to more than one language can
acquire more than one, and a child who cannot hear speech can acquire a signed
language through accessible interaction. Because milestones vary, persistent
lack of babbling, poor response to sound, loss of previously learned words, or
no single words by 16--18 months warrants prompt hearing and developmental
assessment. Brain development is especially
adaptable in early childhood, suggesting a sensitive period during which
language input is particularly influential. This is not a single deadline:
language learning continues throughout life, although age of exposure can
affect pronunciation, automaticity, and some aspects of grammar. Outcomes also
depend on the quality and accessibility of input, development, hearing, and the
child's opportunities to communicate.

Language is the ability to represent and share meaning through a structured
system; speech is only one way to express it. Delayed or inaccessible language
input can restrict development, whereas early access to a natural spoken or
signed language supports communication. This does not mean that children need
formal grammar lessons: responsive interaction gives them examples, feedback,
shared attention, and reasons to communicate. All naturally developed human
languages have complex systems for expressing meaning; no language is inherently
primitive or incomplete.

\subsection*{Is there a Universal Grammar?}

\emph{Universal Grammar} (UG) is the name given, especially in Chomskyan
linguistics, to the hypothesis that humans are born with language-specific
structures or constraints that help determine which grammars are learnable. UG
does not mean that every language uses the same words or visible rules;
languages differ greatly in sounds, word order, and grammatical organisation.
It is also different from the less controversial claim that humans have evolved
general biological capacities that make language learning possible.

UG remains a debated theoretical proposal rather than an established result.
Researchers disagree about how much language-specific structure is inherited
and how much language can be explained by domain-general learning, social
interaction, prediction, memory, and brain development. Evidence that languages
share some patterns does not by itself prove UG, and wide variation among
languages also constrains strong universal claims. A cautious conclusion is
that humans clearly inherit biological capacities and biases that support
language learning, but science has not settled whether these amount to the
specialised grammar proposed by UG.

\section{How can skeletal features help estimate biological sex, and what are the limits?}\index{bones}\index{DNA}\index{muscles}\index{pelvis}\index{puberty}

Bones can show average differences associated with sex-related development, a
pattern called \emph{sexual dimorphism}. In unidentified adult remains, the
pelvis usually provides the strongest skeletal evidence. Several pelvic traits
are considered together, including the shape and breadth of the pelvic inlet and
outlet, the angle and form of the pubic region, and the shape of the sacrum.
People with typical female development often have a relatively wider pelvic
canal, while people with typical male development often have a narrower canal
and more pronounced areas of muscle attachment. These are overlapping population
patterns, not separate anatomical categories. The skull, jaw, and some long
bones may also show average differences in size and shape, but they are
generally less reliable when assessed alone.

These features are not absolute. Human bodies vary widely, and the distributions
for males and females overlap substantially. Nutrition, health, ancestry, age,
physical activity, pathology, and individual genetics all influence the
skeleton. Methods are also population-specific: a formula developed from one
reference population may be less accurate for someone from another population.
Sex estimation is generally less reliable in children because many adult
differences develop during puberty, and it becomes harder when remains are
fragmentary or damaged. Forensic specialists therefore combine visual assessment
and measurements from several bones, use validated methods appropriate to the
available evidence and reference population, and report a probabilistic estimate
with uncertainty rather than claiming that one bone can identify a person's sex
with certainty. When available, DNA and other medical or documentary evidence
may provide more direct evidence about an individual's sex-related biology than
skeletal morphology alone.

Bones can also provide clues about age, stature, injuries, and disease. In a
forensic investigation, DNA, dental records, medical history, and other evidence
may be needed to establish identity. A skeletal estimate describes population-
level patterns in sex-related development; it does not determine a person's
gender identity and may not fully capture variation in sex development.

\section{How many egg cells are usually released during a lifetime?}\index{egg cells (oocytes)}\index{fertilisation}\index{hormones}\index{menopause}\index{puberty}\index{speech}

The direct answer is roughly 400--500 oocytes released by ovulation during a
typical reproductive lifetime, although the number varies widely. The phrase
``produce eggs'' is slightly misleading. The ovarian reserve---the
pool of immature egg cells, or oocytes---forms during fetal development. There
is no clinically established way to replenish this reserve in adult humans.
Rough estimates suggest about 1--2 million oocytes at birth and about
300,000--500,000 at puberty, but these are population estimates and vary
considerably between individuals.

During a typical ovulatory cycle, a group of follicles begins to develop and
usually one dominant follicle releases one secondary oocyte. Some cycles are
anovulatory and release none; occasional multiple ovulation can release two or
more and may result in fraternal twins. Most follicles that begin developing
are naturally lost through \emph{atresia}. Across a typical reproductive
lifetime, roughly 400--500 oocytes are ovulated, but this is a rough
population-level estimate, not a fixed biological quota. The number varies with
cycle patterns, pregnancies, breastfeeding, hormonal contraception, and the
age at which menopause occurs. The great majority of oocytes are lost through
atresia without being released.

In everyday speech, the ovulated cell is called an egg. Technically, it is a
secondary oocyte and completes its final meiotic division only if fertilisation
occurs. The number of oocytes is only one part of reproductive potential: egg
quality also changes with age, and ovarian reserve declines at different rates
in different people. Tests such as anti-Müllerian hormone (AMH) and antral
follicle count can help estimate aspects of ovarian reserve and predict response
to some fertility treatments, but they do not count every remaining oocyte,
directly measure egg quality, or reliably predict whether a person will conceive
naturally.

Hormonal contraception can suppress ovulation while it is being used, but it
does not preserve a fixed number of oocytes or stop their age-related decline;
follicles continue to be lost through atresia. Fertility medicines can stimulate
several follicles to mature in one cycle, so more than one oocyte may be
retrieved, but they do not create a new lifetime supply of oocytes.

\section{How can facial and body features influence perceptions of attractiveness?}\index{breasts}\index{hormones}\index{learning}\index{muscles}\index{skin}

Facial and body features can influence judgements of attractiveness, but those
judgements are neither universal nor reliable biological tests. People may notice
features such as facial symmetry, averageness, skin appearance, expression,
body proportions, or the distribution of body fat and muscle, sometimes without
conscious thought. Appearance is also affected by age, genetics, hormones,
health, posture, clothing, and temporary changes such as pregnancy or illness.
No single feature, including breast size or hip shape, directly measures health,
fertility, or reproductive hormones.

Some evolutionary theories propose that preferences for particular features may
have been influenced, in part, by broad cues related to age or physical
condition. These proposals remain hypotheses about population-level tendencies,
not proof that a feature reliably signals fertility or health. Cultural learning,
availability of potential partners, social status, personal experience, and
individual psychology also shape preferences. Even when an average pattern is
found in a study, it cannot explain every person's attraction or predict an
individual's behaviour. Results can also depend on the population sampled, the
way faces or bodies are measured, the surrounding culture, and whether people
are judging photographs or real social interactions. A statistical average is
therefore not a universal standard of beauty.

Attraction is shaped by familiarity, personal experience, media, cultural
ideals, social context, emotional connection, sexual orientation, gender
identity, and individual differences. Preferences vary between people and
cultures, and can change over time. A visual impression is not the same as
attraction, attraction is not the same as intention to act, and neither is
consent. A body feature therefore never reveals a person's personality, sexual
experience, willingness to have sex, or worth as a human being; consent must be
freely communicated by the person concerned.

\chapter{Heart, Blood, and Circulation}

\section{Why is blood red, and why do people have different blood groups?}\index{blood!groups}\index{antibodies}\index{genes}\index{haemoglobin}\index{placenta}\index{platelets}

\subsection*{Why blood is red}

Blood is red mainly because red blood cells are packed with haemoglobin, an
iron-containing protein whose haem groups absorb some wavelengths of light and
reflect others. Oxygenated blood is usually bright red, whereas deoxygenated
blood is dark red; veins can look blue through the skin because of the way light
is absorbed and scattered, not because the blood inside them is blue. Plasma is
pale yellow, but the large number of red cells dominates the colour of whole
blood.

\subsection*{Why blood groups differ and why they matter}

Blood groups are inherited differences in molecules called \emph{antigens},
which are found especially on the surface of red blood cells. The best-known
systems are ABO and Rh, but systems such as Kell, Duffy, Kidd, and MNS can also
matter. In the ABO system, a person's red cells may carry A antigen, B antigen,
both, or neither; the plasma usually contains antibodies against the ABO antigen
that the person lacks. If incompatible red cells are transfused, these
antibodies can bind to them and cause dangerous destruction of the cells.
Before a transfusion, laboratories therefore identify the relevant antigens,
screen for antibodies, and perform compatibility testing for the specific blood
component being given, since plasma and platelet products have additional
compatibility rules beyond red cells.

Blood-group differences arise from inherited gene variants. Mutation, migration,
chance changes in gene frequency, and possible historical selection by infectious
diseases have contributed to their distribution in human populations. These
differences are normal genetic variation; no common blood group is universally
healthiest, and blood groups do not determine personality, intelligence, or
moral character.

Blood-group differences can also matter in pregnancy, although parental blood
groups being different does not automatically cause a problem. For example, a
D-negative pregnant person can become sensitised to D-positive fetal red cells
after exposure. Anti-D antibodies can cross the placenta and cause haemolytic
disease of the fetus and newborn, particularly in a later D-positive pregnancy.
Antibody screening and preventive anti-D immunoglobulin substantially reduce
this risk, although the exact schedule depends on local guidance and the
person's antibody status.

\section{Why does the heart beat continuously?}\index{heart}\index{aorta}\index{arteries}\index{nervous system!autonomic}\index{blood pressure}\index{carbon dioxide}

The heart is a muscular pump that keeps blood moving through the lungs and the
rest of the body. Continuous circulation is necessary to deliver oxygen and
nutrients, remove carbon dioxide and waste, and maintain blood pressure for the
organs. Special pacemaker cells in the sinoatrial (SA) node spontaneously
depolarise and generate electrical impulses. The impulse spreads through the
atria, is delayed briefly at the atrioventricular (AV) node, and then travels
through the bundle branches and Purkinje fibres so the ventricles contract in a
coordinated sequence. The heart therefore has an intrinsic rhythm and does not
need a new conscious command for each beat.

The chambers work in a repeating sequence rather than all contracting together.
The atria receive blood and contract to help fill the ventricles; the
ventricles then eject blood into the pulmonary artery and aorta. Valves open and
close in response to pressure differences and prevent substantial backward
flow. Because the right and left sides pump in series and the vessels remain
filled, circulation continues even though each chamber alternates between
filling and ejection.

``Continuously'' does not mean that the muscle contracts without pause. Each
cardiac cycle alternates between contraction (systole) and relaxation and
filling (diastole), and cardiac muscle has a refractory period after each
electrical impulse that prevents sustained contraction and allows the chambers
to fill again. Diastole occupies a variable fraction of the cycle and
becomes relatively shorter when the heart rate rises. The autonomic nerves and
circulating hormones adjust the SA node's rate—typically about 60--100 beats
per minute at rest in adults—and the strength of contraction in response to
exercise, sleep, stress, and other demands; they normally regulate the intrinsic
rhythm rather than initiating every beat. Persistent unexplained palpitations,
dizziness, fainting (syncope), or chest pain with an unusually fast, slow, or
irregular pulse warrants medical evaluation for rhythm or conduction disorders.

The heart itself needs oxygen and fuel, supplied by the coronary arteries. Much
of the blood flow through the left coronary circulation occurs during diastole,
when contraction of the ventricular muscle is no longer compressing those
vessels; coronary-flow patterns differ between regions and between the left and
right sides.

\section{Why is the heart slightly left of centre?}\index{heart}\index{diaphragm}\index{embryonic development}\index{lungs}

The heart is not entirely on the left. It lies obliquely in the middle of the
chest, behind the breastbone, within a region called the mediastinum. Its broad
base is relatively posterior, while its pointed lower end, the apex, angles
downward and to the left. As a rough teaching approximation, about two-thirds
of the heart lies left of the body's midline, although the exact position and
the proportion on either side vary among people. The apparent left-sided
location mainly reflects this orientation, not a heart that is attached wholly
to the left side.

During embryonic development, left--right patterning signals help establish the
body axis, while the initially simple heart tube bends, loops, and remodels into
chambers and connecting vessels. This three-dimensional process establishes the
usual orientation of the atria, ventricles, great vessels, and apex. The lungs,
diaphragm, spine, and surrounding organs then help hold the heart in the central
thorax.

The apex normally points toward the left fifth intercostal space, near the
midclavicular line. Movement of the apex against the chest wall can create a
palpable \emph{apical impulse},
although it may be difficult to feel because of body shape, lung inflation,
chest-wall thickness, or the heart's position. Anatomy is not identical in
everyone. In the uncommon condition \emph{dextrocardia}, the heart is oriented
with its apex toward the right. Dextrocardia can occur alone or with other
structural abnormalities; when the organs are also arranged as a broader
mirror-image pattern, the condition is called \emph{situs inversus}.

\chapter{Breathing, Healing, and Immunity}

\section{Why do we breathe faster during exercise, even before blood oxygen falls?}\index{oxygen}\index{ATP (adenosine triphosphate)}\index{brain}\index{breathing}\index{carbon dioxide}\index{heart}

Working muscles use ATP more quickly during exercise. They therefore usually
need more oxygen and produce more carbon dioxide. Total ventilation increases
through a combination of faster breaths and deeper breaths. It often rises
within seconds because the brain's motor commands and signals from moving
muscles and joints anticipate the increased demand, before a large change in
blood gases occurs. This anticipatory response is often called feed-forward
control, and it is later adjusted by chemical feedback.

As exercise continues, central chemoreceptors in the medulla and peripheral
chemoreceptors in the carotid and aortic bodies help adjust ventilation in
response to carbon dioxide, hydrogen-ion concentration (acidity), and, in some
situations, low arterial oxygen. Minute ventilation can rise from roughly
5--8 litres per minute at rest to over 100 litres per minute during intense
exercise. During ordinary exercise at sea level, arterial oxygen is usually
kept close to its resting level, so breathing does not wait for a major fall in
oxygen. Ventilation therefore helps regulate carbon dioxide and acidity as well
as supplying oxygen.

During very intense exercise, the muscles can develop a net acid load faster
than the body can buffer or remove it. Bicarbonate buffers some of these acids,
generating extra
carbon dioxide and stimulating ventilation further. The heart also usually beats
faster and pumps more blood, delivering oxygen to active tissues and carrying
carbon dioxide to the lungs for exhalation. The size and timing of these
responses depend on intensity, fitness, altitude, health, and the person's
individual physiology.

The response therefore anticipates demand and helps regulate oxygen, carbon
dioxide, and acidity rather than responding only to falling blood oxygen.

\section{Why do wounds heal, and why do scars form?}\index{scars}\index{wound healing}\index{blood!clotting}\index{collagen}\index{skin!epidermis}\index{hair}

Wound healing limits bleeding and infection, closes the injury, and restores
the skin's barrier. It proceeds through four overlapping phases. During
\emph{haemostasis}, blood vessels constrict and platelets and fibrin form a
temporary clot. During \emph{inflammation}, immune cells remove microbes and
damaged material and release signals that coordinate repair. During
\emph{proliferation}, keratinocytes cover the surface, fibroblasts produce
extracellular matrix, and endothelial cells help form new small blood vessels.
Myofibroblasts can contract the wound edges, helping close the gap. During
\emph{remodelling}, the provisional tissue is reorganised and strengthened over
months or longer.

If an injury extends beyond the superficial epidermis into the dermis, the body
may close it with collagen-rich repair tissue rather than recreating the
original skin architecture. A scar can therefore restore continuity and
considerable strength, but it is usually less flexible and less fully functional
than uninjured skin, commonly lacking the original arrangement of collagen,
nerves, hair follicles, sweat glands, and pigment cells. Scar appearance and
healing depend on depth, location, wound tension, infection, blood supply, age,
genetics, and other health factors, whereas very superficial injuries limited to
the epidermis can heal with little or no visible scar.

Scars may initially be firm, raised, itchy, red, or painful and can continue to
soften and change colour for months or years. A \emph{hypertrophic scar} remains
within the original wound edges and often becomes less prominent with time. A
\emph{keloid} extends beyond those edges and may continue growing; the tendency
to form keloids is influenced strongly by genetics and previous injury. Both
result from excessive or poorly regulated repair rather than from a failure of
the body to heal. During healing, spreading redness, increasing warmth or
swelling, pus, worsening pain, red streaks, or fever can signal wound infection
and warrant prompt medical assessment.

\section{Why do bones heal but spinal-cord injuries often do not?}\index{bones}\index{inflammation}\index{neurons}\index{scars}\index{spinal cord}

Bone usually heals more effectively because it has a strong blood supply,
renewable matrix-producing cells, and a supportive structural environment. The
adult spinal cord contains specialised neurons and long axons in a tightly
organised central-nervous-system environment that is much less permissive to
regrowth. Neither process is perfect: bone can heal with deformity or fail to
unite, and spinal-cord injury can leave substantial function when pathways are
spared.

\subsection*{How bone repairs itself}

Bone is a living, vascular tissue with cells that can divide and produce new
matrix. After a fracture, bleeding and inflammation create a temporary repair
environment; blood vessels, immune cells, and progenitor cells help form a soft
and then a hard callus. Osteoblasts make new bone, osteoclasts remove and
reshape it, and remodelling can restore much of its strength when the fragments
have adequate blood supply, alignment, and mechanical stability. Bone healing
is not automatic: infection, poor blood supply, excessive movement, large gaps,
or illness can delay it or cause non-union.

\subsection*{Why spinal-cord repair is limited}

The adult spinal cord is part of the central nervous system and has a much more
limited repair environment than bone or a peripheral nerve. Injury can kill
neurons and supporting cells, sever axons, disrupt myelin, and create a
fluid-filled cavity, while secondary inflammation, swelling, and loss of blood
flow can enlarge the original damage. Surviving central neurons have limited
intrinsic ability to regrow their axons. Molecules associated with damaged
myelin and the reactive glial scar can inhibit extension, although the glial
scar also helps contain damage and restore a protective boundary. Even if an
axon regrows, it must survive, become remyelinated, and reconnect with the
correct targets across complex pathways.

\subsection*{Partial recovery and current research}

Spinal-cord injury does not always mean complete or permanent loss of function.
Some axons and neurons may survive, swelling may subside, and rehabilitation
can use neuroplasticity to strengthen spared or alternative pathways. These
mechanisms explain why recovery can occur without rebuilding the entire cord;
functional improvement does not necessarily mean that severed axons have
regrown across the injury.
Rehabilitation and supportive care are established parts of treatment, whereas
cell-based approaches, growth-promoting signals, and engineered scaffolds
remain areas of research rather than reliable cures. Restoration of function
after severe injury remains a major clinical challenge.

\section{Why do we feel pain?}\index{pain}\index{brain}\index{spinal cord}\index{stress}\index{touch}

Pain is an unpleasant sensory and emotional experience associated with actual or
potential tissue damage, but it is not a direct readout of injury. It often
protects us by encouraging withdrawal, rest, or care, yet pain can persist after
tissues have healed or occur without a clear injury. Special nerve endings
called \emph{nociceptors} respond to potentially damaging heat, mechanical
forces, and chemical changes. The neural detection and processing of those
signals is called \emph{nociception}; it is not itself pain. Pain emerges when
the spinal cord and brain interpret sensory information in the context of
attention, emotion, expectation, past experience, and the person's biological
and social circumstances.

\subsection*{Types of pain and sensitisation}

Pain is real even when there is no continuing tissue injury. \emph{Nociceptive
pain} is driven mainly by actual or threatened tissue damage. \emph{Neuropathic
pain} results from a lesion or disease affecting the somatosensory nervous
system. \emph{Nociplastic pain} involves altered pain processing without clear
evidence of ongoing tissue damage or a definable nerve lesion. Peripheral
sensitisation can make injured tissues' nerve endings respond more strongly,
while central sensitisation can increase responsiveness within the spinal cord
and brain.

\subsection*{How pain signals are carried and modulated}

Different nociceptors detect potentially harmful pressure, temperature, or
chemical changes. Faster A-delta fibres often contribute to sharp,
well-localised first pain, while slower C fibres often contribute to delayed,
aching or burning pain; these are general patterns because the spinal cord and
brain combine and modify the signals, rather than treating fibre type as a
fixed pain category.

Incoming signals are filtered in the spinal cord and modulated by descending
pathways from the brain. Attention, threat, expectation, touch, stress, and
endogenous chemicals such as opioid peptides can increase or reduce the final
experience. Conversely, nociception may occur without conscious pain, such as
during unconsciousness or a brief protective reflex. Pain is therefore not a
precise meter of tissue damage, and a person's report remains central to
assessing it. Severe, persistent, recurrent, or unexplained pain deserves
assessment of both the body and the nervous system.

\section{How does the immune system recognise germs, and how do innate and adaptive immunity differ?}\index{immune system!adaptive immunity}\index{immune system}\index{antibodies}\index{immune system!innate immunity}\index{memory}\index{proteins}

\subsection*{Innate immunity: rapid, broad detection}

The immune system uses several layers of detection. Physical barriers and
antimicrobial chemicals act before deeper recognition begins. Innate immune
cells carry pattern-recognition receptors—built-in molecular sensors—that
detect recurring molecular features associated with groups of microbes, such as
parts of bacterial cell walls, as well as signals released by damaged or
stressed host cells. Complement and other soluble proteins provide additional
surveillance and can coat microbes to make them easier for scavenger cells to
engulf. Innate recognition is rapid but relatively broad; it does not identify
an entire microbial species with a single receptor.

\subsection*{Adaptive immunity: specific receptors and memory}

Adaptive lymphocytes carry highly varied receptors generated before an
infection. By chance, some B or T cells have receptors that fit part of the
invading material. If the appropriate activation signals are present, those
matched cells multiply rapidly (\emph{clonal selection and expansion}) to build
an army of identical responders. B-cell receptors can bind intact
three-dimensional parts of an antigen (a foreign target molecule); activated B
cells can become plasma cells that secrete antibodies to neutralise or mark the
threat. Most T-cell receptors recognise short peptide fragments displayed on
cell surfaces by major histocompatibility complex (MHC) molecules—molecular
display trays that show the immune system what is inside or captured by a cell.
CD8 T cells can kill some infected or abnormal cells, while CD4 T cells
commonly help coordinate immune responses, support B cells, or activate other
cells.

\subsection*{How the two arms work together}

Dendritic cells and other antigen-presenting cells help connect the two arms.
They capture material, respond to innate danger signals, and display antigen
fragments on MHC molecules to T cells in lymphoid tissues. A second safety
signal (\emph{co-stimulation}, confirming real danger) and local inflammatory
signals help determine whether that recognition produces an effective response
or tolerance. T cells therefore generally recognise an
antigen in a cellular context rather than binding a free microbe in the same way
as an antibody.

Recognition alone does not always produce an attack. Immune tolerance normally
limits responses to the body's own molecules, and activation depends on context,
including danger signals, co-stimulatory signals, the amount and location of
antigen, and regulatory cells. If an infection is controlled, some B and T cells
persist as memory cells and can respond faster to a later encounter. The
strength and duration of memory vary with the pathogen, infection, vaccine,
and person. The initial adaptive response is generally slower and more specific
than innate recognition, but memory B and T cells can respond rapidly during a
later exposure. Innate immunity is rapid and broad; adaptive immunity uses B and
T cells, has highly specific receptors, and can form memory. The two arms
cooperate, and neither is perfect.

\section{Why do vaccines work?}\index{vaccines}\index{antibodies}\index{DNA}\index{immune system}\index{memory}\index{cell nucleus}

Vaccines expose the immune system to a carefully selected target from a
pathogen. Depending on the vaccine, this may be an inactivated or weakened
microbe, a purified subunit, a detoxified toxin (toxoid), a conjugated antigen,
or genetic instructions that allow cells to make a pathogen component. Genetic
instructions in vaccines such as mRNA vaccines provide a temporary template;
they do not normally enter the cell nucleus or rewrite a person's DNA. These
materials are designed not to cause the full disease, although vaccination can
produce expected local or systemic reactions and, rarely, serious adverse
effects. Some live-attenuated vaccines require special precautions in people
with severe immune suppression, and the relevant precautions depend on the
specific vaccine and person.

Antigen-presenting cells process the vaccine material and activate lymphocytes
with matching receptors. B cells can produce antibodies, while helper and
cytotoxic T cells coordinate or carry out cellular responses. Some activated B
and T cells become long-lived memory cells. An adjuvant is included in some,
but not all, vaccines to strengthen the immune response, and some vaccine
schedules use additional doses to improve or extend protection.

If the real pathogen is encountered later, antibodies and memory B and T cells
can respond more rapidly and strongly. This may prevent infection, shorten the
infection, or make disease less severe, depending on the pathogen, vaccine,
time since vaccination, and person. A vaccine may provide strong protection
against severe disease while providing less consistent protection against
infection or transmission.
Protection is not always complete: immunity can wane, pathogens can change,
exposure dose can matter, and individual immune responses differ. Vaccines can
reduce infection and transmission for some diseases, but for others their
strongest and most consistent benefit is protection against severe illness,
complications, and death. High vaccination coverage can also reduce spread and
help protect people who cannot be vaccinated or who respond poorly to vaccines.

\chapter{Reflexes, Sleep, Ageing, and Development}

\section{Why do we have an appendix?}\index{appendix}\index{antibiotics}\index{digestion}\index{fever}\index{inflammation}\index{large intestine}

The appendix is a narrow, blind-ended, finger-like pouch attached to the caecum,
the first part of the large intestine. It is not required for digestion, and
most people remain healthy after it is removed. Its lymphoid tissue participates
in gut immune activity, particularly earlier in life. The appendix may also
help maintain or reseed parts of the gut microbial community after disturbance,
but the size and importance of this proposed reservoir function in humans remain
uncertain. The appendix may have more than one function, and no single
evolutionary explanation has been established.

The appendix's narrow opening can become blocked by enlarged lymphoid tissue,
hardened stool (a fecalith), a foreign body, a parasite, a growth, or other
causes. Obstruction is common but not required in every case; impaired drainage,
infection, or inflammation can all contribute to \emph{appendicitis}. Untreated
appendicitis can progress to perforation (a hole through the appendix wall), a
walled-off pocket of pus (an abscess), or peritonitis (widespread inflammation
of the abdominal lining). Pain often begins near the navel and later localises
to the right lower quadrant—typically around McBurney's point, roughly a third
of the way from the hip bone toward the navel—though symptoms vary and the
pattern may be atypical in children, pregnancy, or older age. Worsening
abdominal pain, fever, vomiting, or marked tenderness warrants prompt medical
assessment rather than home treatment; symptoms alone cannot establish the
diagnosis. Hospital evaluation may include examination, blood or urine tests,
and imaging such as ultrasound or CT. Treatment
usually involves appendectomy, although antibiotics alone may be appropriate in
selected uncomplicated cases or when surgery is unsuitable.

\section{What are vestigial features of the human body?}\index{vestigial features}\index{appendix}\index{goosebumps}\index{hair}\index{muscles}\index{teeth}

Vestigial features are inherited structures or responses whose ancestral
function has been reduced or altered during evolution. Some are smaller or
simpler than the corresponding ancestral feature, but reduced size is not
required. ``Vestigial'' does not mean completely useless: a feature may retain
a minor, new, or supporting function even when its original role has largely
disappeared. It is vestigial relative to an ancestral function, not merely
because it is unnecessary for an individual to survive.

Examples include the coccyx, which is the remnant of an ancestral tail but now
anchors pelvic muscles and ligaments; the appendix, which contains immune
tissue and may help support gut microbes; and the small muscles around the
outer ear, which most people cannot use to point their ears effectively.
Goosebumps are another example of a retained reflex: raising body hair could
make a furry ancestor look larger or help conserve heat, but it has little
effect on a mostly hairless human. Wisdom teeth are reduced third molars that
can still help with chewing, but often lack enough room to erupt normally,
leaving them impacted or prone to painful gum inflammation
(\emph{pericoronitis}) and crowding. Other classic human vestigial structures
include the \emph{palmaris longus} forearm tendon (absent in a notable minority
of people without loss of grip strength) and the \emph{plica semilunaris}, a
small fold at the inner corner of the eye corresponding to a third eyelid in
some other animals. Although most vestigial features are harmless, knowing when
one causes trouble—such as persistent tailbone pain (\emph{coccydynia}) after a
fall or recurrent pain and swelling around an emerging wisdom tooth—helps
readers recognise when medical or dental care is needed.

These examples do not prove that every nonessential structure is an evolutionary
leftover. Functions can be lost, modified, or shared with other tissues, and
the evolutionary history of a feature must be assessed from anatomy,
development, genetics, and comparative biology rather than from its apparent
usefulness alone.

\section{Why do we get hiccups?}\index{hiccups}\index{brainstem}\index{brain}\index{breathing}\index{diaphragm}\index{larynx and voice}

A hiccup is a brief, involuntary contraction of the diaphragm and sometimes
other inspiratory muscles, followed almost immediately by closure of the
glottis, the opening between the vocal folds. This sudden interruption of
inward airflow produces the characteristic sound. Sensory information from the
chest, diaphragm, stomach, and other areas can travel through pathways including
the vagus, phrenic, and sympathetic nerves to a brainstem reflex network; motor
signals then return to the breathing muscles and larynx. Hiccups also occur
commonly before birth and in infants. Their normal biological purpose in humans
is not known, and proposed explanations remain speculative.

Eating quickly, stomach distension, carbonated drinks, alcohol, sudden
temperature changes, and strong emotions can trigger short-lived episodes.
Most stop spontaneously and are harmless. Hiccups lasting longer than 48 hours
are called \emph{persistent} hiccups; those lasting about a month are often
called \emph{intractable}. Prolonged or frequently recurring hiccups can
interfere with eating, drinking, sleeping, or breathing and should be medically
assessed. Reflux, medicines, metabolic disturbances, and diseases affecting
the chest, nerves, or brain can occasionally be responsible. Severe breathing
difficulty, inability to keep fluids down, confusion, or other acute illness
requires urgent care. The cause is not always found.

\section{Why do we sneeze?}\index{sneezing}\index{allergies}\index{brainstem}\index{coughing}\index{inflammation}

Sneezing is an involuntary protective reflex that helps clear the nasal passages
of mucus, particles, and irritating chemicals. Sensory endings in the nasal
lining, especially fibres carried by branches of the trigeminal nerve, detect
irritation or inflammation and signal brainstem circuits. These circuits
coordinate a deep inhalation, brief closure and reopening of the glottis, and a
rapid, forceful expiration through the nose and mouth. Infections and allergies
can trigger sneezing, but so can dry air, dust, strong smells, and other
non-infectious stimuli.

Some people sneeze after exposure to bright light; this is often called the
\emph{photic sneeze reflex}. Its mechanism and frequency vary among individuals,
and the tendency can run in families. It is not a sign that bright light has
damaged the eyes. Sneezing is mainly a defence for the upper airway and is
distinct from coughing, which more often clears the throat and lower airways;
the two reflexes can overlap when irritation affects several regions.

A sneeze produces droplets and smaller airborne particles. They are not
automatically infectious, but they can carry respiratory pathogens when the
sneezer is infected. Sneezing into a tissue and disposing of it, or into the
crook of the elbow, helps reduce spread; hands should be washed if they become
contaminated. It is safer to let a sneeze out than to forcefully block both the
nose and mouth, which can sharply increase pressure and rarely injure tissues.
Sneezing may expel some irritants but does not remove every particle or cure an
infection, so persistent symptoms should be interpreted in their wider context.

\section{Why do we cough?}\index{coughing}\index{allergies}\index{antibiotics}\index{brainstem}\index{larynx and voice}\index{muscles}

Coughing is a protective reflex that helps clear mucus, dust, fluids, food, and
other irritants from the throat and lower airways. Sensory receptors in the
larynx, trachea, and bronchi send signals through pathways including the vagus
nerve to brainstem circuits. The resulting sequence usually includes a deep
inhalation, closure of the vocal folds, contraction of expiratory muscles to
build pressure, and sudden reopening and forceful expiration. The high airflow
can move material toward the mouth. Coughing may be productive, bringing up
mucus, or dry, without easily visible mucus. It can be reflexive or voluntarily
started, and an over-sensitive cough pathway can itself become part of a
problem.

An acute cough lasts up to about three weeks and commonly accompanies a viral
respiratory infection. A subacute cough lasts roughly three to eight weeks; a
cough lasting more than eight weeks in an adult is generally considered chronic,
although definitions and evaluation thresholds differ in children. Asthma,
allergies, reflux, smoking, infection, aspiration, other airway disease, and
medicines such as ACE inhibitors can cause a longer-lasting cough. A persistent
or unexplained cough deserves medical evaluation, even when no mucus is
produced; duration alone does not establish whether an infection is bacterial
or whether antibiotics are needed. Severe breathlessness, blue or grey lips,
chest pain, choking, confusion, or coughing blood requires urgent medical
attention.

\section{Why do we yawn?}\index{yawning}\index{brain}\index{breathing}\index{headaches}\index{hearing}\index{muscles}

Yawning is a widespread, partly involuntary behaviour involving prolonged
opening of the mouth, stretching of the jaw and face, changes in breathing, and
exhalation. Brain circuits in the hypothalamus (notably the paraventricular
nucleus) and brainstem coordinate the reflex through neurotransmitters such as
dopamine, oxytocin, and orexin. Yawns vary in depth and duration, and a person
can sometimes yawn voluntarily. They are common when sleep pressure is high,
during transitions between sleep and wakefulness, and during changes in
alertness or attention. Contrary to a common myth, ordinary yawning is not
simply a response to a lack of oxygen in the blood.

Yawning can be contagious: seeing, hearing, reading about, or even thinking
about another person yawning can increase the chance of yawning in some people.
This may involve attention, automatic imitation, expectation, and social
context. Although contagious yawning has sometimes been linked with empathy,
the association is inconsistent and it is not a reliable test of empathy or
social concern.

Yawning-like movements occur in many vertebrates and are observed before birth,
so yawning cannot be explained only as social imitation. Proposed functions
include helping shift arousal during transitions between sleep and wakefulness,
stretching the jaw and face muscles, and regulating brain temperature. The
thermoregulatory and arousal accounts remain hypotheses: some findings support
parts of them, but no single function has been established for every yawn.
Yawning is therefore not a dependable way to raise blood oxygen or diagnose an
oxygen problem. A sudden, marked increase in yawning, especially with confusion,
weakness, fainting, severe headache, or other new symptoms, deserves medical
assessment, as excessive yawning can occasionally accompany migraine prodromes,
vasovagal episodes, brainstem strokes or lesions, or certain medicines such as
SSRIs and dopaminergic drugs; ordinary occasional yawning is generally normal.

\section{Why do we get goosebumps?}\index{goosebumps}\index{nervous system!autonomic}\index{hair}\index{muscles}\index{shivering}\index{skin}

Goosebumps, or \emph{piloerection}, occur when postganglionic sympathetic
nerves release noradrenaline onto alpha-1 adrenergic receptors on the tiny
involuntary arrector pili muscles attached to hair follicles, causing them to
contract. The hairs stand up and pull the surrounding skin into small raised
areas. Cold commonly triggers this response, but fear, excitement, music, and
other intense emotions can also activate sympathetic pathways and produce a
similar appearance. Goosebumps are therefore an involuntary body response, not
a conscious choice. The visible bumps result from piloerection, the hair-raising
response; the accompanying subjective feeling of chills can occur with or
without obvious raised hairs. Clinically, prominent piloerection can also
appear during fever-induced rigors, opioid withdrawal, or spinal autonomic
dysreflexia; when chills and goosebumps accompany a high fever, confusion, or
rigors, medical assessment is warranted.
Emotional goosebumps reflect autonomic arousal rather than necessarily a fall
in skin temperature.

In furry mammals, raised hair can trap a layer of insulating air and make the
animal appear larger to a threat. In humans, body hair is relatively sparse and
short, so raising it provides little insulation and rarely changes apparent
body size enough to deter a threat. Cold-induced narrowing of skin blood vessels
is a separate sympathetic response and does more to reduce heat loss; shivering
is also much more important for generating heat. Human piloerection is therefore
a reflex whose protective effects are much smaller than in furry mammals. Its
exact present-day function in humans is uncertain, and it should not be assumed
to be completely functionless simply because its insulating effect is small.

\section{Why do we dream?}\index{dreaming}\index{brain}\index{heart}\index{learning}\index{memory}\index{sleep}

Dreams are reported experiences associated with changing patterns of brain
activity during sleep. Because researchers usually learn about a dream by
awakening a person and asking for a report, the private experience cannot be
measured as directly as heart rate or brain waves. Dreams are often vivid and
more frequently recalled after rapid-eye-movement (REM) sleep—which makes up
roughly 20--25\% of an adult's night and features strong cholinergic brainstem
activation, suppressed noradrenaline and serotonin, active limbic emotion
centres, and temporary muscle paralysis (\emph{atonia})—but they also occur
during non-REM sleep. Failure to remember a dream therefore does not prove
that no dream experience occurred. Dream content can combine fragments of
recent and older memories with emotion, sensations, and imagination into unusual
narratives; it is not simply a complete recording of the day.

Dreaming may be related to memory processing, emotional regulation, the brain's
ongoing predictions about the world, or the interpretation of internally
generated sensory signals. These are active research questions, and no single
function has been established for all dreams. Evidence that sleep supports
learning or emotional regulation does not prove that every dream has a separate
beneficial purpose. A dream's particular story is not a reliable prediction of
the future or a universally encoded symbolic message. A dream may have personal
meaning to the dreamer, but generalized dream dictionaries and fixed symbolic
interpretations are not scientifically validated.

From a health perspective, occasional vivid or unpleasant dreams are normal, and
medicines such as antidepressants, beta-blockers, or alcohol withdrawal can
alter REM sleep and dream intensity. However, frequent distressing nightmares
that disrupt sleep or daytime functioning (nightmare disorder or trauma-related
nightmares in PTSD), or physically acting out dreams through punching, kicking,
or shouting—which suggests REM sleep behaviour disorder due to loss of normal
muscle atonia—warrant medical evaluation.

\section{Why do some people sleepwalk or experience sleep paralysis?}\index{sleep!paralysis}\index{sleep}\index{sleep!sleepwalking}\index{brain}\index{breathing}\index{fever}

Sleepwalking is a non-REM parasomnia, often arising from deep stage-N3 sleep
and more often occurring in the first part of the night. A person partly arouses
while brain systems for movement are active but awareness, memory, and judgement
remain limited. The person may sit up, walk, speak, or perform routine actions
and often remembers
little the next morning. Sleepwalking is especially common in children and
often becomes less frequent with age. Sleep deprivation, irregular schedules,
fever, stress, alcohol, obstructive sleep apnoea, and some medicines can
increase the likelihood.

Sleep paralysis occurs while falling asleep (a \emph{hypnagogic} episode) or
waking (a \emph{hypnopompic} episode). Conscious awareness returns before the
normal REM-related muscle atonia has ended, so the person temporarily cannot
move or speak. Vivid visual, auditory, or bodily hallucinations can occur,
including a frightening sense of pressure on the chest. Episodes usually last
seconds to a few minutes and are not usually dangerous, although they can be
distressing. Breathing and eye movements usually remain possible, even when the
feeling of chest pressure makes breathing seem difficult. They may occur in
isolation or recur with sleep deprivation, irregular schedules, stress, or
narcolepsy. They reflect a temporary sleep-state transition, not a supernatural
event or, by themselves, psychosis.

Either condition may be benign, but sleepwalking can cause injury. Episodes
that are frequent, dangerous, newly developing in adulthood, associated with
major daytime sleepiness, cataplexy, breathing pauses, unusual violent movements,
or severe distress deserve medical evaluation. A clinician may need to distinguish
these events from seizures, REM-sleep behaviour disorder, medication effects, or
sleep-disordered breathing. Practical safety measures include removing sharp or
breakable objects, securing windows and exterior doors, and keeping the sleep
environment clear. A sleepwalker should generally be guided calmly back to bed
rather than restrained or forcibly startled awake.

\section{Why do we age?}\index{ageing}\index{DNA}\index{epigenetics}\index{genes}\index{inflammation}\index{mitochondria}

\subsection*{Cellular and molecular mechanisms}

Ageing is the gradual, system-wide change in structure and function that occurs
with time. It is not controlled by one biological clock. Molecular damage and
altered regulation accumulate, while repair and maintenance become less
effective. Important interacting processes include DNA damage and mutation,
epigenetic change, changes in telomeres, loss of protein quality control,
mitochondrial dysfunction, cellular senescence, reduced stem-cell function,
altered nutrient sensing, chronic low-grade inflammation, and changes in
communication between cells and organs. These processes differ among tissues
and influence one another, so some are causes in one setting and consequences
in another. For example, cellular senescence can temporarily stop damaged cells
from dividing, but the accumulation of senescent cells and their secreted
signals can disrupt nearby tissue. Repair systems do not fail all at once, but
their declining reserve makes tissues less resilient to stress and disease.

\subsection*{Biological age, health, and disease}

Genes, development, infections, nutrition, environment, and behaviour influence
the rate and pattern of ageing, so chronological age is not the same as the
condition of a person's tissues. ``Biological age'' is not one directly observed
quantity; it is an estimate based on selected biomarkers or predictive models,
and different models can give different results. Some age-related changes and
disease risks can be delayed or treated, and healthy life expectancy can be
extended—most reliably through regular physical activity, adequate sleep,
balanced nutrition, avoiding smoking, managing blood pressure and blood sugar,
and routine preventive care—but no intervention has yet been shown to stop or
reliably reverse whole-body human ageing. Across organ systems, normal ageing
typically brings gradual declines in glomerular filtration rate, muscle mass
(\emph{sarcopenia}), bone mineral density, and arterial elasticity, though rates
vary widely among individuals. Ageing itself is a normal biological process, not
a single disease, so a particular symptom, disability, or condition should not
automatically be dismissed as ``just ageing.''

\subsection*{Evolutionary perspectives}

Evolution may also help explain why ageing persists. Natural selection is often
stronger before and during the main reproductive years than late in life, so
harmful effects that appear mainly later may be removed less efficiently. Early
benefits can also be linked to later costs through trade-offs in growth,
reproduction, repair, and energy use—ideas formalised in evolutionary models
such as \emph{antagonistic pleiotropy} and the \emph{disposable soma} theory.
These are broad evolutionary ideas, not a single proven explanation for every
aspect of human ageing.

\section{Why does hair turn grey?}\index{hair}\index{ageing}\index{autoimmune disease}\index{hair!greying}\index{melanin}\index{stem cells}

Hair turns grey or white when the follicle supplies less melanin to newly formed
hair shafts. With ageing, melanocyte stem cells and mature pigment-producing
cells can be depleted, become less active, or fail to maintain pigment
production. A mixture of pigmented and unpigmented hairs looks grey, while an
individual shaft containing very little or no melanin usually looks white. The
visible shaft is already a dead, keratinised structure, so it normally does not
lose pigment suddenly along its length; pigment loss is usually seen as new
growth from the follicle. Patchy changes can occur when different follicles
lose pigment at different times.

Genetics strongly influences when greying begins, and ordinary age-related
greying is not itself a disease. Smoking, some nutritional deficiencies (such as vitamin B12 deficiency),
autoimmune conditions (such as vitiligo, alopecia areata, or thyroid disease),
pigment disorders, and medical treatments can contribute
to premature or patchy loss of pigment in particular people, but they do not
explain most normal age-related greying. Premature greying—often defined as
onset before roughly age 20 to 30 depending on ancestry—or greying that is
abrupt, very patchy, or accompanied by other symptoms is more worth discussing
with a doctor.

Psychological stress can affect follicle biology in ways that may accelerate
pigment loss or reveal hairs that were already growing without pigment, but a
single stressful event does not normally bleach an existing shaft overnight.
Grey hair primarily contains less pigment; its texture may also seem drier or
coarser because of changes in the follicle and scalp, not because the lack of
melanin automatically makes the shaft structurally damaged. Plucking one grey
hair does not make other follicles turn grey, although the regrown hair from
that follicle will usually have the same pigment status.

\section{Why do people have different voices?}\index{brain}\index{breathing}\index{language acquisition}\index{larynx and voice}\index{lungs}\index{pain}

Voice begins when air from the lungs passes through the larynx and makes the
vocal folds vibrate. Pitch, loudness, and voice quality are distinct features:
the length, mass, tension, and vibration frequency of the vocal folds largely
determine pitch; vibration amplitude and outgoing air pressure govern loudness;
and filtering in the throat, pharynx, mouth, and sometimes nasal passages
shapes the balance of harmonics and resonance, creating \emph{timbre}---the
quality that helps one person's voice sound different from another's. The
tongue, lips, jaw, and soft palate shape speech sounds and pronunciation.

People differ in the anatomy of these structures and in how the brain controls
breathing, vocal-fold tension, resonance, and articulation. Growth at puberty
and hormonal changes can alter vocal-fold length and mass; age, illness,
medicines, injury, and training can also change the sound. Language and accent
mainly affect speech patterns and pronunciation, while learned vocal habits
affect both speech and voice quality. Emotion and social context can alter
pitch, loudness, timing, and quality. A person's voice is therefore both
anatomically individual and dynamically controlled.

On average, puberty-related anatomy produces some sex-associated differences in
voice—typical adult fundamental speaking frequencies cluster around roughly
85--180 Hz in males and 165--255 Hz in females—but these are overlapping
population tendencies, not a reliable way to determine an individual's sex or
gender. Persistent hoarseness lasting longer than 2 to 4 weeks, pain while
speaking, or a sudden voice change should be medically assessed (often including
visualisation of the vocal folds by laryngoscopy) to rule out causes such as
vocal-fold nodules, acid reflux irritation, or recurrent laryngeal nerve palsy.

\section{Why do we blush?}\index{blushing}\index{nervous system!autonomic}\index{diarrhoea}\index{fever}\index{menopause}\index{skin}

Blushing is a visible, involuntary increase in blood flow through superficial
skin vessels, especially in the face, ears, neck, and sometimes the upper chest.
It is particularly associated with self-conscious emotions such as
embarrassment, shame, praise, excitement, or being the focus of attention. The
autonomic nervous system changes vascular tone—with evidence implicating
sympathetic vasodilator and beta-adrenergic pathways in facial skin—but the
precise neural and vascular mechanism of emotional facial blushing is unusual
and not fully understood. Heat, alcohol, some medicines, menopause, fever, exercise, and
skin or systemic disorders can also cause facial \emph{flushing}, which is not
necessarily emotional blushing. Emotional blushing is typically brief and
localized, whereas flushing may be more diffuse, prolonged, recurrent, or
associated with heat, medication, hormonal change, or illness; appearance alone
cannot reliably identify its cause.

Humans may be unusually susceptible to conspicuous social blushing compared
with many other animals, but the size and significance of this difference are
not simple to establish. One evolutionary hypothesis is that visible blushing
can signal awareness of a social mistake or vulnerability and sometimes support
social repair. This possible signalling effect is not established as the sole
or universal explanation, and a blush does not reliably indicate guilt,
honesty, attraction, or a person's intentions.

Brief blushing during an emotional moment is usually normal. Recurrent flushing
without a clear trigger, flushing that is persistent or worsening, or flushing
accompanied by wheezing, faintness, swelling, fever, diarrhoea, or other
concerning symptoms deserves medical assessment because medicines (such as
vasodilators or niacin), hormonal changes, skin conditions such as rosacea, and
systemic illnesses including anaphylaxis, mastocytosis, or carcinoid syndrome
can contribute.

\section{Why does menstruation occur?}\index{menstruation}\index{hormones}\index{hypothalamus}\index{implantation}\index{oestrogen}\index{pain}

The menstrual cycle coordinates changes in the ovaries and uterus in preparation
for a possible pregnancy. Pulsatile gonadotropin-releasing hormone (GnRH) from
the hypothalamus signals the pituitary to release
follicle-stimulating hormone and luteinising hormone, which coordinate ovarian
follicle development during the follicular phase and, in many cycles, ovulation,
followed by the post-ovulatory luteal phase. Correspondingly, the uterine
lining passes through a proliferative phase (when oestrogen stimulates
endometrial growth) and a secretory phase (when progesterone released after
ovulation prepares it for implantation). Typical adult cycles last roughly
21--35 days, with menstrual bleeding lasting about 2--7 days and average blood
loss around 30--40 mL (losses above 80 mL are clinically considered heavy).

If implantation does not occur, the temporary corpus luteum normally regresses
and progesterone and oestrogen levels fall. This withdrawal of hormonal support
causes blood vessels in the endometrium to constrict and the superficial lining
to break down. Blood, mucus, and endometrial tissue then pass through the vagina
as menstrual bleeding. If implantation occurs, early placental trophoblast
tissue produces human chorionic
gonadotropin (hCG), which helps maintain the
corpus luteum and hormonal support for the lining. Prostaglandins help the
uterus contract and can contribute to menstrual cramps.

Menstruation is one phase of the menstrual cycle, and cycles do not always occur
monthly or include ovulation. Hormonal contraception can suppress ovulation and
cause scheduled withdrawal bleeding that is not necessarily the same as a
natural menstrual period. Timing, flow, and
symptoms vary with age, health, stress, contraception, breastfeeding, and other
factors. Menstruation is not a process for cleansing toxins; it is the shedding
of uterine lining after a particular reproductive-hormone pattern. Bleeding
outside the expected pattern is not automatically menstruation and may have
other causes. Severe or worsening pain, unusually heavy or prolonged bleeding,
bleeding between periods, bleeding after sex, or a major change from a person's
usual pattern should be discussed with a health professional.

\section{Why can some organs regenerate while others cannot?}\index{regeneration}\index{cancer}\index{skin!epidermis}\index{heart}\index{inflammation}\index{liver}

Regeneration is not an all-or-nothing ability. It ranges from routine
replacement of short-lived cells, as in the skin and intestinal lining, through
repair of damaged tissue, to restoration of the mass and architecture of an
organ. Cell turnover alone is not the same as regenerating a whole organ: for
example, the epidermis is continually replaced, but deeper skin injuries may
heal with scar tissue. The outcome depends on the surviving cells' ability to
divide or change state,
resident stem or progenitor cells, blood supply, nerves, supporting extracellular
matrix, developmental signals, and control of inflammation.

The liver can often restore lost mass remarkably well—regrowing to its
functional size within weeks even after surgical removal of up to roughly 70\%
of healthy liver tissue—by activating surviving hepatocytes and progenitors,
although chronic disease, repeated injury, or extensive scarring can impair this
process. Restored liver mass does not always mean that the original microscopic
architecture or full function has returned. By contrast, most adult
heart-muscle cells and many central nervous-system neurons divide very little,
whereas surface epithelia (such as skin and gut lining) turn over every few
days to weeks and skeletal muscle can repair minor damage via resident
satellite cells. Peripheral nerves can also slowly regrow axons if their
connective-tissue sheath remains intact, unlike the central nervous system
where myelin inhibitors and glial scarring restrict regrowth. Their specialised architecture, electrical connections, and
long-distance wiring are difficult to reconstruct. Some limited cell
replacement and functional recovery can occur, but it is usually insufficient
after severe injury. Functional improvement can also result from surviving
neurons reorganising their connections, which is repair or plasticity rather
than replacement of the original tissue.

Inflammation is not simply the enemy of regeneration: an early, controlled
inflammatory response helps clear damage and initiate repair. Persistent or
excessive inflammation can instead cause fibrosis, disrupt the extracellular
matrix, and interrupt tissue connections. Thus an organ may regain some mass or
function without returning to its original structure.

Regeneration must also be carefully controlled. Cells must divide in the right
place and amount, differentiate into the correct types, connect with blood
vessels and nerves, and stop when repair is complete. Signals that stimulate
cell division can cause abnormal growth or tumours if safeguards fail. This
balance explains why increasing regeneration is a promising research goal but
not a simple matter of adding stem cells.

\section{How does the gut microbiome affect us?}\index{microbiome}\index{antibiotics}\index{nervous system!autonomic}\index{bile}\index{brain}\index{genes}

The gut microbiome is the community of trillions of bacteria, archaea, fungi,
viruses, and other microorganisms living in the digestive tract—most densely in
the colon—together with their genes and chemical products. The important
question is not only which organisms are present, but what they are doing and
where they are located. Microbes that ferment dietary fibre produce short-chain
fatty acids such as butyrate, propionate, and acetate that nourish colonic
cells, support the intestinal barrier, and influence metabolism and immune
activity. Gut microbes also compete with pathogens, modify bile acids, and can
produce or transform some vitamins, medicines, and other dietary compounds.
These effects depend on the species or strain, diet, gut location, and host.

The community and its activity change with diet, medicines (especially
antibiotics), infection, age, environment, and host biology. Broad-spectrum
antibiotics can disrupt protective commensals, sometimes allowing opportunistic
pathogens such as \emph{Clostridioides difficile} to overgrow and cause severe
diarrhoea—a condition in which faecal microbiota transplantation has proven
effective for recurrent cases—though the extent and recovery of microbiome
changes vary. Signals between the gut and brain can pass through
the enteric and autonomic nervous systems, hormones, immune mediators, and
microbial metabolites. These pathways are real, but many claims about mood,
autism, obesity, or other diseases still come from associations or animal
studies rather than proof that a particular microbe causes a human outcome.

There is no single microbiome profile or diversity score that defines health.
``More diverse'' is not automatically better, and a stool sample may not fully
represent microbes attached to the intestinal lining or their activity. Probiotic
effects are strain- and condition-specific; one species or product cannot be
assumed to have the effects of another, and a product that helps one person or
condition may have no measurable benefit in another. Many probiotic organisms
also do not permanently colonise the gut after supplementation. A varied,
fibre-containing diet can support many beneficial microbial functions, but no
single food or supplement reliably cures many unrelated diseases.

\section{Why are humans mostly right-handed?}\index{handedness}\index{brain}\index{genes}\index{language acquisition}

About 90\% of people show a preference for the right hand in many skilled tasks,
but this is an approximate population estimate. The proportion depends on the
task, the definition of preference, the person's age, and cultural history.
Handedness is multidimensional: writing, throwing,
using tools, and other activities do not always involve the same hand. Right-
hand preference is associated with asymmetries in brain networks for
movement—with voluntary fine motor control of each hand largely governed by the
opposite cerebral hemisphere via crossed corticospinal pathways—but it is not
rigidly determined by one brain region. Motor lateralisation and language
lateralisation are related in some people (around 95\% or more of right-handers
and roughly 70\% of left-handers show left-hemisphere language dominance) but
are not a simple one-to-one cause-and-effect relationship.

Many genetic variants, prenatal development, early brain organisation, practice,
culture, and social pressure contribute to handedness. There is no single
``handedness gene.'' Cultural pressure and forced switching have also changed
how handedness is reported in some populations and generations. A person's
preferred hand can be influenced by injury or neurological disease, and a
strong, fixed hand preference appearing unusually early in infancy (before
12--18 months) warrants paediatric assessment because it can occasionally signal
weakness on the opposite side; however, natural left-handedness and
mixed-handedness are normal human variation, not defects.
Hand preference is not identical to strength, coordination, or skill; practice
can improve performance with either hand without necessarily changing the
underlying preference.

One evolutionary possibility is that a strong population bias toward one hand
made shared tools, gestures, or cooperative skills easier to use. Other models
emphasise inherited brain lateralisation or developmental trade-offs. These
remain population-level hypotheses; no single evolutionary explanation has been
confirmed, and the reason right-handedness became common is still not fully
resolved.

\section{Why do identical twins become different?}\index{identical twins}\index{brain}\index{chromosomes}\index{DNA}\index{embryonic development}\index{epigenetics}

Monozygotic (``identical'') twins arise when one fertilised egg or early embryo
gives rise to two embryos. They begin with nearly the same nuclear DNA, but
``identical'' does not mean that every cell has exactly the same genome.
Post-fertilisation mutations and mosaic differences can affect one twin or only
particular cell lineages. Small differences in mitochondrial DNA may be present,
and in XX twins the pattern of X-chromosome inactivation can differ, but neither
is an inevitable feature of every pair.

Their development before birth is not identical. Differences in placental and
amniotic arrangement, blood flow, position, and local conditions can expose them
to different amounts of oxygen and nutrients. Random developmental events can
also influence growth, fingerprints, and which genes are active in each cell.
Epigenetic marks such as DNA methylation can therefore diverge without changing
the DNA sequence, altering gene activity in probabilistic, tissue-specific ways.
Not all identical twins share a placenta, so unequal placental blood flow is an
important possibility rather than a universal explanation.

After birth, infections, diet, sleep, stress, relationships, education,
exposures, and chance further shape the brain and body. Gene activity can
respond to these environments, so similar DNA can contribute to different
outcomes. Twins may therefore differ in appearance, height, fingerprints,
personality, immune responses, and disease risk while remaining more alike than
unrelated people. These differences are probabilistic: genes and environment
usually interact, and neither alone determines every trait.

\chapter{Senses, Memory, Skin, and Teeth}

\section{Why do our eyes receive an upside-down image?}\index{brain}\index{cornea}\index{neurons}\index{retina}\index{touch}

The cornea supplies most of the eye's refractive power (roughly 43 dioptres),
while the lens makes the fine adjustments (roughly 15--20 dioptres at rest,
increasing during near focus) needed to focus light on the retina. Because these
parts act as a converging optical system, the retinal image is inverted in both
vertical and horizontal directions---equivalent to a 180-degree rotation
relative to the scene. The retina does not send a small photograph to the
brain; its neurons encode the pattern of light at different locations and
begin extracting contrast, colour, motion, and edges, transmitting those
retinotopic signals along the optic nerve, optic chiasm, and lateral geniculate
nucleus to the primary visual cortex (V1).

The brain does not need to rotate an internal picture, because the retina sends
neural signals rather than a photograph that could be turned around. From early
development, neural circuits use the stable relationship between retinal
patterns and objects, movement, touch, and actions in the world, so perception is
experienced as normally oriented. No separate physical ``rotation'' occurs in the
brain. The important fact is the consistent spatial mapping, not whether the
signal is described as ``upside-down'' in isolation.

\section{Why can we not see our own nose all the time?}\index{brain}

The nose is close to the eyes and lies within the overlapping binocular visual
fields, so it can form an image on both retinas. However, it is usually a
stable, predictable feature in peripheral vision rather than the focus of
attention. Neural adaptation (including perceptual fading of stationary stimuli,
related to Troxler's fading) and selective attention reduce the prominence of
such unchanging information; the brain does not literally erase the nose from
the retinal signals, and it can still influence visual processing in the
background. Unlike a pathological visual-field defect or neurological neglect—in
which part of the visual scene is lost even when attention is directed toward
it—the nose immediately reappears whenever you look for it.

The nose becomes obvious when attention is deliberately directed toward it,
when the head or eyes move, or when its position changes relative to the scene.
Closing one eye can alter the nose's apparent outline and the way it obstructs
the scene, making it more noticeable for some viewing positions. This is an
example of habituation and selective attention: the brain prioritizes new or
behaviourally important information while reducing awareness of constant,
harmless features. Sensory adaptation may contribute, but the nose is not
physically removed from the visual signals.

\section{Why do humans blink and cry emotional tears?}\index{blinking}\index{tears}\index{nervous system!autonomic}\index{brain}\index{cornea}\index{hormones}

Blinking spreads a thin, three-layered tear film---composed of lipid from
eyelid Meibomian glands, aqueous fluid from the lacrimal gland, and mucin from
conjunctival goblet cells---across the surface of the eye. This film keeps the
cornea moist, supplies oxygen and nutrients to its outer layers, washes away
dust, and helps maintain a smooth optical surface. Reflexive blinking—mediated
by trigeminal (cranial nerve V1) sensory afferents and facial nerve (cranial
nerve VII) motor output—also shields the eye from bright light, air puffs, or
approaching objects, while spontaneous blink rates (normally around 12--20
blinks per minute) shift with attention and often drop during concentrated
reading or screen use. Persistent eye pain,
sudden vision changes, severe light sensitivity, or sudden inability to close
an eyelid should be medically assessed.

Humans produce three overlapping kinds of tears: basal tears that continuously
lubricate the eye, reflex tears triggered by irritants such as wind or onion
vapours, and emotional tears associated with grief, joy, or intense feeling.
Emotional crying involves autonomic and limbic brain circuits, and emotional
tears differ slightly in protein composition from reflex tears. However, popular
claims that crying ``detoxifies'' the body or clears meaningful amounts of
stress hormones are not supported by strong evidence.

\section{Why do smells bring back strong memories?}\index{memory}\index{amygdala}\index{brain}\index{hippocampus}\index{learning}\index{smell (olfaction)}

Odour molecules activate receptors in the nasal epithelium, and signals travel
along the olfactory nerve (cranial nerve I) through the olfactory bulb to the
piriform and other olfactory cortical areas. Unlike vision, hearing, and touch,
olfactory signals reach cortical and limbic regions without first passing
through a thalamic relay, communicating directly with networks involved in
emotion and memory such as the amygdala, entorhinal cortex, and hippocampus.
Smell does not bypass the brain's normal processing, but an odour can become a
particularly effective cue because it is often encountered together with the
sights, people, and feelings of an event. Clinically, a sudden or persistent
loss of smell (\emph{anosmia}) can follow viral infections, nasal blockage, head
trauma, or early neurodegenerative changes and warrants assessment.

During learning, the brain associates the odour with that surrounding context.
Later exposure can reactivate part of the stored network and help retrieve an
autobiographical memory, sometimes with a strong emotional or bodily quality.
This is an example of cue-dependent retrieval, not proof that smell memories
are always more accurate or vivid than memories prompted by other senses. The
effect depends on familiarity, attention, the original context, age, culture,
and individual differences in memory and smell sensitivity. The emotional
intensity or vividness of the recalled event does not guarantee that every
detail is accurate: autobiographical memories are reconstructed and can be
modified by later experience, expectation, and suggestion.

\section{How do taste and smell work together to create flavour?}\index{taste}\index{brainstem}\index{brain}\index{smell (olfaction)}\index{proteins}

Taste buds on the tongue, soft palate, and throat contain specialised receptor
cells that detect five well-established basic taste qualities: sweet, salty,
sour, bitter, and umami (savoury). Salty and sour tastes rely mainly on ion
channels, whereas sweet, bitter, and umami rely on G-protein-coupled receptor
proteins, converting chemical stimuli into nerve signals carried via cranial
nerves VII, IX, and X to the brainstem's nucleus tractus solitarius and onward
to the cortex.

Most of what people experience as rich ``flavour'', however, comes from smell
rather than taste alone. As food is chewed and swallowed, volatile aroma
molecules travel upward behind the soft palate into the nasal cavity---a route
called \emph{retronasal olfaction}---where they activate olfactory receptors.
The brain combines these taste and smell signals with texture, temperature, and
trigeminal sensations such as chilli heat or menthol cooling. This is why a
blocked nose during a cold makes food seem bland even when basic taste buds
still work. A sudden, unexplained loss of smell or taste, especially without
nasal congestion or when persistent, should be medically assessed.

\section{Why do ears ring?}\index{earwax}\index{hearing}\index{tinnitus}

Tinnitus is the perception of ringing, buzzing, hissing, or another sound when
there is no corresponding external source. It is a symptom, not one single
disease. It commonly accompanies hearing loss or loud-noise exposure, but
earwax, middle- or inner-ear disorders, jaw or neck problems, head injury, and
some ototoxic medicines (such as aminoglycoside antibiotics, cisplatin, loop
diuretics, or high-dose salicylates) can also contribute. When cochlear hair
cells or auditory nerve synapses are damaged, reduced input can trigger
maladaptive increases in central auditory gain, helping explain why the sound
can persist after the original trigger has gone.

Most tinnitus is subjective, meaning only the affected person hears it. Rarely,
sound generated by nearby tissues may be detectable by another person; this is
called objective tinnitus. Somatic tinnitus is a different term, usually used
when tinnitus is influenced by signals from the jaw, neck, or other body tissues;
it is often still subjective. A rhythmic sound that follows the heartbeat is
pulsatile tinnitus and may reflect blood-flow or another vascular problem, so it
deserves medical assessment. Brief episodes often settle on their own, but
persistent, one-sided, or distressing tinnitus should be evaluated. Sudden
hearing loss, severe vertigo, facial weakness, or other new neurological symptoms
require urgent medical care.

\section{How do we keep our balance, and what does the inner ear contribute?}\index{inner ear}\index{brainstem}\index{brain}\index{calcium}\index{hair}\index{headaches}

Balance depends on the integration of three partly independent information
systems: the vestibular organs of the inner ear, vision, and somatosensory
signals from the skin, muscles, tendons, and joints. The three fluid-filled
semicircular canals contain sensory hair cells that detect rotational
acceleration in different planes. The utricle and saccule, the otolith organs,
detect linear acceleration and head tilt relative to gravity using hair cells
and tiny calcium-carbonate crystals. Their signals reflect the direction of
gravity and movement of the head, rather than acting as a simple separate
``balance organ.''

Brainstem and cerebellar circuits combine these signals with information from
higher brain regions and rapidly adjust eye movements, muscle tone, and posture.
The vestibulo-ocular reflex keeps a visual target stable while the head moves,
and postural reflexes help keep the body's centre of mass over its support. If
the signals conflict---for example, after spinning or when vision disagrees
with the inner ear---the result may be vertigo (a false sense of movement),
light-headedness, visual instability, or unsteadiness. The brain can adapt to
some gradual changes, but persistent symptoms should be assessed. Sudden severe
imbalance accompanied by weakness, numbness, trouble speaking, double vision,
fainting, or a severe headache requires urgent medical care.

\section{Why does music sometimes give us chills?}\index{nervous system!autonomic}\index{dopamine}\index{goosebumps}\index{heart}\index{memory}\index{shivering}

Music-induced chills, or \emph{frisson}, are brief episodes of strong emotional
and physiological arousal. Music can build expectations and then satisfy,
delay, or violate them; changes in harmony, rhythm, timbre, or loudness can
therefore feel especially meaningful, surprising, or rewarding. Personal
memories, lyrics, performance context, and social meaning can also trigger the
response. Studies link
highly pleasurable passages with coordinated activity in auditory, emotion,
and reward networks—including the ventral striatum (nucleus accumbens), insula,
amygdala, and dopamine-related circuits—though dopamine is not
the sole explanation.

Autonomic activation—engaging sympathetic fibres to the skin's arrector pili
muscles—may produce goosebumps, shivering, changes in heart rate,
and altered skin conductance. These reactions can occur without cold, although
goosebumps themselves are also part of an old thermoregulatory reflex. Many
people do not experience musical chills even when they enjoy music. The
response depends on familiarity, personal associations, musical training,
mood, personality, and culture. Chills may accompany pleasure, awe, sadness,
fear, or other intense states, so they do not prove that a piece is objectively
better or that one universal mechanism is involved. Their precise neural
mechanism and evolutionary history remain under study.

\section{Why do we experience deja vu?}\index{brain}\index{memory}\index{seizures}\index{stomach}

Déjà vu (French for ``already seen'') is the feeling that a present situation is
familiar even though one believes it is probably new and cannot retrieve a
specific earlier memory. One explanation is that familiarity signals in medial temporal structures
(such as the perirhinal or rhinal cortex) are briefly activated without the
matching hippocampal process of conscious recollection. It can also occur when a
current scene subtly resembles an earlier, poorly remembered experience, or when
attention and memory processes are briefly out of sync. It is therefore a transient familiarity signal, not evidence that the
present event has literally happened before. The precise mechanism remains under
investigation.

Occasional, brief déjà vu without impaired awareness is common and is not evidence
of precognition, a past life, or a hidden literal memory. Recurrent, prolonged,
or newly changing
episodes---especially those accompanied by a rising stomach sensation, unusual
smell, intense fear, staring, loss of awareness, confusion, or other neurological
symptoms---deserve medical evaluation. Déjà vu can sometimes occur as an aura of
a focal temporal-lobe seizure, but ordinary déjà vu alone does not establish
that diagnosis.

\section{Why do memories become distorted?}\index{memory}\index{brain}\index{stress}

Memory is reconstructed rather than replayed like a video, but reconstruction
does not mean that every memory is false. People may retain the broad meaning
of an event while losing or altering specific details. Attention and emotion
influence what is encoded (with amygdala and stress-hormone signals such as
cortisol modulating hippocampal--neocortical consolidation); time, interference,
and forgetting then affect what remains available. During recall, the brain
combines stored details with expectations, general knowledge, later information,
and the current context. Ordinary memory distortion differs from clinical
\emph{confabulation}—unintended fabricated memories seen in certain brain
injuries or severe amnestic syndromes such as Wernicke--Korsakoff syndrome.
Repeated retrieval can strengthen access to a memory, but it can also make the
memory temporarily labile, allowing it to be updated during reconsolidation or
contributing to a source error---for example, confusing something heard later
with something actually seen.

Leading questions, suggestion, imagination, social discussion, and stereotypes
can therefore distort recall or contribute to false memories without deliberate
lying. Emotion may preserve some central details while impairing peripheral
ones, and the effect of stress depends on its timing and intensity. A vivid or
highly confident memory is not necessarily accurate: confidence can itself be
influenced by repetition, feedback, and the wording of a question.

\section{Why do we get motion sickness?}\index{motion sickness}\index{brainstem}\index{brain}\index{dehydration}\index{inner ear}\index{joints}

Motion sickness occurs when the brain receives conflicting information about
movement and orientation, especially when incoming signals do not match its
expectations. For example, the inner ear may detect acceleration and rotation
in a moving car while the eyes focus on a stationary book; signals from muscles
and joints may add further information. Brainstem and higher networks respond
to this mismatch with nausea, pallor, sweating, dizziness, and sometimes
vomiting. Susceptibility also varies with genetics, age, migraine tendency,
expectation, anxiety, and previous adaptation, so the same motion does not affect
everyone equally. The idea that such conflict resembles the effects of a
neurotoxin is an evolutionary hypothesis, not a demonstrated immediate cause.

Some people reduce symptoms by looking at the horizon or a distant stable
object, sitting where motion is smaller, facing the direction of travel,
avoiding reading, and getting cool airflow. Gradual repeated exposure can also
help the brain adapt, although habituation is not immediate and does not work
for everyone. Preventive medicines—typically targeting muscarinic acetylcholine pathways (such
as scopolamine) or H1 histamine receptors (such as sedating antihistamines)—are
generally most useful before travel and can cause drowsiness, dry mouth, or
other adverse effects; suitability also depends on age, pregnancy, medical
conditions, and other medicines, so a pharmacist or clinician can help choose
an appropriate option. Persistent vomiting, signs of severe
dehydration, or dizziness that occurs without motion should be medically
assessed.

\section{Why do we feel nauseated and why do we vomit?}\index{nausea and vomiting}\index{nervous system!autonomic}\index{brainstem}\index{brain}\index{dehydration}\index{diaphragm}

Vomiting (emesis) is a coordinated reflex that expels stomach contents through
the oesophagus and mouth. It may protect against some swallowed irritants, but
it is also a symptom of many illnesses and does not reliably remove a harmful
substance. Signals from the gut, vestibular system, blood, and higher brain
areas converge on a brainstem vomiting network. The area postrema acts as a
chemoreceptor trigger zone that can detect certain circulating chemicals, while
other brainstem circuits coordinate the final reflex.

During retching and vomiting, the stomach and upper intestine change their
movement, the lower oesophageal sphincter relaxes, the airway closes briefly,
and the diaphragm and abdominal muscles contract to raise pressure and force
contents out. Nausea is the unpleasant, centrally coordinated sensation that may
come before vomiting, but it can also occur alone; vomiting can occasionally
occur with little preceding nausea. Autonomic changes may cause pallor,
sweating, excess saliva, or altered heart rate.

Infections, food poisoning, medicines, pregnancy, motion, migraine, pain,
reflux, obstruction, metabolic disorders, anxiety, and many other conditions
can trigger nausea or vomiting. Repeated or severe vomiting can cause
aspiration, dehydration, and electrolyte losses. Vomiting is not a dependable
way to remove a poison, and a person should not induce it unless specifically
directed by a poison-control service or clinician. Blood or coffee-ground vomit,
green (bilious) vomit, severe abdominal pain, suspected poisoning, inability
to keep fluids down, neurological symptoms, or serious dehydration requires
urgent medical advice.

\section{Why do joints crack?}\index{joints}\index{cartilage}\index{pain}

Joint sounds have several possible sources. When joint surfaces are rapidly
separated, pressure in synovial fluid changes and dissolved gases can form or
collapse a cavity---a process called cavitation---that may produce a sharp
crack. A tendon or ligament can also move suddenly over a nearby prominence,
while rough cartilage or another irregular surface can produce a grating or
crackling sound called crepitus. A single joint may produce several of these
sounds for different reasons. A joint that has cavitated often cannot be cracked again immediately because the
dissolved gases take roughly 15--20 minutes to re-dissolve into the synovial
fluid and restore equilibrium.

Painless cracking, especially of the knuckles, is usually harmless and repeated
knuckle cracking has not been shown to cause arthritis, although habitual
forceful pulling can occasionally strain surrounding soft tissues. The sound
alone cannot reliably identify its cause, and even rough crepitus does not by
itself prove cartilage damage. Pain, swelling, warmth, locking, reduced
movement, or instability should be medically examined.

\section{Why do fingers wrinkle in water?}\index{nervous system!autonomic}\index{pain}\index{skin}

After several minutes in water, the hairless (glabrous) skin of the fingers,
palms, toes, and soles can develop ridges and grooves. This response is
confined to those surfaces because their outer stratum corneum is thick and
tightly anchored to underlying tissue, and their digit pulp contains
specialised vasculature. Prolonged water exposure hydrates and slightly swells
the outer stratum corneum, but the characteristic ridges are produced mainly by
an autonomic reflex: sympathetic nerves constrict small blood vessels in the
digit pulp, reducing tissue volume beneath the outer layer and pulling it into
folds. Because the response depends on intact sympathetic nerves, it is reduced
or absent after nerve injury---a feature clinicians have occasionally used as a
simple bedside clue to peripheral nerve function.

The wrinkles may channel water away and improve handling of some wet objects,
but the size and consistency of any grip advantage varies with the task and
experimental conditions. Their evolutionary importance is therefore plausible
but not settled. The skin normally returns to its smooth appearance after
drying; persistent wrinkling, colour change, pain, or numbness is a different
problem and may warrant assessment.

\section{Why do bruises change colour?}\index{bruises}\index{haemoglobin}\index{skin}

A bruise (contusion) forms when small blood vessels break and leak blood into
the tissues beneath unbroken skin. It may look reddish at first and then dark
blue or purple as blood accumulates. Macrophages gradually remove the leaked
blood and break down haemoglobin; biliverdin can contribute to green tones,
while bilirubin, hemosiderin, and other chemical changes can contribute to
yellow or brown tones as the bruise fades. These processes overlap, rather than
following a fixed timetable.

The colour sequence and timing are approximate, not a reliable clock, and bruise
colour alone cannot reliably establish when the injury occurred. Most bruises
fade in about two weeks, but deeper or larger bruises can last much longer.
Skin tone, depth, lighting, the amount of bleeding, and the location of the
bruise affect its appearance, and medicines such as anticoagulants or
antiplatelet drugs can make bruising easier or more extensive. Frequent
unexplained bruising, rapidly expanding bruising, unusual bleeding, or a bruise
after a significant head injury should be medically assessed.

\section{Why do scars itch?}\index{itch (pruritus)}\index{scars}\index{collagen}\index{fever}\index{histamine}\index{skin}

Healing skin can itch because the wound is dry or tight, inflammatory and
immune signals activate sensory nerves, and nerve fibres are regrowing or have
become unusually sensitive. A scar's altered collagen, reduced skin appendages,
and changed sensory nerve endings can also alter sensation, sometimes producing
a neuropathic itch that is felt as tingling, crawling, or burning. Histamine may
contribute, but it is not the only chemical pathway involved, so antihistamines
would not necessarily address every scar itch. Itching often settles as healing
and remodelling progress, but it may persist when the scar
remains raised or overactive.

A hypertrophic scar is raised but generally stays within the original wound
boundaries; a keloid extends beyond them. Either can itch, burn, hurt, or
restrict movement. Some redness and tenderness can be normal early in healing,
but persistent or worsening itch, rapid growth, spreading redness, warmth,
drainage, or fever should be assessed rather than repeatedly scratched, because
these findings can indicate infection or a scar requiring specific treatment.

\section{How does the skin sense touch, pressure, and temperature?}\index{skin}\index{touch}\index{brainstem}\index{hair}\index{mechanoreceptors}\index{pain}

The skin contains several distinct types of sensory nerve endings rather than
one all-purpose sensor. Specialised \emph{mechanoreceptors} detect physical
deformation---such as light touch, skin stretch, steady pressure, vibration,
and hair movement---by opening mechanically gated ion channels. Rapidly adapting
receptors (such as Meissner corpuscles for light touch and Pacinian corpuscles
for deep vibration) respond best to changing contact, whereas slowly adapting
receptors (such as Merkel discs and Ruffini endings) signal sustained pressure,
fine edges, and skin stretch. Separate \emph{thermoreceptors} use
temperature-sensitive ion channels to detect warming and cooling, while extreme
heat or cold activates pain-sensing fibres.

Signals from touch and pressure receptors travel along fast, myelinated A$\beta$
nerves via the spinal dorsal column--medial lemniscus pathway to the brainstem
and somatosensory cortex, whereas temperature and pain signals travel via
thinly myelinated A$\delta$ and unmyelinated C fibres along the spinothalamic
tract. Regions
with dense receptor spacing, such as the fingertips and lips, occupy much larger
areas of the sensory cortex than the back, giving them far finer two-point
discrimination. New, persistent numbness, tingling, or loss of temperature
sensation---which increases the risk of unnoticed burns or injuries---should be
medically assessed.

\section{Why do we shed skin?}\index{skin}\index{skin!epidermis}\index{inflammation}\index{keratin}\index{proteins}

The epidermis is continually renewed. Keratinocytes divide mainly in its basal
layer, then migrate upward through the epidermal layers—a full turnover cycle
taking roughly 4 to 6 weeks—flatten, fill with the tough protein keratin, lose
their nuclei and organelles, and become dead protective cells in the outermost
\emph{stratum corneum} called corneocytes. Enzymes gradually weaken the connections between surface
corneocytes, allowing them to separate in a regulated process called
desquamation while new cells replace them from below. This is the controlled
loss of superficial epidermal cells, not loss of the entire skin.

This balanced turnover helps maintain a tough, flexible barrier against water
loss, chemicals, and microbes. The rate varies by body site and changes with
age and skin health. Most shedding occurs as tiny, invisible particles,
although it becomes more obvious with dry skin, dandruff, sunburn, eczema, or
other inflammatory disorders. Visible scaling may reflect altered cell cohesion,
barrier damage, or inflammation as well as increased turnover. Excessive,
painful, or persistently abnormal scaling may need medical assessment.

\section{Why do fingernails grow faster than toenails?}\index{nails}\index{hormones}\index{keratin}\index{pain}\index{skin}

Nails grow from the living nail matrix beneath the skin at their base. Matrix
cells divide, become keratinized, and push the existing plate forward; the nail
plate itself is made of hard keratin and is not living tissue. Fingernails often
grow about 3 millimetres per month, whereas toenails often grow about
1--1.5 millimetres per month. Thus, replacing a fingernail may take roughly
4--6 months, while replacing a toenail may take about 12--18 months.

The difference probably reflects several local factors rather than one proven
cause. Fingernails may receive different blood flow and temperature exposure,
and fingers experience more frequent movement and minor mechanical stimulation
than toes. Matrix biology, circulation, and mechanical exposure may all
contribute, but circulation alone is not a complete explanation. Growth also
varies with age, genetics, hormones, nutrition, illness, medicines, season, and
the individual nail; the dominant hand may grow somewhat faster. Cutting or
filing the free edge does not make the matrix produce nail faster. These are
population averages, not a precise health test. New pain, swelling,
detachment, marked discoloration, or a persistent change in one nail's growth
should be assessed by a clinician, particularly after injury or when infection
is possible.

\section{Why does sweat smell?}\index{sweating}\index{metabolism}\index{proteins}\index{skin}

Fresh eccrine sweat is mostly water and salts and usually has little smell. In
the armpits and groin, apocrine glands—which become active at puberty under
androgen influence—release a more complex fluid containing proteins and lipids;
these secretions are initially not strongly odorous. Skin microbes, especially
bacteria and sometimes yeasts, metabolize them and other skin substances into
volatile molecules that produce body odour (\emph{bromhidrosis}). The final odour
depends on the mixture of secretions, skin microbes, body site, genetics, and
environment; eccrine sweat can also carry odorous compounds after microbial
breakdown. Feet can also
smell because warm, enclosed footwear keeps the skin moist and encourages
microbial metabolism.

Washing, drying the skin, and changing breathable clothing can reduce odour. A
deodorant may mask odour or reduce the microbes and compounds that produce it,
whereas an antiperspirant reduces sweating; some products do both. Persistent
excessive sweating, a sudden unusual systemic odour (such as a sweet fruity
scent in diabetic ketoacidosis or an ammonia-like odour in severe kidney
failure), or odour accompanied by skin changes or other symptoms merits medical
advice.

\section{Why do we have earwax?}\index{earwax}\index{hearing}\index{pain}\index{skin}

Earwax, or cerumen, is a mixture of secretions from glands in the outer ear
canal and shed skin cells. It coats and lubricates the canal, helps keep its
skin from drying and cracking, traps dust and small particles, and has mildly
acidic and antimicrobial properties. The canal is normally self-cleaning:
skin gradually migrates outward, aided partly by jaw movements, carrying wax
and old cells toward the opening.

The outer ear can be wiped gently, but objects should not be inserted into the
canal. Cotton swabs may push wax deeper, irritate the skin, or injure the
eardrum, while excessive cleaning can remove protective wax and cause irritation.
Impacted wax that causes pain, discharge, ringing, dizziness, or hearing loss
should be evaluated. If removal is needed, a clinician or pharmacist can
recommend a safe method; drops and irrigation are not suitable for every person,
particularly when the eardrum may be perforated, ear tubes are present, or there
has been ear surgery.

\section{Why do teeth not grow back?}\index{teeth}\index{embryonic development}\index{regeneration}

Human teeth are built during embryonic development and childhood through
carefully coordinated signals between the oral epithelium and underlying
mesenchyme. Humans usually form two dentitions---deciduous teeth followed by
permanent teeth. After the permanent tooth buds have developed, the organized
dental-lamina tissue and developmental signals needed to initiate another
complete tooth are generally no longer available. Rare extra teeth can occur,
but a lost permanent tooth is not normally replaced as it is in some animals
that continually make new teeth.

The enamel-forming cells, called ameloblasts, are lost when a tooth erupts, so
enamel cannot regenerate itself after substantial loss. Early mineral loss can
sometimes be repaired by remineralisation, but this does not rebuild a cavity.
Living cells called \emph{odontoblasts} inside a tooth can produce a limited
layer of tertiary dentin in response to injury, helping protect the pulp. These repair processes do not
replace a whole tooth, restore a missing enamel layer, or reliably recover the
original shape and function after major damage. Biological tooth regeneration
remains an area of research.

\section{Why do wisdom teeth cause problems?}\index{teeth}\index{bones}\index{inflammation}\index{pain}

Wisdom teeth are the third molars, usually erupting in late adolescence or
early adulthood. They are more likely to become impacted when there is not
enough space, when they are angled against the neighbouring molar, or when
bone or gum blocks their path. A proposed evolutionary explanation is that
changes in jaw growth and diet left less room for a late-erupting tooth, but
this is not the only explanation and jaw size and tooth position vary greatly
among individuals.

An impacted tooth can remain symptom-free, but a partially erupted or impacted
tooth can trap plaque and food, causing
pericoronitis (inflammation around the tooth), decay, gum disease, an abscess,
damage to the neighbouring tooth, or occasionally a cyst. Many impacted teeth
remain disease-free, and a healthy, functional, cleanable wisdom tooth does
not automatically require extraction. Removal is generally considered when
there is disease, repeated infection, damage, impaired function, or another
clear dental reason. Extraction itself can cause pain, infection, bleeding,
nerve injury, or damage to nearby structures, so the expected benefit should be
balanced against those risks. The decision should be based on symptoms,
examination, risk assessment, and appropriate imaging rather than impaction alone;
an apparently healthy tooth may instead be monitored.

\chapter{Organs and Internal Regulation}

\section{Why can the liver regenerate while the heart and brain regenerate poorly?}\index{brain}\index{heart}\index{liver}\index{regeneration}\index{collagen}\index{muscles}

Most mature liver cells, especially hepatocytes, normally remain in a resting
state (the G0 phase) but can re-enter the cell cycle when tissue is lost.
Signals from injured cells, immune cells, blood vessels, and the extracellular
framework coordinate this response, allowing the remaining hepatocytes to
enlarge and divide.

After a suitable partial surgical removal (resection)—even up to roughly
two-thirds of a healthy liver—or limited injury, the liver can restore
much of its mass and function.

Adult heart-muscle cells have very limited renewal (roughly 1\% per year or
less in adulthood), usually not enough to replace large losses after a major
injury; collagen-rich scar tissue stabilises the heart but cannot contract.
Some brain regions show limited cell generation, but most mature neurons rarely
divide, glial scarring limits axonal regrowth, and the brain's complex
connections are difficult to recreate. Surviving heart cells and neural
networks can remodel, and the brain can use neuroplasticity (rewiring existing
connections) to support partial functional recovery without replacing all lost
cells. The liver's ability is therefore unusually strong,
not a general property of all organs.

\section{Why do we need two kidneys?}\index{kidneys}\index{blood pressure}\index{diabetes}\index{erythropoietin}\index{hormones}\index{hypertension}

Two kidneys provide reserve capacity—receiving roughly 20--25\% of resting
cardiac output and maintaining a combined glomerular filtration rate (GFR) of
about 90--120 mL/min in young healthy adults—for filtering wastes and excess
water, maintaining salt and acid--base balance, regulating blood pressure
through fluid handling and hormones such as renin, and making hormones such as
erythropoietin and the active form of vitamin D.

A healthy person can often live normally with one functioning kidney because
the remaining kidney enlarges and increases its filtration (compensatory
hypertrophy and hyperfiltration, often reaching roughly 70--80\% of two-kidney
baseline GFR). This adaptation is useful but not unlimited:
a solitary kidney has less reserve and may be more vulnerable to high blood
pressure, protein in the urine, or later kidney disease, particularly if it is
itself injured or the person has diabetes or hypertension. Appropriate medical
follow-up helps protect the remaining kidney.

\section{Why is urine yellow?}\index{urine}\index{fever}\index{kidneys}\index{pain}\index{blood!red blood cells}\index{vitamins}

Urine is usually yellow because it contains urobilin and related pigments,
often grouped clinically under the term urochrome, formed
ultimately from haem (the iron-containing pigment in haemoglobin) in old red
blood cells. Haem is converted to bilirubin, which reaches the intestine and is
changed largely into urobilinogen and other pigments. Some urobilinogen is
reabsorbed, returned to the blood, filtered by the kidneys, and oxidized to
urobilin; this contributes to the yellow colour of urine. With typical adult
urine output of about 1--2 litres per day, the amount of water changes how
strongly the pigment is seen: dilute urine is nearly colourless or pale yellow,
whereas concentrated urine is darker yellow or amber.

Food, medicines (such as bright-yellow riboflavin/vitamin B2 or orange-red
rifampicin), vitamin supplements, and disease can also alter urine colour—for
example, true blood in the urine (\emph{haematuria}) vs. dark tea-coloured
muscle pigment (\emph{myoglobinuria}). Red, cola-coloured, persistently cloudy,
or unusually dark urine cannot be diagnosed by colour alone and should be
assessed, especially when accompanied by pain, fever, jaundice (yellowing of
the skin or eyes), or reduced urination. The significance depends on
diet, medicines, hydration, laboratory findings, and other symptoms.

\section{How do kidneys make urine?}\index{kidneys}\index{urine}\index{bladder}\index{glucose}\index{hormones}\index{pelvis}

Each kidney contains roughly a million microscopic filtering units called
\emph{nephrons}, although the number varies. In each nephron, blood enters a
filtering tuft (\emph{renal corpuscle}), where pressure drives water and small
dissolved substances across the glomerular filtration barrier into Bowman's
space while normally retaining blood cells and most proteins. Together, the two
kidneys filter roughly 180 litres of fluid per day, almost 99\% of which is
reclaimed along the renal tubule (passing sequentially through the proximal
convoluted tubule, loop of Henle, distal convoluted tubule, and collecting
duct). Along these segments, useful water and solutes such as glucose and
sodium are selectively reabsorbed into the blood, while additional ions,
wastes, and some medicines are secreted from blood into the tubule.

The loop of Henle helps establish a salty concentration gradient in the kidney's
inner core (\emph{medulla}), while surrounding vessels help preserve it. The
collecting duct uses that gradient to make urine more dilute or concentrated.
Antidiuretic hormone mainly adjusts water reabsorption in the collecting duct,
while aldosterone promotes sodium reabsorption and potassium secretion in later
parts of the nephron. The final 1--2 litres of daily urine drain through the
collecting ducts to the renal pelvis (the kidney's central funnel) and ureters,
are stored in the bladder, and leave through the urethra. The final urine
contains water and wastes such as urea, creatinine (a muscle-breakdown marker
clinically used to estimate kidney filtration), excess ions, and drug
metabolites; finding significant glucose or protein in urine can be an early
clue to diabetes or glomerular disease.

\section{How does the bladder store and release urine?}\index{bladder}\index{urine}\index{nervous system!autonomic}\index{brainstem}\index{fever}\index{kidneys}

The urinary bladder is a muscular, expandable reservoir that stores urine
arriving continuously from the kidneys via the ureters. During filling,
sympathetic nerve signals keep the smooth muscle of the bladder wall (the
\emph{detrusor}) relaxed to accommodate increasing volume at low pressure, while
the internal urethral sphincter (involuntary smooth muscle) and external
urethral sphincter (voluntary skeletal muscle supported by the pelvic floor)
stay closed to maintain continence.

As the bladder stretches, sensory nerves send signals to the sacral spinal cord
and onward to the pontine micturition centre and cerebral cortex, creating the
conscious urge to urinate. At an appropriate time, voluntary control allows
parasympathetic pathways to contract the detrusor while the sphincters and
pelvic floor relax, emptying urine through the urethra. Common problems such as
overactive bladder, prostate enlargement, pelvic-floor weakness, or nerve
injury can disrupt this coordination. Blood in the urine, painful urination,
fever or flank pain, sudden inability to pass urine, or new incontinence
accompanied by severe back pain, leg weakness, or numbness requires urgent
medical assessment.

\section{What does the spleen do?}\index{spleen}\index{antibodies}\index{bone marrow}\index{bones}\index{fever}\index{liver}

The spleen is both a blood-filtering and an immune organ. Its red pulp removes
old or damaged red blood cells and recycles useful components such as iron; its
white pulp contains lymphocytes and antigen-presenting cells (immune sentinel
cells that display fragments of microbes to T cells) that help detect
and respond to microbes in the blood. The spleen also stores roughly a third
of the body's platelets and helps remove certain antibody-coated cells and
particles from circulation. During
fetal development, the spleen can also contribute to blood-cell production,
although this is not its main role after birth.

People can survive without a spleen because the liver, bone marrow, lymph nodes,
and other immune tissues can partly compensate, but they do not fully reproduce
its filtering and immune functions (in people without a spleen, red blood cells
often retain tiny nuclear remnants called \emph{Howell--Jolly bodies} on a
blood smear). Loss or absence of splenic function (\emph{asplenia}) increases
the risk of rapidly progressing, sometimes life-threatening infection,
especially from encapsulated bacteria such as pneumococcus, meningococcus, and
\emph{Haemophilus influenzae}. Vaccination reduces this risk but does not remove it;
an individualized prevention plan and prompt medical assessment of fever are
important after splenectomy or with a non-functioning spleen.

\section{Why do we shiver when cold?}\index{shivering}\index{brainstem}\index{brain}\index{hypothalamus}\index{metabolism}\index{muscles}

Shivering is a rapid, involuntary series of skeletal-muscle contractions. Cold
sensors in the skin and the brain signal temperature-regulating circuits in the
hypothalamus and brainstem, which activate this response. Each contraction uses
chemical energy, much of which is released as heat rather than useful external
movement. Sympathetic signals also constrict blood vessels in the skin,
reducing warm blood flow to the surface and slowing heat loss.

Shivering is only one defense. Behavioural responses such as seeking shelter,
adding clothing, and curling up reduce heat loss, while non-shivering
thermogenesis—heat production without muscle contraction, driven by uncoupling
protein 1 (UCP1) in mitochondria-rich brown adipose tissue (brown fat)—contributes
more in infants. Shivering consumes energy, cannot fully offset continued cooling,
and may weaken or stop when core temperature falls into moderate-to-severe
hypothermia (hypothermia begins below 35\,\textdegree C, and shivering often
ceases around 30--32\,\textdegree C). Severe or prolonged cold exposure can
therefore cause hypothermia even if shivering initially makes a person feel
warmer; confusion, unusual drowsiness, poor coordination, or loss of shivering
requires emergency help and careful rewarming.

\section{Why do we get freckles and moles?}\index{freckles and moles}\index{cancer}\index{hormones}\index{melanin}\index{skin}

Freckles (ephelides) are small, flat areas in which pigment-producing cells
(melanocytes) make and transfer more melanin pigment to nearby skin cells,
usually in people with inherited sun-sensitive traits. Ultraviolet exposure
makes them darker, and they often fade when exposure decreases; the number of
melanocytes is not usually increased (unlike sun-induced \emph{solar
lentigines}, or ``age spots,'' which involve increased melanocytes and persist).
Moles, or melanocytic nevi, are usually benign localized clusters
or growths of melanocytes. They may be present from birth or develop during
childhood, adolescence, or adulthood. Genetics, hormones, and sun exposure can
influence their number and appearance, but sunlight is not the sole cause.
Freckles are not the same as moles, and their presence often signals greater
sensitivity to ultraviolet exposure; sun protection can reduce additional
darkening and cumulative skin damage but does not eliminate skin-cancer risk.

Most freckles and moles are harmless. A new or changing mole or other pigmented
spot should be checked if it shows the ABCDE warning pattern: asymmetry, an
irregular border, varied colour, a diameter larger than 6 millimetres (about
the size of a pencil eraser) or enlarging size, or evolution in size, shape,
symptoms, or appearance.
different from the others (the ``ugly duckling'' sign) also warrants
dermatologic assessment. These signs are a guide rather than a diagnosis, and a
concerning spot should be assessed even if it does not fit every letter.

\section{Why are some people colour-blind?}\index{colour vision}\index{ageing}\index{brain}\index{chromosomes}\index{optic nerve}\index{retina}

The more precise term is colour-vision deficiency. Normal colour vision begins
with three classes of light-sensing cone cells (photoreceptors)---short- (S,
blue), medium- (M, green), and long-wavelength (L, red) cones---that respond
preferentially to different ranges of light. Retinal and brain circuits compare
the relative strength of these signals, rather than reading each cone in
isolation. Inherited changes can make one cone photopigment abnormal or absent,
most often causing difficulty distinguishing some red and green shades. Many
inherited red--green deficiencies are X-linked, affecting roughly 8\% of XY
males compared with about 0.5\% of XX females, although anyone can inherit or
acquire colour-vision deficiency.
Less common inherited forms affect blue--yellow discrimination or nearly all
colour vision.

Colour-vision deficiency can also be acquired through disease or injury of the
retina, optic nerve, or brain, through cataracts, ageing, or some medicines. It
can affect one eye more than the other and may progress when its cause is an
ongoing eye or neurological disorder.
Most affected people still see colour, but distinguish fewer hues or confuse
particular pairs; seeing only shades of grey is a rare, severe form. New or
progressive colour-vision loss should be evaluated by an eye-care professional.

\section{Why does the voice change during puberty?}\index{puberty}\index{androgens}\index{breathing}\index{hormones}\index{larynx and voice}

Pubertal hormones enlarge the larynx and change the length, mass, and tension
of the vocal folds (the ``vocal cords''). Longer, heavier folds usually vibrate
more slowly, lowering the basic pitch (fundamental frequency---clustering around
roughly 85--180 Hz in adult males and 165--255 Hz in adult females after
puberty); the size of the larynx and the shape of the vocal tract also alter
resonance and voice quality. Pitch often falls as the folds lengthen (roughly
doubling in male puberty) and become heavier, but this is an average tendency,
not a fixed outcome. The average change is often larger with the greater
androgen exposure typical of male puberty, but voices change across all sexes
and hormone patterns, with wide individual variation.

During rapid growth, the nervous system must coordinate breathing, vocal-fold
movement, and a changing larynx. This adjustment can temporarily produce voice
cracks or breaks. Learned vocal habits and changing control of resonance also
influence the final voice. The timing and extent of change vary widely, and a
persistent or painful voice change outside an expected developmental period
should be assessed.

\section{Why do males develop an Adam's apple?}\index{androgens}\index{cartilage}\index{hormones}\index{larynx and voice}\index{lungs}\index{puberty}

The ``Adam's apple'' is the laryngeal prominence, where the two plates of the
thyroid cartilage meet at the front of the larynx. This cartilage and
prominence exist in everyone, but the androgen surge typical of male puberty
causes the larynx to grow substantially larger and the cartilage plates to
meet at a sharper angle (typically around 90 degrees rather than roughly 120
degrees), often making the prominence project further forward. Even so, size and
visibility vary widely with neck anatomy, soft-tissue thickness, genetics, body
build, and hormone exposure, so the prominence is not a reliable indicator of
sex or gender.

The thyroid cartilage forms a protective framework around the larynx and
anchors the vocal folds, which lengthen and thicken during pubertal growth to
produce a deeper voice. The prominence itself does not generate sound; the
vocal folds inside the larynx do that when air from the lungs makes them
vibrate. A new, growing, or tender lump in the front of the neck---such as a
thyroid nodule or cyst---is distinct from the thyroid cartilage and should be
medically assessed.

\section{Why does menopause occur?}\index{menopause}\index{hormones}\index{menstruation}\index{oestrogen}\index{progesterone}\index{sleep}

The ovaries begin life with a finite reserve of follicles. Most are naturally
lost through natural follicle breakdown (\emph{atresia}), and the remaining
follicles gradually become fewer and less responsive. Ovarian oestrogen and
inhibin production therefore become less consistent, ovulation becomes
irregular and eventually stops, progesterone production falls because ovulation
is no longer occurring, and pituitary follicle-stimulating hormone (FSH) rises
as feedback is reduced.

Natural menopause is the final menstrual period---typically occurring between
ages 45 and 55, with an average around 51---and is identified retrospectively
after 12 consecutive months without menstruation when no other cause explains
the absence. The preceding period of fluctuating cycles and
hormone levels is called perimenopause. Ovulation can still occur during
perimenopause, so pregnancy remains possible until menopause has been
established. Hot flushes and night sweats, sleep or mood changes, and vaginal or
urinary symptoms are common but not universal; over the longer term, lower
oestrogen also accelerates bone-density loss (osteoporosis risk) and alters
cardiovascular risk, and troublesome symptoms can be treated. Menopause can also occur earlier after ovary-removing surgery;
chemotherapy, radiation, or some medicines can temporarily or permanently
impair ovarian function. Menstrual cessation before age 40 needs assessment for
primary ovarian insufficiency or another cause, and any bleeding after 12 months
without periods should be medically assessed.

\section{Why can some tumours contain hair, bone, or teeth?}\index{bones}\index{cancer}\index{hair}\index{teeth}\index{cartilage}\index{embryonic development}

Some germ-cell tumours called \emph{teratomas} arise from cells that retain
pluripotent developmental potential (the capacity to form many different cell
types). Their descendant cells can differentiate abnormally into tissues from
several early embryonic cell layers---ectoderm, mesoderm, and endoderm---rather
than forming a normal, organized embryo. A teratoma can therefore contain mature
skin, hair follicles, cartilage, bone, fat, glandular tissue, or tooth-like
structures in an anatomically disorganized mass. These are differentiated
tissues, not a complete developing embryo, and not every germ-cell tumour is a
teratoma. Teratomas occur most often in the ovary or testis, but can also arise
along midline developmental paths (such as the chest or tailbone area) and can
occasionally cause complications such as ovarian torsion; blood markers such as
AFP and hCG help evaluate immature or mixed germ-cell tumours.

Mature teratomas are usually benign, whereas immature teratomas are more likely
to behave malignantly; a teratoma can also contain a separate malignant
component. Thus, the presence of recognizable hair or bone does not by itself
show whether a tumour is harmless. Diagnosis and behaviour depend on pathology,
including microscopic features, as well as location, spread, and clinical
assessment.

\section{Why is the total length of human blood vessels estimated to be so great?}\index{arteries}\index{capillaries}\index{carbon dioxide}\index{oxygen}\index{veins}

The vascular system is so extensive because cells can obtain oxygen and
nutrients, and remove carbon dioxide and wastes, only over short diffusion
distances. Repeated branching brings microscopic capillaries close to most
living tissue, while arteries and veins distribute and return blood through a
controllable network. The network's value depends on proximity to cells and
exchange surface area as well as length; capillary density differs greatly
between tissues, and not every capillary is maximally perfused at the same
moment.

The often-quoted 60,000-mile figure is a rough estimate of the combined length
of arteries, veins, and especially capillaries in an average adult, though more
recent estimates place the total substantially lower (roughly 6,000--12,000
miles, or about 9,000--19,000 kilometres). It is not a
directly measured length or one continuous tube. The true total varies with
body size, tissue type, capillary density, what counts as a vessel, and the
method or assumptions used to estimate it. The number is therefore a popular
order-of-magnitude illustration, not a precise anatomical constant. It should
not be interpreted as the length of a single vessel or as evidence that blood
travels through every capillary at once.

\section{How should the strength of human muscles be compared?}\index{muscles}\index{heart}\index{joints}

There is no scientifically meaningful single winner because ``strongest'' can
mean maximum force, force relative to muscle size, power, endurance, or total
work. These measures should be tested under standardized conditions and with a
specified contraction type, because a muscle's force is not the same as the
external force or joint torque produced after leverage is considered. The
masseter can produce high force relative to its size and contributes to a
powerful bite, but bite force also depends on jaw shape and leverage. The
gluteus maximus and the quadriceps group can generate very large forces for hip
extension, standing, climbing, and jumping; the quadriceps is a group of
muscles rather than one muscle. The heart is exceptional for continuous work,
but its endurance is not directly comparable with the peak force of skeletal
muscle.

The result also depends on muscle architecture, physiological cross-sectional
area, joint leverage, neural activation, training, and the person being tested.
The common claim that the tongue is the strongest muscle has no clear scientific
basis; the tongue is a coordinated group of muscles, and no single test defines
one universal winner.

\section{How did human development produce two lungs but one liver?}\index{liver}\index{lungs}\index{embryonic development}\index{heart}

The respiratory outgrowth of the embryonic foregut (the upper part of the early
digestive tube) divides into right and left lung buds, which then branch into
two airways and two extensive gas-exchange organs. Pairing provides a large
exchange surface, distributes the organs on both sides of the chest around the
heart, and gives some reserve if one lung is damaged. One healthy lung can often
sustain life, but it cannot always provide the same exercise capacity or reserve
as two lungs.

The liver develops mainly as one large organ from a hepatic diverticulum, an
outgrowth of the foregut that grows into surrounding embryonic connective
tissue (the septum transversum) and vitelline veins. It performs integrated metabolic, storage, synthetic,
digestive, and chemical-processing tasks and has considerable functional
reserve. It can restore much of its mass after some injuries, but this does not
mean that it is a second liver in reserve. Its two major lobes, together with
smaller anatomical subdivisions such as the caudate and quadrate lobes, are parts
of one connected organ, not independent livers. The difference reflects
embryological patterning and functional requirements; it does not establish a
universal rule that every vital organ must be duplicated.

\section{Why are our two sides not perfectly symmetrical?}\index{embryonic development}\index{genes}\index{handedness}\index{heart}\index{liver}\index{lungs}

The embryo begins with a broadly bilateral plan, but left--right asymmetry is
actively established during development. Signals involving the embryonic node,
motile cilia, and a broader network of genes including \emph{NODAL} and
\emph{PITX2} help position and shape internal organs: the heart apex, stomach,
and spleen are generally to the left, while most of the liver is to the right.
The two lungs are also not identical---the right normally has three lobes and
the left two---and rare developmental conditions such as situs inversus or
heterotaxy can reverse or disrupt this internal arrangement.

The face, limbs, and other paired structures are not exact mirror images.
Differences in growth, muscle use, posture, handedness, genetics, measurement
position, and small random developmental events can all contribute to external
asymmetry; these factors do not by themselves determine the programmed
placement of internal organs. Mild, long-standing asymmetry is common, whereas
new or rapidly worsening asymmetry---especially with pain, weakness, swelling,
or a neurological symptom---should be medically assessed.

\section{Why do people produce so much saliva?}\index{saliva}\index{antibodies}\index{nervous system!autonomic}\index{calcium}\index{dehydration}\index{digestion}

Saliva is produced by the parotid, submandibular, sublingual, and many minor
glands. It moistens and binds food, lubricates speaking and swallowing, begins
starch digestion, and carries dissolved molecules to taste receptors. Enzymes
can also begin limited fat digestion. Bicarbonate buffers acids, while proteins
and antibodies help protect the mouth and teeth from microbes and decay; calcium
and phosphate can also support remineralisation of early enamel damage. Most
saliva is swallowed and its water and dissolved substances are subsequently
reclaimed by the gut, so normal production does not deplete the body. Continuous
flow also helps clear food and acids between meals.

Autonomic nerves regulate secretion. Chewing, food smells, and anticipation
increase flow; sleep and dehydration reduce it. Nausea can cause a sudden
increase, which lubricates the mouth and may help dilute or buffer acid before
vomiting. Healthy adults produce roughly 0.5--1.5 litres per day, with
substantial variation. Persistent dry mouth or apparent excessive drooling can
have medical, dental, medication-related, or neurological causes; drooling may
reflect increased production, difficulty swallowing, or difficulty keeping the
mouth closed; these possibilities require different evaluations.

\section{Why can the brain feel no pain even though it processes pain from the body?}\index{brain}\index{pain}\index{brainstem}\index{cerebral cortex}\index{headaches}\index{inflammation}

Most brain tissue (\emph{parenchyma}) lacks nociceptors, the sensory nerve
endings that detect potentially damaging stimuli. This does not mean that the
head cannot hurt: the scalp, the membrane covering the skull bones
(\emph{periosteum}), muscles, the protective membranes around the brain
(\emph{meninges}, especially the dura), and selected large intracranial blood
vessels and venous structures are pain-sensitive.
Irritation, inflammation, stretching, or pressure involving these structures
can produce headache, and they are the main sources of pain during
neurosurgery; manipulating brain parenchyma itself is usually not painful. For
this reason, neurosurgeons can perform awake craniotomies---after numbing the
scalp, periosteum, and meninges with local anaesthetic---to map speech or
movement areas safely without causing pain in the brain tissue itself.

Pain from the rest of the body is detected by peripheral nociceptors and sent
through the spinal cord and brainstem to networks including the thalamus and
cerebral cortex, where the experience of pain is constructed and influenced by
attention, emotion, and context. Nociception is detection of potentially
damaging stimuli; it is not identical to conscious pain. Thus, the brain can
process nociceptive signals and construct pain without having nociceptors
throughout its own main tissue.

\section{Why does the heart beat automatically without conscious control?}\index{heart}\index{adrenal glands}\index{adrenaline}\index{nervous system!autonomic}\index{hormones}\index{muscles}

The heart is intrinsically myogenic (meaning it generates its own beat within
muscle tissue): specialized pacemaker cells in the sinoatrial (SA) node in the
right atrium slowly build electrical charge (\emph{depolarize}) on their own
because of their membrane ion currents (typically firing at 60--100 beats per
minute at rest). When they reach threshold, they generate an
electrical impulse that spreads through the atria, is delayed briefly at the
atrioventricular (AV) node to allow ventricular filling, and then travels
through the bundle branches and Purkinje fibres. This electrical sequence
coordinates contraction of the ordinary atrial and ventricular muscle cells.

The autonomic nervous system and circulating hormones continuously adjust this
intrinsic rhythm. Sympathetic activity and adrenaline usually speed it up,
whereas parasympathetic vagal activity (via the vagus nerve) usually slows it;
neither is needed to initiate every normal beat. Other cardiac cells can act as
slower backup pacemakers if the SA node fails (roughly 40--60 beats per minute
at the AV junction, or 20--40 beats per minute in ventricular tissue), but
adequate oxygen, fuel, and electrolyte balance are still required. ``Automatic'' therefore means that the normal rhythm
is generated within the heart, not that the heart can function independently of
the body.

\section{What are the advantages and trade-offs of having two kidneys but one heart?}\index{heart}\index{kidneys}\index{lungs}

Two kidneys filter blood in parallel, providing extra filtering capacity and
partial redundancy if one is damaged or removed. A person can often live with
one kidney---which is why living kidney donation is possible---because the
remaining kidney enlarges and increases filtration, though its reserve is not
unlimited and long-term blood-pressure and kidney-function monitoring remain
important.

The heart is one organ with two synchronised sides. The right side pumps blood
to the lungs, and the left side pumps oxygenated blood to the body. Because
these two circuits are connected in series, one coordinated heart keeps flow
closely matched between them; two separate hearts pumping slightly out of sync
would rapidly cause blood to back up in the lungs or drop systemic pressure.
Because both sides share one organ and depend on each other, failure of either
side can compromise the circulation as a whole. This arrangement reflects
developmental and evolutionary history, not a guarantee that heart disease will
be harmless.

\chapter{Immunity, Cancer, and Sleep}

\section{Why can a person survive with one kidney, one lung, or part of the liver?}\index{kidneys}\index{liver}\index{lungs}\index{blood pressure}\index{breathing}\index{diabetes}

These organs have reserve capacity, so losing some tissue does not necessarily
make their essential functions fail immediately. After one kidney is lost, the
remaining kidney usually enlarges and its individual filtering units
(\emph{nephrons}) filter more, often preserving roughly 65--75\% of two-kidney
filtration capacity. This compensation can maintain adequate filtration, but it
is not unlimited; blood pressure, diabetes, pre-existing kidney disease, and
other factors affect long-term risk.

After removal of one lung (pneumonectomy), the remaining lung, chest mechanics,
circulation, and breathing muscles can make more effective use of the available
gas-exchange surface; the remaining lung may also expand, especially during
growth, though clinicians monitor for mediastinal shift or pulmonary
hypertension. It does not regrow a second lung, however, so while one lung may
support ordinary activity, maximal exercise capacity or reserve may be reduced.

After a limited surgical removal (\emph{resection}) of a healthy liver—leaving
at least roughly 25--30\% of healthy liver volume, or up to about 70\%
resected—surviving liver cells (\emph{hepatocytes}), stimulated in part by
increased portal blood-flow signalling, can enlarge and divide, restoring much
of the liver's mass and function. This is
restoration of tissue mass, not necessarily an exact reconstruction of the
missing lobe or its original microscopic arrangement.

These outcomes depend on age and health, the amount and quality of tissue left,
and the reason for its loss. Cirrhosis, emphysema, vascular disease, or
pre-existing kidney damage can greatly reduce the available reserve; medical
follow-up remains important even when daily activities are possible.

\section{Why does the immune system sometimes attack the body's own cells?}\index{immune system}\index{autoimmune disease}\index{hormones}\index{proteins}

The immune system normally maintains self-tolerance through central and
peripheral layers of control. During development in the thymus and bone marrow
(\emph{central tolerance}), many strongly self-reactive T and B immune cells
are deleted, edited, or silenced; in the tissues (\emph{peripheral tolerance}),
other self-reactive cells are restrained by regulatory T cells (specialised
peacekeeping cells), inhibitory signals, lack of co-stimulation, or limited
exposure to particular target molecules (\emph{antigens}). These safeguards are
effective but not perfect, so some self-reactive cells normally remain in
healthy people.

In a susceptible person, genetic variants (most notably HLA/MHC immune-gene
types), infections, tissue injury, altered self-proteins, hormones (many
autoimmune diseases, such as lupus or thyroiditis, occur more often in women),
or environmental exposures can provide activation signals that break or bypass
these restraints.

Activated T cells, autoantibodies, complement (a cascade of blood proteins that
punches holes in or marks targets), and inflammatory cells may then damage
tissues or alter their function—producing conditions such as type 1 diabetes,
rheumatoid arthritis, or lupus. Autoantibodies can destroy cells, form
inflammatory antibody--antigen clumps (\emph{immune complexes}), or block or
stimulate receptors without directly killing the target cell. Transient self-reactivity is not
necessarily autoimmune disease: disease generally means that self-directed
immunity causes persistent or clinically important injury. The initiating
trigger remains unknown for many conditions.

\section{How does the immune system distinguish self from non-self?}\index{immune system}\index{antibodies}\index{autoimmune disease}\index{cancer}\index{proteins}

There is no single ``self detector.'' During development in the thymus (including AIRE-guided screening against
body-wide proteins), many strongly self-reactive T and B cells are deleted or
restrained. In the tissues, T cells recognise short protein fragments
(\emph{peptides}) displayed on cell-surface display trays called major
histocompatibility complex (MHC) molecules: MHC class I is widely displayed by
almost all nucleus-containing (\emph{nucleated}) cells and activates CD8 killer
T cells, whereas MHC class II is mainly displayed by specialised sentinel
(\emph{professional antigen-presenting}) cells and activates CD4 helper T cells.
MHC is not a simple identity card:
healthy cells display many normal self peptides, whereas infected or transformed
cells may display abnormal or newly produced peptides. B-cell receptors and
antibodies can also bind three-dimensional shapes on intact molecules, while
regulatory cells and inhibitory signals normally prevent responses to ordinary
self molecules. Natural-killer (NK) cells also monitor patterns such as reduced
or altered MHC class I and can respond to ``missing-self'' signals, especially in
infected or transformed cells.

Innate immune cells also use built-in pattern-recognition receptors (such as
Toll-like receptors). They respond to conserved microbial signatures
(pathogen-associated molecular patterns, or PAMPs) and stress signals released
by injured or dying self cells (damage-associated molecular patterns, or DAMPs).
Adaptive lymphocytes usually need both antigen recognition and confirmatory
co-stimulatory and cytokine signals before resting (\emph{naive}) T cells become
fully activated. The system is therefore a layered, context-dependent probability
judgment rather than an infallible yes-or-no detector: infections, cancer,
transplants, and autoimmune disease can all expose its limitations.

\section{Why does the immune system normally tolerate intestinal bacteria?}\index{immune system}\index{inflammation}

The immune system does not ignore intestinal bacteria; it contains resident
microbes while limiting damaging inflammation. A mucus layer and tightly joined
epithelial cells keep most organisms in the hollow interior (\emph{lumen}) and
mucus, away from the underlying tissue. Antimicrobial peptides and mucosal
antibodies called secretory immunoglobulin A (IgA) restrict microbial attachment
and penetration without trying to sterilize the intestine. Gut sentinel cells
(dendritic cells and macrophages, sampling microbes via specialised Peyer's
patches and M cells) and peacekeeping regulatory T cells—using calming signals
such as IL-10 and TGF-$\beta$—promote a measured, largely non-destructive
response to many established resident microbes. Microbial metabolites such as
short-chain fatty acids (notably butyrate) can also support epithelial
integrity and immune regulation, although their effects depend on the
metabolite, microbial species, dose, location, and host context.

This balance permits colonization by generally beneficial or harmless microbes
while preserving the ability to attack invasive pathogens. Disrupting this
ecosystem—for instance, when broad-spectrum antibiotics wipe out protective
bacteria and allow opportunistic \emph{Clostridioides difficile} overgrowth—can
trigger severe inflammation. A microbe can be
harmless while contained in the lumen yet harmful after entering tissue, and
resident communities can include opportunistic organisms. Barrier injury,
changes in the microbial community, genetic susceptibility, or defective
immune regulation can allow bacteria or their products to enter tissue and
contribute to infection or chronic intestinal inflammation. The gut is
therefore a site of active tolerance, containment, and surveillance, not
immunological indifference; tolerance is selective and can fail.

\section{Why does inflammation protect us but also cause disease?}\index{inflammation}\index{atherosclerosis}\index{cancer}\index{pain}\index{proteins}\index{scars}

Inflammation is a coordinated response to infection, physical injury, toxins,
or other danger signals. It can be triggered by microbes or occur as non-infectious (\emph{sterile})
inflammation after injury, crystals, or cell damage. Acute inflammation changes
local blood flow and vessel permeability—producing the classic visible signs of
redness (\emph{rubor}), warmth (\emph{calor}), swelling (\emph{tumor}), pain
(\emph{dolor}), and sometimes reduced function (\emph{functio laesa})—recruits
early-responding neutrophils followed by macrophages and plasma proteins, helps
contain or kill invaders, removes damaged material, and initiates repair.
Once the threat is
controlled, the response normally shifts toward resolution: inflammatory-cell
recruitment falls, surviving cells clear debris, and repair and tissue-balance
programs take over. Resolution is an active biological process, not merely the
passive disappearance of inflammation; even after symptoms improve, repair may
leave scar tissue or altered function.

If the response is too intense, directed at harmless or self tissues, or fails
to resolve, the same chemical messengers (\emph{cytokines}), enzymes, reactive
molecules, and fluid that were protective can injure healthy cells. Persistent
inflammation can lead to ongoing pain, tissue destruction, abnormal repair,
scarring (\emph{fibrosis}), impaired organ function, and, in some settings,
increased cancer risk. Conversely, strongly suppressing inflammation with
medicines can relieve disease but also carries a higher risk of infection. It contributes to diseases
including rheumatoid arthritis, asthma, inflammatory bowel disease,
atherosclerosis, and some neurodegenerative conditions; inflammation is often
one contributing process rather than the sole cause. A response can also be
harmful when it becomes widespread, as in severe systemic infection, even if
the original stimulus is localized.

\section{Why does fever occur, and when can it be beneficial?}\index{fever}\index{breathing}\index{cancer}\index{dehydration}\index{seizures}\index{shivering}

Fever is a regulated rise in core temperature. During infection or another
inflammatory condition, cytokines such as interleukin-1, interleukin-6, and
tumour-necrosis factor stimulate prostaglandin E$_2$ pathways that temporarily
raise the hypothalamic set point. The body then constricts skin vessels, reduces
heat loss, and may cause chills and shivering until the higher set point is
reached. This differs from hyperthermia, in which temperature rises without a
raised set point because heat production or heat gain exceeds heat loss;
hyperthermia requires different urgent management.

A modest, short-lived fever may inhibit the growth of some microbes and enhance
aspects of immune activity, although its possible benefit depends on the illness
and is balanced by increased metabolic demand and discomfort. It is not a reason
to ignore concerning symptoms. Fever is a symptom, not proof
that an infection is present, and temperature alone does not reliably indicate
how serious an illness is. Fever reduction is mainly for comfort or particular
clinical circumstances; it does not treat the underlying cause and should be
guided by symptoms, age, and underlying conditions; treating a fever for comfort
does not necessarily interfere with recovery.

Prompt medical assessment is important for a rectal temperature of 38~$^\circ$C
or higher in an infant younger than 3 months, or for fever with confusion,
breathing difficulty, seizure, severe dehydration, or a markedly unwell
appearance. Persistent or very high fever also warrants medical advice.

\section{Why do some wounds heal without scars while others produce permanent scars?}\index{scars}\index{wound healing}\index{blood!clotting}\index{collagen}\index{skin!epidermis}\index{hair}

Wound healing has overlapping stages: clotting stops bleeding; inflammation
removes damaged tissue and helps control microbes; proliferation restores the
surface and builds temporary vascularised repair tissue (\emph{granulation
tissue}); and remodelling reorganizes the collagen over months (shifting from
early type III collagen to stronger type I collagen, though mature scar tissue
typically reaches only about 70--80\% of uninjured skin's tensile strength). If
an injury is limited to the epidermis and its supporting underlying sheet
(\emph{basement membrane}) remains intact, epidermal cells can migrate and
divide to restore nearly normal architecture, leaving little or no visible scar.
Keeping wounds clean, moist, protected from tension and infection, and shielded
from strong sun during healing helps minimise unnecessary scarring. This
does not necessarily mean that microscopic tissue is identical to uninjured
skin.

A deeper wound that disrupts the dermis damages its collagen framework, elastic
fibres, and often hair follicles or sweat glands. Repair must then fill the gap
quickly with collagen-rich fibrous tissue rather than recreate every original
layer and appendage. Thus, a full-thickness adult skin wound generally leaves
some scar. The scar may flatten and fade over months or years, but it usually
remains structurally different and less flexible than normal skin. Wound size
and depth, separated or tense edges, infection, delayed healing, location, age,
genetics, skin type, and tissue context all influence the result. Some fetal
skin and oral-mucosal wounds can heal with less scarring than typical adult
skin, so the outcome is not determined by depth alone. Excessive or disorganized
repair can produce a raised hypertrophic scar that stays within the original
wound margins, or a keloid that extends beyond them.

\section{Why do bones heal whereas damaged cartilage heals poorly?}\index{bones}\index{cartilage}\index{joints}

Bone is a living, usually vascular tissue. After a fracture, blood vessels,
marrow and outer bone-lining (\emph{periosteal}) cells, signalling molecules,
and bone-building (\emph{osteoblasts}) and bone-remodelling (\emph{osteoclasts})
cells enter the injury, while immobilisation and the surrounding framework
provide mechanical guidance. The repair typically forms a temporary soft
callus, followed by a hard callus over roughly 6--12 weeks that is gradually
remodelled into mature, organized bone. The early callus is not instantly
equivalent to uninjured bone. Poor blood supply, infection, excessive movement,
a large gap, poor alignment, or other disease can still cause delayed healing
or non-union.

Cartilage is not one uniform tissue, but most joint-lining (\emph{articular} or
\emph{hyaline}) cartilage has no blood vessels. Its sparse cells
(\emph{chondrocytes}) are embedded in a dense surrounding scaffold
(\emph{extracellular matrix}) rich in type II collagen, and have limited ability
to migrate to an injury, so nutrients and repair signals arrive mainly by
diffusion. Slow matrix turnover and repeated joint loading further limit repair,
which is why damaged joint cartilage can predispose to osteoarthritis over time.
A small, superficial defect may heal only incompletely, although some cartilage
types and peripheral lesions have better potential than central articular
cartilage. If damage reaches the underlying bone, marrow-derived cells can
enter, but the repair is often rougher, type I collagen-rich
\emph{fibrocartilage} rather than the smoother, more durable hyaline cartilage. Other cartilage types differ; for example, only the
outer region of a meniscus has a useful blood supply, so healing is better near
that vascular edge than in its inner region.

\section{Why do humans age if cells can divide and replace themselves?}\index{ageing}\index{genes}\index{inflammation}\index{mitochondria}\index{neurons}\index{proteins}

Cells do not all divide continually, and replacing cells is not the same as
restoring an entire tissue. Many neurons and other specialized cells divide
little or not at all; even renewing tissues depend on stem cells whose numbers,
quality, and supporting tissue microenvironments (\emph{niches}) can change
with age. Turnover can preserve a tissue's broad structure while failing to
remove every molecular defect in surviving cells or perfectly reproduce the
original supportive scaffolding (\emph{extracellular matrix}) and connections.
Replacement cells are also produced within an ageing niche: they
can inherit accumulated changes and encounter altered signals, blood supply,
matrix, and neighbouring cells. Turnover therefore does not reset a tissue to a
young state.

Ageing is therefore a gradual, interacting process rather than the result of
one defect. Changes in genome maintenance, gene regulation, mitochondrial
function, protein quality control, nutrient sensing, and cell-to-cell
communication can accumulate over time. Shortening of protective chromosome ends
(\emph{telomeres}) limits the replicative capacity of some dividing cells (the
Hayflick limit), but it is not an explanation for all ageing. Permanent
cell-cycle arrest (\emph{cellular senescence}) can help prevent damaged cells
from dividing and can support some repair in the short term; persistent
senescent cells, their inflammatory secretions (the senescence-associated
secretory phenotype, or SASP), chronic inflammation, declining stem-cell
function, and weaker clearance of damaged material can later impair tissue
function. Regular physical activity, good sleep, and avoiding smoking help
preserve tissue reserve and reduce clinical frailty, which is distinct from
ordinary healthy ageing. Eventually maintenance and repair no longer keep pace with damage, even though
cell division continues in many tissues. Cell lifespan also reflects tissue
function: intestinal cells exposed to friction and microbes are replaced
quickly, whereas many neurons preserve specialised connections for decades
(though neurons are not immortal, and replacement is limited when they die).

\section{Why does cancer become more common with increasing age?}\index{cancer}\index{ageing}\index{DNA}\index{genes}\index{hormones}\index{inflammation}

Most cancers develop through a multistep process—often requiring several
independent driver alterations (commonly two to eight in many solid tumours).
Over time, long-lived stem and immature precursor (\emph{progenitor}) cell
lineages have more opportunities to acquire cancer-driving DNA changes, and they
accumulate effects from environmental exposures, inflammation, and ordinary
replication. Ageing can also alter gene regulation and the tissue environment:
repair and growth-control signals become less reliable, aged growth-arrested
(\emph{senescent}) cells and chronic inflammation can change neighbouring-cell
behaviour, and immune surveillance may become less effective. These changes
can make it easier for an abnormal clone that has already arisen to survive and
expand—producing the steep, roughly exponential rise in many epithelial cancers
with age. This is a process of \emph{clonal selection} (cells with survival or
growth advantages outcompeting neighbours), while the effects of ageing on
immune surveillance vary among people and tissues. Age does not simply make mutations
appear in every cell.

These mechanisms increase probability; they do not make cancer inevitable.
Inherited variants, tobacco and other exposures, infections, hormones, chance,
and the specific tissue all influence risk. Some cancers are common in
childhood or young adulthood, and the apparent increase in older adults also
partly reflects that more people now survive long enough to develop
age-associated disease and that screening and diagnosis vary by age. Age is the
strongest overall risk factor for cancer, but it is not a universal rule for
every cancer type.

\section{Why does not everyone develop cancer despite many DNA mutations?}\index{cancer}\index{DNA}\index{chromosomes}\index{epigenetics}\index{genes}

DNA damage occurs frequently, but much of it is repaired before becoming a
permanent mutation. Of the mutations that persist, many are harmless background
(\emph{passenger}) changes: they do not alter an important gene, occur in a
region that does not affect cell behaviour, or remain confined to a short-lived
cell or a small clone. Some mutations can give a cell a growth advantage and
produce a small cluster of identical cells (\emph{clonal expansion}) without
causing invasive cancer; others are harmful to the cell and cause it to stop
dividing or die. Thus, a mutation is not by itself evidence of cancer.
Internal safety brakes (cell-cycle checkpoints and permanent growth arrest, or
\emph{senescence}), programmed cell death, tissue barriers, and immune
surveillance provide additional defenses.

Cancer generally requires one or more persistent driver alterations—for
example, disabling tumour-suppressor genes such as \emph{TP53} or
\emph{BRCA1/2}, or activating growth-promoting oncogenes such as \emph{RAS}—that,
in a particular tissue, enable abnormal survival or growth and help the clone
evade growth suppression, cell death, or immune control. These alterations can
include DNA sequence mutations, copy-number changes, chromosome rearrangements,
or chemical gene-switching (\emph{epigenetic}) changes. The number and identity of necessary drivers differ among
cancer types. Competition from neighbouring
cells and immune surveillance may eliminate or contain abnormal cells, but
neither defense is perfect. Even an inherited cancer-predisposing variant
usually raises risk rather than guaranteeing cancer. Having many mutations is
therefore not equivalent to having the biological combination and opportunity
needed for cancer.

\section{Why can cancer arise from almost any tissue but remain rare in some tissues?}\index{cancer}\index{chromosomes}\index{DNA}\index{genes}\index{hormones}\index{inflammation}

Cancer can arise from many tissues because many tissues contain cells or
progenitor cells that can acquire cancer-driving alterations and, if enough
growth controls fail, form a clone. These alterations may involve DNA sequence,
chromosome structure, copy number, or gene regulation. Risk nevertheless differs
greatly among tissues because they vary in the number and lifespan of susceptible
stem or progenitor cells, normal cell
division rates, DNA repair, exposure to carcinogens, hormones, inflammation,
immune surveillance, and the signals and physical architecture of the local
tissue environment. The identity and developmental state of the cell of origin
also matter: the same mutation can have different effects in different cell
types.

Cell division is therefore one contributor, not a complete explanation—and
even strong statistical correlations with stem-cell division counts do not
replace tissue-specific biology. Tissues with relatively few divisions of
mature cells, such as heart muscle (where cardiomyocytes are largely terminally
differentiated), rarely develop primary tumours, whereas tissues with rapid
turnover and high environmental exposure, such as skin, lung, and colon,
develop cancer much more often; even within the gut, small-intestine cancers
are far rarer than colorectal cancers despite abundant cell division. Cancer incidence reflects several stages:
initiation, continued growth, invasion, and eventual detection. Rare cancers
are not impossible; they may require an uncommon combination of genetic changes
and tissue conditions, although some are also difficult to detect or classify
accurately. Thus, low observed incidence can reflect low biological risk,
under-detection, or both.

\section{Why does the body sometimes fail to recognise and destroy cancer cells?}\index{cancer}\index{inflammation}

Cancer cells arise from the body's own cells, so much of their molecular
identity is treated as self. Some cancers display novel protein fragments created by
tumour-specific alterations (\emph{neoantigens}) or abnormally expressed
molecules (tumour-associated antigens) that immune cells can recognize, but the
strength and quality of that recognition vary among tumours and patients. The
immune system may therefore fail to detect a tumour, or may detect it without
being able to eliminate it. Surviving clones can be selected for immune-evasion
traits: they may lose or alter tumour antigens, reduce antigen presentation,
resist immune-induced cell death, or express immune ``braking'' molecules
(checkpoint ligands such as PD-L1) that dampen T-cell activity.

Tumours can also recruit suppressive immune cells and release factors that
exclude or exhaust effective T cells. Abnormal tumour vessels, dense
surrounding scar-like tissue (\emph{stroma}), physical barriers, and a hostile
metabolic environment can further limit immune access and function. These
processes are part of \emph{cancer immunoediting} (the ongoing evolutionary
tug-of-war between the immune system and developing tumours): immune pressure
can remove sensitive cells while allowing less visible or more resistant clones
to persist; different cells within one tumour can also display different
antigens. Checkpoint inhibitors and other immunotherapies
can enhance antitumour activity in some settings, but they do not work for
every tumour and can also produce immune-related inflammation or organ injury.

\section{Why does lack of sleep impair the brain so quickly?}\index{brain}\index{sleep}\index{circadian rhythm}\index{learning}\index{memory}

Wakefulness creates homeostatic sleep pressure. Adenosine is one chemical
associated with this pressure and generally rises in the brain during waking,
while the circadian clock continues to promote or oppose alertness according
to the time of day. As sleep pressure rises, attention and arousal become
unstable. Because attention and decision-making depend on tightly coordinated,
high-energy neural networks, their performance can deteriorate before there is
obvious structural injury. Brain networks supporting working memory,
decision-making, emotional control, and learning no longer coordinate as
reliably, producing slower
reactions, more errors, and brief microsleeps that a person may not notice.

Sleep also provides state-dependent opportunities for memory consolidation,
metabolic regulation, and cellular maintenance. Staying awake cannot fully
substitute for those processes. Effort, stimulation, or caffeine may temporarily
increase alertness; caffeine partly blocks adenosine receptors and can improve
some tasks, but it does not reliably restore the judgment, learning, and
vigilance of adequate sleep and may disrupt later sleep. Repeated short sleep
can cause deficits to accumulate. The degree of impairment varies with the
person, task, time of day, and prior sleep, but prolonged deprivation can
severely impair functioning. Microsleeps make driving, working at heights, and
operating machinery unsafe even when a person feels able to continue.

\section{How does the brain store a memory physically?}\index{brain}\index{memory}\index{genes}\index{hippocampus}\index{learning}\index{neurons}

Learning changes the strength and sometimes the structure of synapses—a key
mechanism known as long-term potentiation (LTP), involving NMDA and AMPA
glutamate receptors—by altering receptor numbers, tiny receiving branches
(\emph{dendritic spines}), neuronal excitability, and patterns of gene and
protein activity. Short-lived memories
can depend on temporary changes in electrical activity and synaptic chemistry;
longer-lasting memories generally require gene expression and protein synthesis
that stabilize cellular and synaptic changes. A memory is therefore best
understood as a distributed network pattern, often called an engram, in which
particular groups of neurons and their connections are changed. An engram is a
useful scientific model of such a change, not a complete recording stored in
one neuron or one anatomical location.

For many episodic and spatial memories, the hippocampus (a seahorse-shaped
structure deep in the temporal lobe) helps encode and initially bind information
from different sensory systems; bilateral hippocampal damage famously impairs
the ability to form new long-term episodic memories (\emph{anterograde amnesia})
while sparing older memories and motor skills. With time and repeated
reactivation, aspects of some memories may become less dependent on the
hippocampus and more reliant on distributed cortical networks (the brain's outer
layers), although the exact process varies with memory type and remains an
active area of research.
Procedural skills and emotional memories also rely on partly different circuits,
and one experience can have sensory, spatial, emotional, and action-related
components distributed across them. Retrieval reactivates and reconstructs the
network pattern; this can make the memory temporarily malleable before it is
stabilized again, so remembering can
both preserve information and subtly modify it.

\section{Why do we forget?}\index{memory}

Forgetting can begin before a durable memory is formed: if attention is
divided, information may never be encoded strongly. A memory that was encoded
can later become harder to retrieve because its context or cues have changed,
because similar memories interfere with one another, or because the underlying
neural representation has weakened or been remodelled. Recalling a memory also
reconstructs it, so later experiences can update or distort what is retained.

Failure to recall does not necessarily mean that the memory has been erased;
an effective cue or a change in state can sometimes restore access. Recovery
with a cue does not guarantee that the resulting recollection is complete or
accurate, because cues can also evoke a related memory or a source error. Some
forgetting may be an active biological process, in which weakening selected
representations reduces interference and leaves the system more flexible, but
its importance differs by memory type and remains under study. Normal everyday
forgetting is not the same as progressive, functionally significant memory loss
caused by conditions such as dementia. The relative contributions of decay,
interference, retrieval failure, and active forgetting remain active research
questions.

\section{How can the brain recognise a familiar face in a fraction of a second?}\index{brain}\index{memory}

Visual processing is massively parallel. Face detection can begin within tens
of milliseconds under favourable conditions, while judging familiarity and
identifying an individual generally unfolds over the following tens to hundreds
of milliseconds (with face-selective brainwaves such as the N170 peaking around
170 milliseconds). Early visual areas extract edges, contrast, and motion;
face-selective regions analyse the arrangement of facial features and their
identity. A distributed network across the lower back and side of the visual
brain (the occipital face area, fusiform face-selective cortex, and superior
temporal sulcus) coordinates this analysis; anterior temporal and memory
systems can then associate the face with a person's name, history, and
emotional significance. Damage to these occipitotemporal regions can cause
\emph{prosopagnosia} (``face blindness''), where vision is intact yet familiar
faces cannot be recognised. The fusiform region is important, but it is not a
single face-recognition module operating alone.

Repeated exposure tunes these networks to the relatively stable relationships
among a familiar person's features across changes in lighting, viewpoint, and
expression. Familiarity can therefore emerge rapidly through pattern matching,
before deliberate reasoning is complete. It is an inference influenced by visual
evidence, context, expectation, and memory, rather than an infallible readout.
Poor viewing conditions, disguises, unusual viewpoints, and similar-looking
people can produce mistaken recognition. Recognizing a face as
familiar is distinct from identifying the person and from retrieving their name
or other facts; these stages can take longer or fail independently.

\section{Why can damage to a tiny brain region change speech, personality, or memory?}\index{brain}\index{memory}\index{speech}\index{language acquisition}\index{learning}\index{neurons}

Brain functions are implemented by networks, not by isolated pieces of tissue.
A small region may contain neurons that perform a specialized computation, act
as a bottleneck or hub, or carry a crucial long-range wiring bundle (a
\emph{white-matter tract}) between larger areas—for instance, a tiny stroke in
the internal capsule can cause widespread weakness on one side of the body.
A lesion can therefore interrupt communication across a much
larger network. Swelling, altered blood flow, or loss of normal input can also
temporarily disturb structurally intact regions; remote functional suppression
caused by loss of normal input is often called diaschisis. It is not the same as
direct destruction of the distant region.

The effect depends more on the lesion's location, depth, timing, and disruption
of connectivity than on its visible size. Pre-existing network reserve and
individual anatomy also influence the result. Damage in a language network (such
as Broca's area, Wernicke's area, or the arcuate fasciculus connecting them)
can cause impaired language production or comprehension (\emph{aphasia}); injury
to medial temporal or thalamic circuits can impair memory, while orbitofrontal
or prefrontal injury can markedly alter judgement, behaviour, or personality. These effects do not mean that speech,
personality, or memory resides in one tiny ``centre''; they reflect the failure
of connected circuits. Recovery varies because the brain can reduce swelling,
restore function in partially injured tissue, and recruit other pathways, but
reorganization may be incomplete and outcomes differ among individuals.

\section{What does neuroscience know about consciousness and neural activity?}\index{consciousness}\index{brainstem}\index{brain}\index{cerebral cortex}\index{memory}\index{sleep}

Neuroscience strongly links ordinary consciousness to organised activity in the
brain. It distinguishes the level of consciousness, such as wakefulness, from
its contents, such as a sight, thought, memory, or feeling. Brainstem arousal systems (including the ascending reticular activating system,
or ARAS) and thalamic relays help support wakefulness, while interacting
networks across the cerebral cortex and thalamus---especially two-way
thalamocortical loops and long-range signalling linking sensory areas with
frontal, parietal, and posterior cortical regions---help construct particular
experiences. Clinically, doctors assess disorders of consciousness (such as
coma, unresponsive wakefulness/vegetative state, and minimally conscious state,
while carefully distinguishing \emph{locked-in syndrome}, where awareness is
intact despite near-total paralysis) using bedside tools such as the Glasgow
Coma Scale.

Sensory input, bodily signals, memory, attention, and communication among brain
regions all shape what a person experiences. Sleep, anaesthesia, brain injury,
and electrical or magnetic stimulation can change consciousness along with
brain activity, providing converging evidence for this dependence. Leading
frameworks---such as global workspace theories, which emphasise widespread
cortical broadcasting of information, and integrated information theories, which
focus on how strongly brain circuits combine and differentiate signals---make
testable proposals, yet none is complete or universally accepted. EEG, imaging,
and related methods measure neural signals indirectly; they do not directly
measure the private experience itself. Scientists still do not know exactly how
physical activity produces subjective experience, so finding a neural correlate
is not the same as fully explaining consciousness.

\chapter{Brain and Nervous System}

\section{Why cannot we consciously control most internal organs?}\index{nervous system!autonomic}\index{bladder}\index{blood pressure}\index{brain}\index{breathing}\index{carbon dioxide}

The heart, digestive tract, blood vessels, pupils, and many glands contain
cardiac or smooth muscle and secretory cells that are regulated mainly by the
autonomic nervous system (balancing sympathetic ``fight-or-flight'' and
parasympathetic ``rest-and-digest'' branches, such as the vagus nerve,
coordinated by the hypothalamus and brainstem medulla), the enteric nervous
system, local circuits, and hormones rather than by the voluntary motor system. Their control loops
continuously sense variables such
as blood pressure, carbon dioxide, temperature, and gut contents and adjust
activity to maintain homeostasis. This automatic regulation can continue during
sleep and frees attention for interaction with the outside world. Local
feedback is also faster and more reliable for moment-to-moment regulation than
conscious commands would be.

Involuntary does not mean completely independent of the brain or impossible to
feel. People can consciously sense signals such as heartbeat, bladder fullness,
gut discomfort, or breathing effort without being able to command the underlying
organs directly. Emotions, posture, exercise, breathing, learned habits, and some
forms of biofeedback can alter autonomic activity, but most people cannot
directly command their heart or intestine to contract in a chosen pattern.
Breathing is a clear exception: its rhythm is normally automatic, but its
skeletal muscles can be voluntarily controlled for a time. Rising carbon
dioxide and other chemical signals eventually create a powerful involuntary
drive to breathe, although trained individuals can extend voluntary control.

\section{Why does the brain consume about 20 percent of resting energy?}\index{brain}\index{glucose}\index{glycogen}\index{neurons}\index{neurotransmitters}\index{oxygen}

In a typical adult at rest, the brain is about 2\% of body mass yet receives
roughly 15\% of cardiac output (about 750--1000 mL of blood per minute) and uses
roughly 20\% of the body's resting metabolic energy and oxygen consumption
(consuming roughly 120 grams of glucose per day). This is an
approximate average, not 20\% of every person's total daily energy expenditure;
the value varies with age, metabolic state, and how it is measured. A major cost
is maintaining electrical ion gradients across neuronal membranes during ongoing
electrical and synaptic activity. Neurons and supportive cells (\emph{glial
cells}) also recycle neurotransmitters, support synapses, and perform cellular
housekeeping.

Because the brain contains only small local fuel reserves, including stored
carbohydrate in supportive astrocyte cells (\emph{astrocyte glycogen}), it
depends on a continuous blood supply of oxygen and glucose—making it vulnerable
within minutes to interrupted blood flow (ischaemia) or severe low blood sugar
(hypoglycaemia); during prolonged fasting it can also use fat-derived
\emph{ketone bodies}.

\section{How much does a human brain weigh?}\index{brain}\index{learning}

An average adult human brain weighs about 1.2--1.4 kilograms, or roughly
3 pounds. A full-term newborn's brain weighs about 350--400 grams and grows
rapidly during the first two years of life, driven mainly by myelination,
growth of dendrites and synapses, and expansion of glial cells rather than the
birth of large numbers of new neurons. Brain weight varies naturally with age,
body size, sex, and individual development, and gradual, mild tissue loss can
occur with normal ageing, though marked shrinkage reflects specific disease.
These averages are therefore reference estimates, not a target or diagnostic
limit.

Brain weight alone is not a reliable measure of intelligence, personality, or
ability. Modest associations between brain measures and some cognitive outcomes
can appear across groups, but they do not predict an individual's abilities
well. What matters for function includes the organisation and connectivity of
neural circuits, along with development, learning, health, and many environmental
factors.

\section{Why do we have two cerebral hemispheres?}\index{brain}\index{language acquisition}\index{learning}

The cerebral hemispheres develop as two connected halves of the forebrain,
consistent with the body's broader bilateral organization. Having two partly
independent systems permits parallel processing and some specialization while
retaining communication between them. During development and learning, some
functions become more lateralized; for example, language is left-dominant in
roughly 95\% of right-handed people and about 70\% of left-handed people, while
some aspects of visual-spatial attention are more often right-dominant.
Dominance is relative rather than absolute, varies among
individuals, and can be modified by learning, development, or injury. These are
population tendencies, not rigid rules, and most complex abilities use networks
in both hemispheres.

The main bridge of nerve fibres between the two sides (\emph{corpus callosum}),
together with smaller cross-bridges (\emph{commissures}), exchanges information
between the hemispheres and helps coordinate perception, movement, learning,
and emotion—which is why a stroke on one side or surgical disconnection of the
corpus callosum can produce distinct, side-specific deficits. The hemispheres are therefore neither completely separate nor
simply duplicate copies. Parallel processing and specialization are plausible
functional advantages, but the precise evolutionary history of this arrangement
is not explained by one settled theory. The popular ``left-brain'' and
``right-brain'' personality categories are not scientifically valid.

\section{Why do many motor and sensory nerve pathways cross sides?}\index{brainstem}\index{brain}\index{smell (olfaction)}\index{spinal cord}\index{vision!visual field}

The crossing is a feature of how particular nerve pathways develop, not a
single cable deliberately routed from one side to the other. About 75--90\% of
the main movement fibres from the cortex (\emph{corticospinal} fibres) cross in
the lower brainstem at the medullary pyramids (the \emph{pyramidal decussation})
to form the lateral corticospinal tract, so the left motor cortex contributes
mainly to skilled movement on the right side and vice versa. Sensory pathways
also cross, though at different levels: fine touch and joint-position fibres
(the dorsal column--medial lemniscus pathway) cross in the medulla, whereas pain
and temperature fibres (the spinothalamic tract) cross soon after entering the
spinal cord—a pattern clinicians use to pinpoint whether an injury lies in the
spinal cord, brainstem, or brain. Visual pathways are arranged so that each
hemisphere primarily represents the opposite visual field, not the opposite eye.

Axon-guidance signals establish these crossings during development. Their
functional and evolutionary explanation is not settled; proposed accounts
include coordination of bilateral body and visual maps and advantages of
specialized, efficiently organized circuits. The rule is not universal:
smell (\emph{olfactory}) pathways remain largely on the same side
(\emph{ipsilateral}), some motor and sensory fibres remain uncrossed, and the
two sides communicate extensively. Thus ``the left
brain controls the right body'' is a useful summary of a major pattern, not an
absolute rule for every nerve, sensation, or movement. Crossings are partial and
modality-specific, so both hemispheres generally receive some information from
both sides of the body.

\chapter{Heart, Kidneys, Digestion, and Metabolism}

\section{Why does the heart almost never develop primary cancer?}\index{cancer}\index{heart}\index{muscles}

Primary cardiac tumours are exceptionally rare, and malignant primary tumours
are rarer still (about 75\% of primary heart tumours are benign, most often
myxomas). Most adult heart-muscle cells (\emph{cardiomyocytes}) have very
limited renewal and divide very little, which reduces one route by which
replication errors and uncontrolled cell multiplication (\emph{clonal
expansion}) can initiate cancer; it does not make them biologically incapable
of acquiring abnormal changes. This cannot be the whole explanation, because
the heart also contains fibroblasts, endothelial cells, smooth-muscle cells,
immune cells, and other cell types that can proliferate. Proposed additional
factors include tissue-specific growth controls, the surrounding structural
scaffolding (\emph{extracellular-matrix} properties), immune surveillance, and
the mechanical and metabolic environment of constantly contracting tissue, but
their relative contributions remain uncertain.

Primary tumours do occur, including benign myxomas (gelatinous growths often
arising in the left atrium) and rare malignant sarcomas or lymphomas. Tumours
that spread from other organs (such as lung, breast, or melanoma) are far more
common than primary cardiac tumours and may involve the heart muscle
(\emph{myocardium}), inner lining, or outer sac (\emph{pericardium}). When a
cardiac mass does occur, it can cause shortness of breath, chest pain, heart
rhythm disturbances, fluid around the heart, or blood-clot/tumour emboli, so
new unexplained cardiac symptoms merit evaluation. Thus ``rare primary heart cancer'' does not mean that the heart
is completely resistant to cancer; it means that cancers arising from cardiac
tissue are uncommon and that the heart may be affected, when it is affected by
cancer, as a site of secondary spread. The biological reasons for the rarity
of primary cardiac cancer are therefore plausible but not fully established.

\section{Why can heart muscle continue contracting briefly outside the body?}\index{heart}\index{muscles}\index{nervous system!autonomic}\index{brain}\index{calcium}\index{glucose}

The heart has an intrinsic, or self-generated (\emph{myogenic}), rhythm.
Pacemaker cells in the sinoatrial node slowly build electrical charge
(\emph{depolarize})—typically 60--100 times per minute at rest, with slower
backup rates in the AV junction (about 40--60/min) and Purkinje fibres (about
20--40/min)—through the coordinated activity of membrane ion channels;
intracellular calcium cycling also contributes to their self-firing
(\emph{automaticity}). Their impulses spread through the conduction system and
cell-to-cell electrical pores (\emph{gap junctions}). Calcium entry and calcium
release in working cardiomyocytes then activate the contractile proteins.
Initiation of a beat therefore does not require the brain or external nerves,
although autonomic nerves and hormones normally adjust the rate and force.
Pacemaker tissue is automatic; most working heart muscle is excitable and
contracts when it receives a propagated electrical impulse rather than
initiating the normal rhythm itself.

After removal—as during heart transplantation, where cold preservation
solutions typically limit safe cold ischaemia to roughly 4--6 hours—an isolated
heart may beat briefly using residual oxygen and stored energy. Sustained
beating, however, requires an appropriate perfusion system that maintains
oxygen, glucose, electrolytes, temperature, pH, and waste removal. Without
circulation, oxygen and fuel are rapidly depleted, ionic gradients fail, and
oxygen-starvation injury (\emph{ischaemic injury}) and dangerous rhythm
disturbances (\emph{arrhythmias}) develop. Rhythmic
contraction alone is not equivalent to a functioning circulation: the valves,
chambers, pressure generation, and perfusion must also remain coordinated.

\section{How does the heart know when to speed up during exercise?}\index{heart}\index{adrenal glands}\index{adrenaline}\index{brain}\index{breathing}\index{cardiac output}

The heart does not consciously ``know'' that exercise has begun. The response
starts before arterial blood oxygen has fallen substantially. As movement is
initiated, ``central command'' from motor-related brain regions acts in parallel
on cardiovascular control centres. It withdraws parasympathetic braking via the vagus nerve (\emph{vagal braking})
of the sinoatrial node and increases sympathetic drive. This feed-forward signal
is supplemented by feedback from mechanically and metabolically sensitive
nerve endings in working muscles---the exercise pressor reflex. These sensory
fibres (group III and IV afferents) respond to forces and to local changes
associated with contraction, including potassium, acidity, and other
metabolites; they are not a simple direct measurement of blood oxygen.

Sympathetic nerves and circulating adrenaline increase pacemaker rate and,
through adrenaline-sensitive beta-adrenergic signalling, usually increase
contractile force, helping cardiac output rise from roughly 5 litres per minute
at rest to 20--30 litres per minute or more during vigorous exercise. Venous
return also rises: the contracting muscles compress veins, breathing produces
a thoracic pressure pump, and the resulting increase in filling stretches heart
chambers to pump more strongly (the \emph{Frank--Starling mechanism}). At higher
intensities, stroke volume may plateau, so further increases in cardiac output
depend more on heart rate. Blood-pressure sensors in major arteries and the
heart (\emph{baroreflexes}) continue to regulate the response; during exercise
the arterial baroreflex is reset around a higher pressure. Local chemical
signals dilate vessels in active muscle, while vasoconstriction or different
control operates in other vascular beds.

Chemoreceptor input becomes more important when exercise is severe or gas
exchange is impaired. The final response varies with intensity, posture,
temperature, fitness, medications (for example, beta-blockers blunt the heart
rate rise), age, and cardiovascular disease (such as chronotropic incompetence,
where the heart cannot speed up appropriately with exertion).

\section{Why does blood not normally clot inside healthy vessels?}\index{blood!clotting}\index{atherosclerosis}\index{collagen}\index{inflammation}\index{platelets}\index{proteins}

The endothelial lining normally presents a smooth, anticoagulant surface. Its protective sugar-protein coating (\emph{glycocalyx}) and intact cell layer
limit platelet contact with the collagen and tissue factor beneath the vessel
wall. Endothelial cells release platelet-calming substances such as nitric oxide
and prostacyclin, and express built-in blood-thinning systems involving heparan
sulfate--antithrombin, thrombomodulin--protein C and protein S, and
tissue-factor-pathway inhibitor. Surface enzymes (\emph{ectonucleotidases})
break down platelet-activating ADP. Endothelial tissue-type plasminogen
activator also supports clot breakdown (\emph{fibrinolysis}), converting
plasminogen to plasmin so fibrin strands can be removed after they form.

Normal blood flow dilutes and carries away activated factors, and shear helps
limit the persistence of some platelet-adhesive von Willebrand factor. Flow
alone, however, does not prevent thrombosis: stasis, disturbed flow, and
abnormal shear can instead encourage it. When a vessel is injured or becomes
dysfunctional, collagen and von Willebrand factor recruit platelets, while
exposed tissue factor initiates coagulation; local procoagulant signals then
overcome the endothelial restraints. Endothelial dysfunction, atherosclerosis,
inflammation, sluggish flow (such as prolonged immobility during long travel or
hospitalisation, which can trigger a deep-vein thrombosis), or increased
coagulability can cause thrombosis without an obvious external wound. These three contributors---abnormal vessel wall,
abnormal flow, and increased coagulability---are often summarized as
Virchow's triad.

\section{How does the body stop bleeding without clotting the entire circulation?}\index{blood!clotting}\index{collagen}\index{platelets}\index{proteins}\index{wound healing}

Hemostasis is a rapid, localized sequence. The injured vessel constricts to
reduce blood loss. Injury exposes underlying collagen (\emph{subendothelial collagen}) and the
bridging protein von Willebrand factor, allowing platelets to adhere, activate,
change shape, and aggregate into a temporary plug. Exposed tissue factor and
the surface of activated platelets then assemble coagulation reactions that
generate the clotting enzyme \emph{thrombin} locally. Thrombin converts soluble
\emph{fibrinogen} to insoluble \emph{fibrin} strands, activates additional
platelets and clotting factors, and creates a mesh that stabilizes the plug.
Platelet contraction can also draw the wound edges closer together.

Several features keep this reaction from spreading through the circulation.
The intact endothelium around the injury remains antiplatelet and
anticoagulant. Antithrombin, the thrombomodulin--protein C and protein S
system, tissue-factor-pathway inhibitor, dilution and clearance of activated
factors, and blood flow all limit propagation. Endothelial tissue-type
plasminogen activator supports the regulated generation of plasmin, which
degrades fibrin as repair proceeds; plasminogen-activator inhibitors help
prevent premature, widespread fibrinolysis. The clot is then remodelled and
removed as the vessel regains integrity. If these localizing and limiting
systems fail---or if disease (such as haemophilia), medication (such as
antiplatelets or anticoagulants), or an inherited defect alters them---
excessive bleeding or unwanted thrombosis can result.

\section{Why are arteries thick-walled while veins contain valves?}\index{arteries}\index{veins}\index{capillaries}\index{heart}\index{muscles}\index{oxygen}

Arteries generally receive higher, more pulsatile pressure than veins (mean
systemic arterial pressure averages roughly 90--100~mmHg, compared with much
lower pressures of only a few mmHg in large central veins). Their
relatively thick walls, especially their smooth-muscle and elastic layers,
resist distension, regulate diameter, and allow elastic recoil to smooth the
pulse between heartbeats. Veins carry blood at much lower pressure, so their
walls are generally thinner and more compliant; they can hold a large fraction
of the circulating blood and act as capacitance vessels. Veins often have a
relatively substantial outer connective-tissue layer even though their
muscular middle layer is thinner.

Many veins of the limbs contain one-way valves. By dividing a long hydrostatic
column into shorter segments, they reduce the tendency of gravity to drive
blood backward; when these valves weaken over time, blood can pool and stretch
superficial leg veins into varicose veins or cause chronic venous insufficiency. Contraction of surrounding skeletal muscles then compresses
the veins and, with the valves, pumps blood toward the heart; the respiratory
pump and sympathetic changes in venous smooth-muscle tone also contribute.
Valves are not present in every vein, because some veins rely mainly on short
distances, pressure gradients, or surrounding structures.

Finally, ``artery'' and ``vein'' describe flow direction relative to the heart,
not oxygen content: pulmonary arteries carry oxygen-poor blood at lower pressure
than systemic arteries, while pulmonary veins carry oxygen-rich blood.
Capillaries have a different, thin-walled design suited to exchange rather than
pressure handling.

\section{Why does blood not simply pool in the legs when we stand?}\index{nervous system!autonomic}\index{brain}\index{breathing}\index{dehydration}\index{heart}\index{muscles}

Some blood does shift into the compliant veins of the legs and abdomen when we
stand---often several hundred millilitres after moving from lying to standing.
The body limits the resulting fall in venous return, but does not eliminate the
pooling. Valves in many leg veins prevent retrograde flow, but do not remove
the hydrostatic pressure or make a motionless person immune to pooling.
Contractions of the calf, thigh, and, during some movements, abdominal muscles
compress veins and propel blood toward the heart. This skeletal-muscle pump is
especially important during walking; quiet standing provides much less pumping,
although postural sway can help. Rhythmic breathing changes pressure in the
chest relative to the abdomen and also assists venous return through the
respiratory pump.

At the same time, reduced stretch of arterial baroreceptors triggers increased
sympathetic activity. This raises heart rate, increases contractility, and
constricts arterioles and venous capacitance vessels---especially in the
splanchnic circulation---helping maintain arterial pressure and brain
perfusion. Sympathetic constriction of veins within leg muscles is limited, so
prolonged immobility, dehydration, autonomic disease, medications, or weak leg
muscles can overwhelm these defenses and cause dependent swelling, orthostatic
symptoms, or increased venous-stasis risk. The effectiveness of these
mechanisms varies with posture, movement, hydration, vascular compliance, and
health.

\section{Why does blood pressure change when we suddenly stand up?}\index{blood pressure}\index{nervous system!autonomic}\index{brain}\index{breathing}\index{dehydration}\index{heart}

Standing suddenly creates a hydrostatic pressure difference. Gravity shifts
several hundred millilitres of blood into distensible veins in the legs and
abdomen and increases filtration of some plasma into the tissues. Less blood
temporarily returns to the heart, so cardiac filling, stroke volume, cardiac
output, and often arterial pressure fall. Pressure measured in the dependent
feet, however, rises because of the same hydrostatic column; pressure in the
brain can briefly fall if compensation is inadequate.

The reduced stretch of baroreceptors in the carotid sinuses and aortic arch is
detected within seconds. Their reflex response withdraws vagal braking and
increases sympathetic activity, raising heart rate and contractility and
constricting arterioles and capacitance vessels, particularly in the
splanchnic circulation. Leg-muscle activity and breathing further assist
venous return. Over a longer interval, the kidneys and hormones such as renin,
angiotensin II, and aldosterone help retain salt and water, but this is not the
immediate correction.

A brief change can be normal; orthostatic hypotension is a measured fall of at
least 20~mmHg in systolic pressure or 10~mmHg in diastolic pressure within
3 minutes of standing, whether or not symptoms are present. Dehydration,
medications, blood loss, autonomic disease, or other causes can impair
compensation and produce dizziness, visual dimming, or fainting.

\section{How can the kidneys filter 180 litres daily but produce only 1--2 litres of urine?}\index{kidneys}\index{urine}\index{amino acids}\index{glucose}\index{hormones}\index{proteins}

The 180 litres is an approximate volume of glomerular filtrate, not 180 litres
of whole blood discarded from the body. A filtration rate of about 120 mL/min
would yield about 173 litres per day, which is conventionally rounded to
roughly 180 litres; the value varies with body size, age, and kidney function.
The glomerular barrier filters water and many small, unbound solutes from
plasma. Blood cells and most large proteins normally remain in the circulation,
although very small amounts of protein may cross and are usually reclaimed or
degraded by the tubules.

As the filtrate travels through the renal tubules, transporters and osmotic
gradients return about 99\% of the filtered water and useful solutes such as
glucose, amino acids, sodium, and bicarbonate to the blood. Other substances,
including hydrogen ions, potassium in regulated segments, drugs, and drug
metabolites, may be actively secreted into the tubule. Urea and other solutes
can be reabsorbed to varying degrees. The small, selectively adjusted
remainder is excreted as urine. Daily urine output is therefore variable:
1--2 litres is a common resting range, not a fixed requirement, and it can be
much higher or lower with hydration, hormones, illness, or medication. For
example, glucose is normally almost completely reclaimed, but can appear in
urine when its filtered load exceeds tubular transport capacity.

\section{How do kidneys know how much water and salt to retain?}\index{kidneys}\index{brain}\index{hormones}\index{urine}

The kidneys do not ``know'' consciously; they respond to several sensors and
hormones. Hypothalamic concentration sensors (\emph{osmoreceptors}) detect blood solute
concentration (\emph{osmolality}, normally kept tightly around 275--295
mOsm/kg) and help control thirst and antidiuretic hormone (ADH, also called
vasopressin). ADH inserts water pores (\emph{aquaporin-2}) into collecting-duct
cells, allowing more water to return to the blood (as after sweating or eating a
salty meal) and producing a smaller amount of more concentrated urine.
Volume depletion can stimulate ADH and thirst even when osmolality is not high,
because preserving circulation becomes the priority.

Sensors in the kidney and cardiovascular system detect renal perfusion, blood
pressure, and circulating volume. Reduced pressure, sympathetic stimulation,
or reduced sodium chloride delivery to a specialised kidney sensor
(\emph{macula densa}) can activate renin release. The
renin--angiotensin--aldosterone system promotes sodium reabsorption; aldosterone
also increases potassium and, in some segments, hydrogen-ion secretion. Water
often follows retained sodium, but water balance is regulated separately by ADH
and depends on collecting-duct permeability. When the body has extra volume,
atrial and brain natriuretic peptides suppress renin and aldosterone and promote
salt excretion (\emph{natriuresis}), while other signals increase water
excretion. Local pressure-and-salt feedback in each nephron
(\emph{tubuloglomerular feedback}) also adjusts filtration. Clinically,
disruptions of these pathways—such as excessive ADH (SIADH), ADH deficiency or
resistance (diabetes insipidus), or blood-pressure drugs that block the
renin--angiotensin--aldosterone system or diuretics—directly alter salt and
water balance. The
kidneys therefore balance water, sodium, potassium, and acids continuously
rather than using one single water-or-salt switch.

\section{Why is urine normally very low in microbes while the intestine is full of them?}\index{urine}\index{bladder}

Compared with the intestine, the urinary tract is a low-nutrient, continuously
flushed environment. The upper tract normally has a very low microbial burden,
and urine flows through the ureters to a bladder that is emptied regularly.
This removes organisms entering through the urethra. Variable pH and
osmolality, urea and other urinary factors, the glycan-coated urothelium,
local immune defences, and forward flow all make persistent colonization more
difficult. The lower urinary tract may contain a low-abundance resident
community, so a sensitive test does not imply that every healthy urinary tract
is completely sterile. Results also depend on collection and testing methods.

The intestine, in contrast, receives food and secretions that provide abundant
nutrients, retains contents, and offers a large mucus-lined surface on which a
dense community can grow. Its epithelial and immune defences contain most
organisms within the lumen without trying to sterilize it.

Even with these defences, bacteria can still ascend the urinary tract---especially
when flow is obstructed, urine is retained, a catheter is present, or host
defences are impaired---causing a urinary tract infection. Bacteria in a
specimen may represent contamination or asymptomatic bacteriuria rather than
infection; symptoms, inflammatory findings, culture quality, and clinical
context determine the significance.

\section{Why can dehydration become dangerous faster than starvation?}\index{dehydration}\index{brain}\index{diarrhoea}\index{glucose}\index{glycogen}\index{kidneys}

Water is the main solvent and transport medium of the body's extracellular and
intracellular fluids. A substantial deficit reduces plasma volume, strains the
circulation, disturbs electrolyte concentrations, impairs brain function, and
can reduce kidney perfusion and waste excretion. It also reduces sweating and
other heat-loss mechanisms, so heat illness can develop quickly during
exertion or hot weather. Water and electrolytes are continually lost in urine,
breath, sweat, and stool, and the body has only a limited reserve. The exact
electrolyte disturbance depends on what is lost and what is replaced: sweating
often causes both water and sodium loss, whereas drinking excessive plain water
can dilute plasma sodium and cause dangerous hyponatraemia. Severe dehydration can become life-threatening rapidly—sometimes within days,
and faster in extreme heat, infection, vomiting, diarrhoea, infancy, or older
age—and warning signs such as minimal or absent urination, confusion, sunken
eyes, rapid heartbeat, or low blood pressure warrant urgent care.

By contrast, after food intake stops, hormonal changes allow the body to use
glycogen briefly (typically depleted within roughly 24 hours) and then stored
fat, while lowering energy expenditure and producing ketone bodies that supply
much of the brain's energy after metabolic adaptation (allowing survival for
weeks in many healthy adults, depending on body reserves and hydration). Some glucose production and protein breakdown still continue, and
eventually fat, protein, vitamin, and mineral reserves become insufficient.
These energy stores generally last longer than the body's immediately usable
water reserve, so a person may survive longer without food than without water.
There is no universal survival timetable: heat, activity, illness, age, body
reserves, ongoing fluid loss, and electrolyte or vitamin deficiencies can change
the outcome. This comparison is not a safe fasting or dehydration plan.

\section{Why does stomach acid not normally digest the stomach itself?}\index{stomach}\index{cell membrane}\index{digestion}\index{inflammation}\index{oxygen}\index{proteins}

The stomach is exposed to both hydrochloric acid and the protein-digesting
enzyme pepsin, so it relies on a dynamic protective barrier rather than being
immune to digestion. Surface cells secrete a continuous, adherent mucus gel
containing bicarbonate. The gel slows acid diffusion and limits pepsin contact,
while bicarbonate creates a near-neutral microenvironment (pH near 7) at the
cell surface, protecting it from the highly acidic stomach interior (pH around
1--2), where pepsin is active. Tight junctions limit back-diffusion of acid,
and the cell membrane and its surface phospholipids add further resistance.
Good mucosal blood flow carries away diffused acid, supplies oxygen, and helps
maintain the local chemical gradients; prostaglandins support this flow and
stimulate mucus and bicarbonate secretion.

The surface epithelium is renewed rapidly, and nearby cells can spread over
small defects (restitution) before new cells divide, so minor injury can be
repaired. These defenses are protective, not perfect. \textit{Helicobacter pylori}
can promote inflammation and impair barrier function, while nonsteroidal
anti-inflammatory drugs reduce protective prostaglandin production; smoking,
severe physiological stress, and other factors can add risk. Acid and pepsin
can then contribute to inflammation, erosions, or ulcers, usually together
with impaired defence and healing rather than as an isolated effect of acid
alone; vomiting blood, black tarry stools, or sudden severe abdominal pain
requires urgent medical evaluation for bleeding or perforation.

\section{Why is the small intestine several metres long?}\index{small intestine}\index{bile}\index{digestion}\index{lymphatic system}\index{proteins}\index{vitamins}

The adult small intestine is commonly around 5--6 metres long, although its
length varies with the person and with how it is measured. Its length provides
substantial area and time for digestion and absorption, while segmentation,
peristalsis, and other muscular contractions mix the contents and move them
gradually through the duodenum, jejunum, and ileum. These regions are not
interchangeable: the duodenum receives bile and pancreatic secretions, the
jejunum absorbs much of the available nutrient load, and the ileum recovers
important remaining nutrients, bile salts, and vitamin B12 (so disease or
removal of the terminal ileum can lead to B12 deficiency and bile-acid
diarrhoea).

Length is only part of the solution. Circular folds, finger-like villi, and
microscopic microvilli amplify the absorptive surface area roughly several
hundred-fold to tens of square metres. Digestive enzymes,
transport proteins, blood vessels, and lymphatic lacteals then help move
absorbed substances into the body; many lipids enter lymph before the blood.
The intestine is coiled to fit this large functional surface inside the
abdomen. Its exact length does not by itself predict absorption, and severe
loss, obstruction, or mucosal damage can cause malabsorption, including
short-bowel syndrome after major resection.

\section{Why do intestinal bacteria influence metabolism, immunity, and possibly brain function?}\index{brain}\index{immune system}\index{metabolism}\index{antibiotics}\index{bile}\index{genes}

Intestinal bacteria digest dietary fibre and resistant starch that human
enzymes cannot fully break down, producing short-chain fatty acids such as
acetate, propionate, and butyrate. Butyrate is an important fuel for many colon
cells; these metabolites can also influence gut hormones, liver metabolism,
immune signalling, appetite, and energy handling. Bacteria transform bile
acids and other dietary or host molecules, and some synthesize vitamins, but
the nutritional importance of these products varies and is not a substitute
for a balanced diet. These effects depend on species, genes, diet, location,
dose, and host physiology.

Microbial molecules and cell-surface patterns interact with the intestinal
barrier and immune cells, helping shape immune tolerance while supporting
defence against pathogens. The gut and brain communicate through the vagus and
enteric nervous systems, hormones, microbial metabolites, and immune signals;
some metabolites or their derivatives enter the circulation and influence
distant tissues. Inflammation, barrier changes, and altered tryptophan or bile
acid metabolism are among proposed routes, but they do not establish a simple
one-microbe-to-one-behaviour pathway.

Many human findings about microbiome effects on brain function and disease are
correlational, and results may be influenced by diet, medication, geography,
and illness. Evidence from animal studies does not automatically translate to
humans. A microbiome association does not by itself prove causation or show
that a probiotic, antibiotic, or dietary change will provide a treatment.

\section{Why can the intestine absorb nutrients without allowing most microbes into the bloodstream?}\index{amino acids}\index{inflammation}\index{liver}\index{lymphatic system}\index{vitamins}

The intestine is a selectively permeable barrier rather than an impermeable
wall. Absorptive enterocytes use specific transporters, channels, and some
controlled vesicular processes to move digested sugars, amino acids, ions,
vitamins, and water across the epithelium; some lipids are packaged as
fat-transport particles (\emph{chylomicrons}) that enter intestinal lymph vessels
(\emph{lacteals}) before reaching the bloodstream. Tight junctions regulate the
spaces between cells (\emph{paracellular route}), while the intact epithelial
layer normally prevents larger microorganisms from crossing in significant
numbers. Small microbial products can cross in controlled amounts
and are processed by local immune cells or, after portal drainage, by the liver.

Several additional defences reinforce this selectivity. Mucus and the
epithelial glycocalyx limit microbial contact, specialised Paneth-cell
antimicrobial peptides inhibit organisms near the crypts, and secretory IgA
binds and helps contain microbes in the lumen. Immune cells in the underlying
connective tissue layer (\emph{lamina propria}) survey what crosses the barrier.
Specialised sampling cells (\emph{M cells}) can deliver sampled luminal
material to immune tissues, but this is controlled immune sampling, not
unrestricted entry of microbes into the blood. Inflammation, infection,
ischemia, malnutrition, or physical injury can loosen tight junctions or damage
epithelial cells, allowing microbial products and occasionally whole organisms
to translocate and potentially cause systemic inflammation or infection.

\section{Why does the body maintain blood glucose within a narrow range?}\index{glucose}\index{adipose tissue}\index{adrenal glands}\index{adrenaline}\index{brain}\index{consciousness}

The body keeps glucose within a relatively narrow, changing range (fasting
blood glucose in healthy adults typically lies around 70--99~mg/dL, or
3.9--5.5~mmol/L, with values below 70~mg/dL generally considered hypoglycaemia)
because different tissues depend on a reliable fuel supply, while excessive
glucose alters proteins, osmotic balance, and blood vessels. Red blood cells lack
mitochondria and depend almost entirely on glucose; the brain normally uses
substantial glucose too, although it can use ketone bodies during prolonged
fasting. Severe hypoglycaemia can quickly impair brain function and cause
seizures or loss of consciousness. Persistent hyperglycaemia promotes tissue
injury over time, while severe acute elevation can cause osmotic dehydration,
hyperosmolar illness, or, in appropriate metabolic settings, ketoacidosis.

After a meal, pancreatic beta cells release insulin, assisted by gut-derived
hormone signals (\emph{incretins}). Insulin promotes glucose uptake through
cell-surface GLUT4 glucose transporters in skeletal muscle and adipose tissue,
encourages glycogen and fat storage, and suppresses hepatic glucose production.
When glucose falls, pancreatic alpha cells release glucagon; the liver breaks
down glycogen and makes new glucose from non-carbohydrate precursors
(\emph{gluconeogenesis}). Glucagon has little direct effect on skeletal muscle glycogen.
Adrenaline, cortisol, and growth hormone provide additional
counter-regulatory support in stress or prolonged fasting. This feedback
system is dynamic rather than perfectly fixed, and the liver acts as a buffer
between meals.

\section{Why can chronically high glucose damage eyes, kidneys, nerves, and vessels?}\index{glucose}\index{kidneys}\index{arteries}\index{atherosclerosis}\index{blood pressure}\index{capillaries}

Persistently high glucose drives several interacting forms of injury. It can
attach non-enzymatically to proteins, creating advanced glycation end products
that stiffen or cross-link tissues and promote inflammation. Excess glucose
also increases oxidative stress and can disrupt the polyol, protein-kinase-C,
and hexosamine pathways. Over time, these changes can damage endothelial
cells, thicken basement membranes, alter vessel permeability, and narrow or
weaken small blood vessels, reducing tissue blood flow and repair.

In the retina, damaged capillaries can leak, cause macular swelling, and
stimulate abnormal new-vessel growth; traction or bleeding can threaten
vision. In the kidney, glomerular pressure and barrier injury can produce
albumin leakage and progressive nephron loss. Nerves can be injured directly
by metabolic stress and indirectly when the small vessels supplying them
deliver too little oxygen and nutrients—commonly producing a symmetric
``glove-and-stocking'' numbness or tingling in the feet and hands. Larger-vessel
atherosclerosis also contributes to coronary, cerebral, and peripheral artery
disease. Risk depends on the duration and severity of hyperglycaemia and on
blood pressure, lipids, smoking, genetics, and other factors; regular clinical
monitoring—such as HbA1c blood tests, urine albumin-to-creatinine checks, and
annual eye and foot examinations—helps detect injury early. Good diabetes
management reduces risk but does not make complications impossible. Acute, very high glucose can
additionally cause dehydration and metabolic emergencies, which are distinct
from these slowly developing complications.

\section{How do Type 1 and Type 2 diabetes differ?}\index{diabetes}\index{adipose tissue}\index{autoimmune disease}\index{breathing}\index{genes}\index{glucose}

Both Type 1 and Type 2 diabetes cause chronically elevated blood glucose (commonly
diagnosed when fasting plasma glucose is $\ge$ 126~mg/dL or HbA1c is $\ge$
6.5\% on confirmed testing), but their underlying mechanisms differ. Type 1 diabetes usually results from
autoimmune destruction of insulin-producing pancreatic beta cells, leading to
little or no insulin production and requiring lifelong insulin replacement. It
often begins in childhood or young adulthood, though it can arise at any age.
Type 2 diabetes combines \emph{insulin resistance}---reduced responsiveness to
insulin in skeletal muscle, liver, and adipose tissue---with a progressive
inability of beta cells to produce enough insulin to compensate. It is more
common in adults, influenced by genetics, age, body-fat distribution, and
lifestyle, and is managed with lifestyle measures, non-insulin medicines, and
sometimes insulin.

Other forms exist, including gestational diabetes during pregnancy and rare
single-gene forms, and clinical features can overlap. Without enough insulin,
people with Type 1 diabetes can rapidly develop \emph{diabetic ketoacidosis},
a life-threatening buildup of acids and ketones, whereas severe insulin
resistance and dehydration in Type 2 diabetes can trigger a hyperosmolar
hyperglycaemic state (extremely high blood glucose without marked ketoacidosis).
Extreme thirst, frequent urination, unexplained weight loss, vomiting,
abdominal pain, fruity breath, rapid breathing, or confusion requires urgent
medical care.

\section{Why does exercise improve insulin sensitivity?}\index{insulin}\index{adipose tissue}\index{ATP (adenosine triphosphate)}\index{calcium}\index{capillaries}\index{glucose}

Contracting skeletal muscle needs more ATP and can move GLUT4 glucose
transporters to its cell surface through contraction-sensitive pathways
involving calcium and AMP-related energy sensing. This uptake can occur partly
independently of insulin, so active muscle can remove glucose from blood even
when insulin signalling is impaired, although the response is not necessarily
normal in every metabolic disease. After exercise, insulin sensitivity often
remains increased while muscle replenishes its glycogen stores; this effect
commonly lasts for hours and sometimes into the following day, depending on
the activity and the person.

Repeated aerobic and resistance exercise can increase muscle mass, capillary
density, mitochondrial capacity, GLUT4 availability, and the amount of tissue
able to store and use glucose. Training can also improve glycogen-synthesis
responses and reduce visceral fat and related inflammatory signalling, improving
whole-body insulin action in muscle, liver, and adipose tissue. The size and
duration of the benefit vary with exercise type, training status, nutrition,
illness, and medications; the effect diminishes when training stops. People
using insulin or insulin-releasing medicines may develop low blood glucose
during or after exercise, while very intense exercise can transiently raise it;
individualized medical guidance may be needed.

\section{Why does muscle use much more glucose during exercise?}\index{glucose}\index{muscles}\index{ATP (adenosine triphosphate)}\index{glucagon}\index{glycogen}\index{insulin}

Exercise sharply raises ATP demand in active muscle. Contraction-sensitive
signals move GLUT4 glucose transporters to the muscle-cell surface, partly
through pathways that do not require insulin, while local blood flow increases.
Muscle also breaks down its own glycogen through glycogenolysis, a rapid
carbohydrate reserve whose glucose is used locally rather than released into
the blood (because skeletal muscle lacks the enzyme glucose-6-phosphatase).
Sympathetic signals and changes in insulin and glucagon help
coordinate this supply with the wider circulation.

The fuel mixture is not fixed. High-intensity exercise generally relies more on
carbohydrate from blood glucose and muscle glycogen, because carbohydrate can
provide ATP at a high rate. Lower-intensity or prolonged exercise permits a
larger contribution from fatty-acid oxidation, although carbohydrate remains
important. Duration, training, nutrition, workload, and glycogen availability
all alter the balance. When glycolysis outpaces mitochondrial oxidation, some
pyruvate is converted to lactate; lactate can be exported and recycled into
glucose by the liver (the Cori cycle) or used by other tissues, rather than
being merely waste.

The liver helps maintain circulating glucose through glycogen breakdown and
gluconeogenesis, including recycling some lactate, but its output may not fully
match muscle uptake. Thus the important change is increased fuel use by working
muscle; it is not a universal rule that every exercise session uses much more
glucose specifically.

\section{Why does the body store energy as glycogen and fat?}\index{glycogen}\index{adipose tissue}\index{ATP (adenosine triphosphate)}\index{glucose}\index{liver}\index{metabolism}

Glycogen is a rapidly mobilised, branched glucose polymer stored mainly in the
liver and skeletal muscle. Liver glycogen helps stabilise blood glucose between
meals, whereas muscle glycogen is a local fuel reserve that can supply ATP
quickly, including when energy demand is too high for aerobic metabolism alone.
Storing glucose as glycogen also avoids the extreme osmotic effect of storing
the same number of glucose molecules separately. Glycogen stores are limited (roughly 100 grams in the liver and 400 grams in
skeletal muscle in a typical adult), and glycogen binds substantial water (about
3 grams of water per gram of glycogen, which explains why rapid early weight
changes on low-carbohydrate diets mostly reflect water shifts), so this form of
storage is relatively bulky.

Most long-term energy is stored as triacylglycerol in adipose tissue. Fat yields
about 9~kcal per gram, compared with about 4~kcal per gram for carbohydrate,
and does not bind comparable amounts of water, allowing a much larger reserve
in less mass. It is mobilised and oxidised more slowly and requires adequate
oxygen, so it cannot support maximal rapid ATP production as well as
carbohydrate. Fatty acids generally cannot be converted into net glucose in
humans, although glycerol and certain unusual fatty-acid substrates can
contribute to gluconeogenesis; therefore fat cannot directly replace glucose
for all glucose-dependent tissues. Keeping both stores combines rapid,
readily accessible carbohydrate with dense, durable energy storage. Adipose
tissue also provides insulation and secretes signalling molecules.

\section{Why cannot humans synthesise certain vitamins and amino acids?}\index{amino acids}\index{vitamins}\index{metabolism}\index{proteins}\index{skin}

Human cells can make some nutrients from simpler molecules, but they lack the
complete biochemical pathways needed to make others, or cannot make them fast
enough to meet physiological needs. This is not a general inability to make
protein: adults can synthesise many amino acids, but nine amino acids---
histidine, isoleucine, leucine, lysine, methionine, phenylalanine, threonine,
tryptophan, and valine---are essential and must be supplied by food. Other
amino acids can become conditionally essential during growth, severe illness,
prematurity, or when a precursor or metabolic pathway is limited; arginine,
cysteine, tyrosine, glutamine, and glycine may be conditionally important in
such settings.

Most vitamins are dietary requirements because human metabolism cannot
construct their molecular structures in sufficient amounts. There are useful
exceptions and partial sources: ultraviolet B light enables skin to make
vitamin D, the body can convert some tryptophan to niacin, and intestinal
bacteria produce some forms of vitamin K. These sources depend on exposure,
precursors, gut function, and absorption and may be insufficient; they do not
eliminate the need for a balanced diet. Vitamins are organic micronutrients,
not sources of energy, and their requirements vary with age, growth, pregnancy,
illness, genetics, sunlight exposure, and the nutrient concerned. ``Essential''
means required from outside the body under the relevant conditions, not that
the nutrient is more important than a nonessential one.

\section{Why do humans need oxygen continuously?}\index{oxygen}\index{ATP (adenosine triphosphate)}\index{brain}\index{glucose}\index{heart}\index{mitochondria}

Most cells with mitochondria use oxygen as the final electron acceptor in the
cellular energy-harvesting assembly (\emph{electron-transport chain}). This
permits the chain to maintain an electrochemical proton gradient across the inner
mitochondrial membrane, which drives ATP synthesis (\emph{oxidative
phosphorylation}) and produces roughly 30--32 molecules of ATP per glucose—far
more than the 2 ATP yielded by oxygen-free (\emph{anaerobic}) glycolysis; when
oxygen falls, pyruvate is converted to lactate, whose concentration rises when
production outpaces clearance.

Oxygen is therefore needed continuously at the level of the whole organism,
although not every human cell uses it directly: mature red blood cells lack
mitochondria and rely on glycolysis, and some tissues can tolerate brief
periods of low oxygen. Tissues differ in their tolerance. Skeletal muscle can
increase anaerobic ATP production briefly, whereas the brain and heart are
especially vulnerable because of high continuous energy demands and limited
fuel reserves. Continuous oxygen use requires ventilation, adequate blood
oxygen carriage, circulation, and delivery to cells; a failure of any link—such
as severe anaemia, circulatory shock, or poisoning of mitochondrial respiration
by carbon monoxide or cyanide—can cause tissue injury.

\section{Why can brain cells suffer irreversible injury after minutes without oxygen?}\index{brain}\index{oxygen}\index{ATP (adenosine triphosphate)}\index{breathing}\index{calcium}\index{glucose}

The brain has a high metabolic rate, depends mainly on aerobic use of glucose,
and stores very little fuel. Severe hypoxia, or more commonly ischaemia that
removes oxygen and glucose while allowing wastes to accumulate, can therefore
lower ATP levels within minutes. When ATP-dependent ion pumps fail, neurons
lose their electrical gradients, depolarise, swell as sodium and water enter,
and release excessive glutamate. Excess glutamate can drive calcium into
neurons, contributing to acidosis, mitochondrial injury, reactive oxygen
species, and activation of enzymes that damage cell structures. This cascade,
called excitotoxic injury, can progress to cell death.

There is no single time limit at which every brain cell becomes irreversibly
injured, though loss of consciousness typically occurs within seconds of
complete circulatory arrest and irreversible neuronal injury often begins within
roughly 4--6 minutes at normal body temperature (cooling, as in therapeutic
hypothermia, slows metabolism and can extend this window). Vulnerability depends
on the severity and duration of oxygen and blood flow loss, glucose
availability, temperature, pre-existing disease, and the speed and quality of
resuscitation. Electrical activity may stop before every
cell is dead, and partially perfused tissue may remain recoverable. Rapid
restoration of breathing and circulation gives salvageable tissue the best
chance of recovery, although inflammation, mitochondrial injury, and
reperfusion-related oxidative stress can cause damage to continue after blood
flow returns.

\section{Why does oxygen not damage us despite being chemically reactive?}\index{oxygen}\index{DNA}\index{inflammation}\index{mitochondria}\index{proteins}\index{stress}

Cells use oxygen in tightly controlled chemical reactions rather than allowing
it to react freely with cell components. Even so, mitochondrial respiration and
other reactions produce reactive oxygen species (ROS), such as superoxide and
hydrogen peroxide. Superoxide dismutases convert superoxide to hydrogen
peroxide, while catalase, glutathione peroxidases, and thioredoxin systems help
detoxify peroxides. Repair and disposal systems also remove oxidised proteins,
lipids, and DNA. These defences reduce damage; they do not make oxygen
chemically harmless or prevent every oxidant reaction.

ROS are not simply harmful waste products. In controlled amounts they help
with cell signalling, adaptation to stress, and the destruction of microbes by
immune cells. They become dangerous when production exceeds antioxidant,
repair, and removal capacity; this imbalance is called oxidative stress and
can damage membranes, proteins, and genetic material. Transition metals can
convert peroxides into especially reactive radicals. Hyperoxia, inflammation,
smoking, toxins, and ischaemia followed by reperfusion can upset this balance
and cause tissue injury. Thus oxygen is essential within a regulated range,
but a higher concentration is not automatically safer.

\section{How can the lungs provide enormous gas-exchange surface inside the chest?}\index{lungs}\index{alveoli}\index{capillaries}\index{carbon dioxide}\index{oxygen}

The airways branch repeatedly into smaller passages that end in hundreds of
millions of tiny air sacs called alveoli. The branching and folding create a
very large gas-exchange surface---roughly 50--100 $\mathrm{m}^{2}$ in commonly
cited adult estimates---inside a compact chest. The conducting airways mainly
move air; exchange occurs in the respiratory zone. Type-I alveolar cells and
the capillary endothelium form a remarkably thin alveolar--capillary barrier
(typically only about 0.2--0.5 micrometres thick), so oxygen and carbon dioxide
cross only a short distance. A dense capillary
network brings blood alongside much of this surface.

The alveoli must remain open while retaining a moist lining for gas exchange.
Type-II alveolar cells release surfactant, which lowers surface tension and
reduces the tendency of small alveoli to collapse during exhalation (a lack of
surfactant causes neonatal respiratory distress syndrome in premature infants);
these cells can also help restore injured epithelium. The combination of huge area,
a short diffusion distance, continual ventilation, and appropriately matched
blood flow allows rapid gas exchange in a small volume. Regional differences
in ventilation and perfusion still exist, and emphysema, fibrosis, fluid, or
airway disease can reduce effective surface area or increase diffusion
distance.

\section{Why does fluid not normally accumulate inside the alveoli?}\index{alveoli}\index{capillaries}\index{heart}\index{inflammation}\index{lungs}\index{lymphatic system}

Alveoli are not completely dry: a very thin lining fluid is needed for
surfactant and gas exchange. Normally, the alveolar--capillary barrier and the balance between outward
blood-vessel pressure (\emph{hydrostatic} force) and inward protein-driven
water pull (\emph{oncotic} force) limit movement of fluid from capillaries into
the lung. Small amounts that enter the surrounding tissue spaces
(\emph{interstitium}) are removed by pulmonary lymphatics. Type-II alveolar
cells and other epithelial cells actively move sodium out of the air space
through epithelial sodium channels (ENaC) and a cell-base sodium--potassium
pump; water follows the resulting osmotic gradient. This continuous balance
keeps the air spaces open enough for gas exchange without letting them become
flooded.

Pulmonary oedema develops when fluid entry exceeds clearance---for example,
when left-sided heart failure raises pulmonary capillary pressure (pulmonary
capillary wedge pressure rising well above its normal ~8--12~mmHg), when low
blood protein reduces the inward oncotic pull, or when infection,
inflammation, inhaled toxins, or severe systemic injury makes the barrier
abnormally leaky. Excess fluid can dilute surfactant, lengthen the diffusion
path, and fill alveoli, reducing gas exchange and blood oxygen—producing urgent
symptoms such as severe breathlessness, difficulty breathing while lying flat
(\emph{orthopnea}), or pink frothy sputum. The mechanism
therefore may be predominantly hydrostatic or permeability-related, and the
distinction matters clinically.

\chapter{Skin, Muscles, Development, and Hormones}

\section{Why do we have five fingers and five toes?}\index{embryonic development}

Humans inherited the basic five-digit limb pattern from the five-digited
(\emph{pentadactyl}) limbs that became widespread among early four-limbed
vertebrates (\emph{tetrapods}). The earliest tetrapods did not
all have exactly five digits, so five is an inherited evolutionary pattern, not
the inevitable starting number. During embryonic development, interacting
signals in the limb bud---including the limb-tip growth ridge (the apical
ectodermal ridge), the limb-patterning zone of polarising activity (ZPA), and
morphogen signals such as sonic hedgehog and BMP-driven programmed cell death
between digits—organise limb growth, digit identity, and digit separation.
Changes in these processes can produce fewer or additional digits. Evolution
has repeatedly modified the pattern: horses retain one main functional toe with
reduced side digits, birds have specialised wing digits, and some animals have
more than five.

Five digits are therefore a historically inherited and developmentally stable
pattern, not a number evolution deliberately designed as a universal optimum.
In humans, the arrangement supports grasping, manipulation, walking, and load
distribution, but those functional advantages help explain why it has been
retained rather than why it first evolved. Extra digits (\emph{polydactyly}) or
fused digits (\emph{syndactyly}) can occur as isolated harmless variations or
occasionally as part of an inherited syndrome, showing that digit number and
form are not biologically fixed.

\section{Why are fingerprints unique?}\index{fingerprints}\index{ageing}\index{genes}\index{scars}\index{skin}\index{identical twins}

Fingerprints are patterns of friction ridges that form before birth—mainly
between roughly weeks 10 and 24 of gestation—as fetal fingertip pads
(\emph{volar pads}) regress and epidermal ridges interlock with underlying
dermal papillae. Genes help establish broad features such as ridge direction
and pattern type, but the fine details are also shaped by local growth,
pressure, movement, tissue mechanics, and the prenatal environment. The combination of these influences
makes ridge endings, splits, and pores highly distinctive---even in identical
twins, whose prints are more alike than those of unrelated people but still
distinguishable. The exact developmental process is complex, and there is no
single gene that specifies an entire fingerprint.

The pattern is supported by deeper skin structures, so ordinary growth and
surface renewal preserve its arrangement throughout life. A superficial scrape
usually heals with the same ridge pattern, whereas a deep cut, burn, or skin
disease can replace part of it with scar tissue. Dry skin, ageing, pressure,
or a poor impression can merely make a print harder to read. Thus ``unique''
means highly distinctive in practical comparison, not that absolute uniqueness
has been mathematically proved or that every partial-print forensic match is
free of error; fingerprint evidence has quality and interpretation limits.

\section{Why does skin continuously shed without becoming thinner?}\index{skin}\index{ageing}\index{collagen}\index{skin!epidermis}\index{inflammation}\index{keratin}

Cells in the deepest, basal layer of the epidermis continually divide. Some
remain as stem or progenitor cells to maintain the renewing population, while
others move toward the surface, differentiate, accumulate tough keratin, lose
their nuclei and organelles, and become corneocytes. Connections between these
cells are gradually loosened so that they are shed in tiny, usually unnoticed
flakes. New cells from below replace them, so the outer epidermis is often
renewed over roughly a month, with timing varying by body region, age, and
other conditions.

Skin thickness stays approximately stable because cell production,
keratinization, and upward replacement are balanced by surface loss. The
balance is not perfect: friction can produce a thicker callus, while burns,
inflammation, ageing, and diseases such as psoriasis can alter production or
desquamation and cause visible scaling. The deeper dermis, which contains
collagen, elastin, vessels, nerves, and glands, also contributes to total skin
thickness and changes independently of the epidermis.

\section{Why does hair grow continuously in some locations but only to limited lengths elsewhere?}\index{hair}\index{hormones}

Individual hair follicles do not grow continuously. Each cycles through an
anagen (growth) phase, a brief catagen (regression) phase, a telogen (resting)
phase, and an exogen shedding phase. The exposed hair shaft is mostly dead,
keratinized material; growth occurs in living cells at the follicle bulb.
Because neighbouring follicles are at different stages, the scalp appears to
grow hair continuously even though each individual hair eventually stops and
falls out.

The length a hair can reach depends mainly on how long its follicle remains in
anagen, together with its growth rate and how much the shaft breaks. Scalp
follicles—about 85--90\% of which are normally in anagen at any given time,
growing roughly 0.3--0.4~mm per day—may grow for several years, whereas eyebrow
and eyelash follicles usually grow for only a few months. Body-hair cycles vary
by location, and age, hormones, genetics, illness, nutrition, medicines, and
follicle disorders (such as temporary shedding in \emph{telogen effluvium} or
progressive follicle shrinking in \emph{androgenetic alopecia}) can change them. The apparent length limit is therefore set by the follicle's
cycle and shaft durability, not by a fixed limit in the exposed hair itself.

\section{Why do humans have relatively little body hair compared with other primates?}\index{hair}\index{hormones}\index{puberty}\index{skin}\index{sweating}\index{body temperature}

Humans are not truly hairless. We have roughly as many follicles as other
primates (commonly estimated at 2--5 million across the body), but much of our
body hair is fine, short vellus hair rather than thick terminal hair. This is a difference in follicle type and hair characteristics,
not simply the absence of follicles. Less insulating hair allows air to reach
the skin more easily and makes evaporation of sweat a more effective way to
lose heat, especially during sustained activity in hot, dry conditions.

Thermoregulation, together with the evolution of extensive sweating and reduced
body insulation, is a leading explanation for this evolutionary change, but it
is not the only proposed influence. Parasite pressure, sexual selection,
developmental changes, and other ecological factors may also have contributed;
the history cannot be reduced to one proven cause. Hair remains useful for
local protection, sensation, limited insulation, and scalp protection from
sunlight, while pubertal androgens convert vellus follicles into thicker
terminal hair in regions such as the armpits, pubic area, and often the face
and chest. In cold
conditions the reduced hair covering provides less insulation, and in humid
conditions it provides less cooling benefit because sweat evaporates less
effectively.

\section{Why do humans sweat so much?}\index{sweating}\index{dehydration}\index{hypothalamus}\index{skin}

Humans rely heavily on sweating for evaporative cooling. We have roughly 2--4 million eccrine
sweat glands distributed over much of the body, capable of producing over a
litre of sweat per hour during intense heat and exertion, and warm-sensitive
pathways in the hypothalamus activate them through sympathetic, mainly
cholinergic nerves during heat exposure or exercise. Emotional or adrenergic signals can also
activate sweating in the palms, soles, and other areas. When sweat changes from
liquid to vapour, it carries heat away from the skin and the blood flowing
beneath it. This helps humans remain active in heat, particularly in dry,
moving air.

Eccrine glands initially secrete a fluid containing water and electrolytes; as
it passes through the duct, sodium and chloride are reabsorbed to varying
degrees. Rapid sweating leaves less time for reabsorption and can therefore be
saltier. Sweating is effective only when enough of it evaporates. Humidity,
still air, heavy clothing, or sweat that drips away can greatly reduce cooling.
The response also causes water and salt loss, so heavy sweating can lead to
dehydration, electrolyte disturbances, and heat illness—ranging from heat
exhaustion (heavy sweating, weakness, nausea) to life-threatening heatstroke
(markedly elevated core temperature with confusion or collapse), while replacing
heavy sweat losses with plain water alone can trigger dangerous low blood
sodium (exercise-associated hyponatraemia). Acclimatization can
increase sweat production while improving salt conservation. Apocrine glands
contribute mainly to body odour after bacterial breakdown; they are not the
main glands responsible for temperature control.

\section{Why does shivering warm the body?}\index{shivering}\index{adipose tissue}\index{ATP (adenosine triphosphate)}\index{brainstem}\index{hypothalamus}\index{muscles}

Shivering is a rapid, involuntary pattern of alternating skeletal-muscle
contractions. The contractions use ATP, but little of their energy becomes
useful movement; much is released as heat. Cold-sensitive pathways in the skin
and body core signal temperature-control regions of the hypothalamus, which
coordinate brainstem and spinal motor circuits so many muscles can generate
heat at the same time. Shivering can raise resting metabolic heat production several-fold, but
it consumes oxygen and fuel and is less efficient than simply reducing heat
loss.

Cold also constricts blood vessels in the skin. This does not produce heat, but
it reduces warm blood flow to the surface and slows heat loss, so it often
works alongside shivering. Other responses include increased metabolic heat
production, including non-shivering thermogenesis driven by uncoupling protein 1
(UCP1) in brown adipose tissue, which is especially important in infants.
Shivering cannot continue indefinitely; below the hypothermia threshold of
$35\,^{\circ}\mathrm{C}$—and especially in severe hypothermia—it may weaken or
stop as muscles and the nervous system lose function, a sign of worsening
rather than recovery.

\section{How does the human body maintain its temperature?}\index{nervous system!autonomic}\index{fever}\index{hormones}\index{hypothalamus}\index{metabolism}\index{shivering}

The preoptic and other hypothalamic regions act as a temperature-control
network. They receive signals from the blood and temperature-sensitive nerves
in the skin and deeper tissues, compare them with a regulated range, and
coordinate autonomic nerves, hormones, and behaviour to balance heat production
with heat loss. This is a negative-feedback system: a change away from the
usual range triggers responses that tend to push body temperature back toward
that range, though the correction is not perfect.

When the body is too warm, skin blood vessels widen and eccrine sweating
promotes evaporative cooling; a person may also reduce activity or seek shade.
When it is too cold, skin vessels narrow, shivering and increased metabolism
produce more heat, and behaviour adds clothing or shelter. Brown-fat
thermogenesis is particularly relevant in infants and can contribute in adults.
Heat is continuously made by metabolism and lost through radiation, convection,
conduction, and evaporation. Humidity, wind, clothing, and contact with water
alter these routes, so temperature is regulated dynamically rather than held at
one exact value.

About $37\,^{\circ}\mathrm{C}$ is an average reference, not a universal target:
normal values vary with measurement site, time of day, activity, age, hormones,
and the individual. In fever, immune mediators cause the hypothalamus to raise
the regulated target, producing chills until that target is reached. In
heatstroke, body temperature rises despite a largely unchanged target because
heat gain exceeds the body's ability to lose heat; this is hyperthermia, not
simply a fever.

\section{Why do muscles become stronger after being stressed?}\index{muscles}\index{bones}\index{fatigue}\index{learning}\index{proteins}\index{sleep}

When a muscle is repeatedly exposed to a challenging but recoverable load, it
adapts to handle that demand more effectively. Mechanical and metabolic
signals activate cellular growth pathways (such as mTORC1 signalling) and
resident muscle stem cells (\emph{satellite cells}) that increase the
production and maintenance of contractile proteins, which can enlarge muscle
fibres and improve their force-producing capacity. Changes in fibre arrangement, leverage, and the
ability to tolerate fatigue can also affect performance. Tendons, connective
tissue, and bone may adapt, although they do so at different rates.

Strength is not just muscle size. Early gains often come from the nervous
system learning to recruit and coordinate motor units more effectively; later
gains can include greater muscle size and movement-specific skill. Small
amounts of muscle damage may accompany unfamiliar or strenuous exercise, but
damage and soreness are neither the goal nor a requirement for progress.
Consistent training with gradually appropriate loading, adequate energy and
protein, sleep, and recovery support adaptation. Excessive stress without adequate rest can instead
impair recovery or cause overtraining and tendon/muscle injury, and strength
can decline again when training is substantially reduced.

\section{Why does unused muscle rapidly shrink?}\index{muscles}\index{ageing}\index{inflammation}\index{insulin}\index{proteins}

Muscle is maintained partly in response to regular mechanical loading and
neural activity. When it is used much less—for example during bed rest,
casting, weightlessness, or paralysis—the signals that support muscle-protein
synthesis weaken, while protein breakdown through systems such as the
ubiquitin--proteasome and autophagy pathways can increase. The balance shifts
so that more protein is removed than replaced; muscle fibres shrink, producing
disuse atrophy—with measurable cross-sectional loss often appearing within 1 to
2 weeks of complete immobilisation or bed rest. Fibre type, blood flow, insulin sensitivity, and connective
tissue can also change.

Strength can fall faster than muscle size because the nervous system also loses
some recruitment and coordination, and the muscle's contractile quality can
change. Denervation is especially severe because the nerve-derived signal is
lost, whereas bed rest or casting leaves the nerve supply intact. The speed
and severity of loss vary with age, illness, nutrition, inflammation, and
duration of inactivity, and disuse atrophy differs from age-related muscle loss
(\emph{sarcopenia}) or disease-driven wasting (\emph{cachexia}). Much disuse atrophy can improve with gradual reloading
and rehabilitation, but recovery may be slow or incomplete after prolonged
denervation, ageing, or severe disease.

\section{Why do muscles become sore after unfamiliar exercise?}\index{muscles}\index{fatigue}\index{pain}\index{stress}\index{urine}

Delayed-onset muscle soreness (DOMS) usually follows unfamiliar, unusually
intense, or eccentric exercise, in which an active muscle lengthens under load.
The unfamiliar mechanical stress can disturb muscle and connective-tissue
structures and trigger local repair and immune responses. These changes,
together with increased sensitivity of local pain-sensing nerves, can make the
muscle tender and stiff. Soreness often begins 12--24 hours later, peaks about
24--72 hours after exercise, and then fades over several days; later sessions
of similar exercise often cause less soreness as the tissue adapts.

DOMS is different from the burning or fatigue felt during exercise and is not
caused by lactic acid remaining in the muscle. Its exact mechanism is complex
and not fully settled, and soreness is neither a reliable measure of workout
quality nor proof that muscle growth occurred. Severe focal pain, marked
swelling, profound weakness, inability to use a limb, or dark urine after
extreme exercise is not typical DOMS and needs medical assessment because
serious muscle injury or rhabdomyolysis is possible.

\section{Why are some muscles voluntary while others work automatically?}\index{muscles}\index{nervous system!autonomic}\index{bladder}\index{breathing}\index{homeostasis}\index{hormones}

The terms ``voluntary'' and ``automatic'' describe how a muscle is controlled,
not two completely separate kinds of muscle. Skeletal muscle (striated, voluntary-controlled fibres) is usually
activated by somatic motor neurons and can be recruited by conscious commands,
so it moves the limbs and helps maintain posture. Cardiac muscle (also striated,
with cells linked by intercalated discs and gap junctions) has pacemaker cells
in the sinoatrial (SA) node that initiate each heartbeat, while the autonomic
nervous system and hormones adjust its rate and force. Smooth muscle (non-striated,
often electrically coupled by gap junctions) in the gut, airways, blood vessels,
bladder, and uterus can respond to stretch, local chemicals, hormones, and
autonomic or enteric nerves; some smooth muscle also has intrinsic rhythmic
activity. Its contraction is not simply switched on by conscious choice.

The distinction is not absolute. Breathing, walking, posture, reflexes, and
shivering use skeletal muscle without conscious attention, although some of
these actions can also be influenced deliberately. Cardiac and smooth muscle
are not independent of the body either: local circuits, nerves, and hormones
continually modulate their intrinsic activity to maintain homeostasis. Thus
muscle type and control mode overlap, but they are not identical categories.

\section{How does the body know where a limb is with the eyes closed?}\index{ageing}\index{brain}\index{cerebellum}\index{fatigue}\index{joints}\index{muscles}

Proprioception is the sense of limb position and movement. Muscle spindles
embedded within skeletal muscles report muscle length and how fast it changes,
while Golgi tendon organs at muscle--tendon junctions report tension. Receptors
in joint capsules and skin add information about joint angle, pressure, and
skin stretch. Signals travel rapidly along large myelinated sensory fibres in
the dorsal columns of the spinal cord to the sensory cortex for conscious
awareness, and along spinocerebellar tracts to the cerebellum for unconscious
coordination, where they are combined with vestibular and visual cues.

This estimate allows posture and coordinated movement without vision, but it
is not infallible. Fatigue, numbness, ageing, alcohol, vitamin B12 deficiency,
or peripheral neuropathy can impair proprioception; vibrating a tendon can also
temporarily alter spindle signals and create the illusion that a limb is
moving. Loss of proprioceptive input produces sensory ataxia, in which balance
worsens markedly as soon as the eyes are closed (the basis of the clinical
Romberg test).

\section{Why cannot humans regrow an amputated limb while some animals can?}\index{bones}\index{genes}\index{inflammation}\index{muscles}\index{nails}\index{regeneration}

Some animals, especially salamanders such as axolotls, can regenerate a limb.
After amputation they form a blastema---a temporary, highly proliferative group
of cells with contributions from local tissues and, in some cases, altered or
dedifferentiated cells. Signals from nerves, immune cells, surrounding tissue,
and the remaining limb coordinate the growth, pattern, and positioning of new
skin, muscle, bone, nerves, and blood vessels.

Humans can repair many tissues, and very distal fingertip injuries---especially
in children and when the nail-matrix stem-cell region remains---may regenerate
some skin, nail, bone, and soft tissue. A major adult limb injury, however,
usually triggers pro-fibrotic macrophage and immune responses and a collagen
scar rather than a lasting blastema supported by sustained Wnt and FGF
growth-factor signalling.
Human cells, immune responses, nerves, extracellular matrix, and positional
signals do not provide the same coordinated regenerative environment, and
scarring can rapidly seal off the wound. The difference is therefore not one
missing gene or a complete lack of developmental information, but a complex
failure to organise all the required tissues safely and in the correct pattern.
Regenerative-medicine research is investigating these mechanisms while
limiting the risk of uncontrolled growth.

\section{How can a fetus develop an entire body from one fertilised cell?}\index{embryonic development}\index{genes}\index{placenta}

The fertilised egg, or zygote, first undergoes cleavage divisions, producing
many smaller cells without needing a miniature body plan inside each one. Early
cells can still give rise to both the embryo and supporting tissues. As a
blastocyst forms around day 5 after fertilisation, its inner cell mass supplies
the cells that will form the embryo, while the outer trophoblast contributes
mainly to extra-embryonic structures, including much of the fetal component of
the placenta. The
placenta also contains maternal tissue and mediates exchange and signalling
between mother and developing organism.

Cells then receive chemical and physical signals from their neighbours, their
position, and the maternal environment. These signals alter gene-regulatory
networks, causing cells to specialise. During gastrulation around week 3, they move into three
foundational layers---ectoderm, mesoderm, and endoderm. Further cell division,
movement, communication, growth, tissue folding, and controlled cell death
shape these layers into organs and body systems—with organ formation
(\emph{organogenesis}) peaking during weeks 3 to 8, when developing tissues are
especially sensitive to teratogens such as alcohol, certain medicines, and
infections. Development is therefore a
coordinated, self-organising process guided by genetic information and its
environment, not a simple assembly from a fixed microscopic blueprint.

\section{How does an embryonic cell know what type of cell to become?}\index{embryonic development}\index{epigenetics}\index{genes}\index{liver}\index{muscles}\index{skin}

An embryonic cell does not consciously choose its fate; over development,
cells narrow from \emph{totipotent} (able to form all embryonic and placental
tissues) to \emph{pluripotent} (able to form all embryonic cell types) to more
specialised \emph{multipotent} stem cells. Each cell receives chemical
and physical signals from nearby cells, its position, and its mechanical
environment. Diffusible signalling gradients (\emph{morphogens}, such as sonic
hedgehog patterning different nerve-cell types along the developing spinal
cord) and pathways such as Wnt, Hedgehog, Notch, and TGF-beta help cells
interpret their location. Receptors convert them into changes in
gene-regulatory networks, switching some genes on and others off. The cell's
earlier history, DNA-packaging (\emph{chromatin}) state, division pattern,
and ability to respond also affect how it interprets the same signal.

These changes first guide cells toward broad groups such as the three germ
layers, then toward specialised fates such as nerve, muscle, liver, or skin.
Feedback from neighbouring cells and stable gene-switching
(\emph{epigenetic}) changes help maintain each choice as cells divide, while
later signals can refine or sometimes alter it. Cell fate is guided rather than perfectly predetermined: some cells remain
flexible or can compensate after injury, whereas mistimed signalling or gene
regulation can cause developmental abnormalities.

\section{How do organs know when they have reached the correct size?}\index{cancer}\index{hormones}\index{regeneration}

There is no single ruler inside an organ, and the organ does not judge its size
consciously. Size emerges from feedback that balances cell growth, cell
division, differentiation, cell size, and programmed cell removal. Signals
from neighbouring cells, hormones, nutrients, blood supply, the extracellular
matrix, and mechanical stretch influence this balance. For example, when cells
become crowded or tissue tension changes, growth-braking pathways such as
Hippo--YAP/TAZ can slow further growth; after tissue loss—as when the liver
regrows after partial surgical removal or a remaining kidney enlarges to
compensate for a lost partner—repair and growth-factor signals can temporarily
promote it. Some organs also adjust cell size or function without
adding many cells.

Different organs use different combinations of these controls, so size is an
approximate, functional outcome that scales partly with the rest of the body,
not a single fixed target. The same feedback helps maintain proportions and
repair ordinary wear, but regeneration is organ-specific and often incomplete.
Genetic changes, chronic injury, altered mechanical load, endocrine disorders,
or cancer can disrupt the balance and cause undergrowth, cell enlargement
(\emph{hypertrophy}), increased cell number (\emph{hyperplasia}), scarring
(\emph{fibrosis}), or uncontrolled tumour growth.

\section{Why are identical twins not completely identical?}\index{identical twins}\index{chromosomes}\index{DNA}\index{embryonic development}\index{epigenetics}\index{fingerprints}

Monozygotic (``identical'') twins arise when one early embryo separates, so they
usually begin with nearly the same nuclear DNA. Their placental arrangement
(whether they share a single outer chorion or inner amniotic sac, depending on
when the embryo split) can differ, and monochorionic twins sharing a placenta
may receive unequal blood flow or nutrient exposure before birth—at the extreme,
causing twin-to-twin transfusion syndrome (TTTS). Mitochondrial DNA may not be exactly the same,
and mutations or chromosome changes can arise after the embryo forms and occur
in some tissues of one twin but not the other. Random X-chromosome
inactivation can also make gene expression differ between cells in female
twins.

Gene activity can differ even when the DNA sequence is nearly the same. These
epigenetic differences may arise during development and often increase with
age, illness, and different exposures. Random developmental events and the
twins' separate experiences can further affect appearance, fingerprints,
immune responses, microbiomes, behaviour, and disease risk. Such factors
influence complex traits without determining them on their own. Thus
``identical'' describes common origin and close genetic similarity, not perfect
biological sameness.

\section{Why does the immune system usually tolerate a fetus carrying paternal genes?}\index{genes}\index{immune system}\index{hormones}\index{placenta}\index{proteins}

The fetus carries some genetic markers from the father that are not present in
the mother, but maternal tolerance is an active, local process rather than a
complete immune shutdown. Maternal blood normally bathes the outer placental layer (the \emph{chorionic
surface}), while specialised placental cells (\emph{trophoblast} layers) and the
placental barrier separate it from fetal blood and most fetal tissues.
Trophoblast cells use an unusual pattern of immune-display molecules, lacking
classic HLA-A and HLA-B markers while expressing tolerogenic HLA-G and local
immune-calming pathways (such as IDO and PD-L1) that limit damaging immune
activation without making the interface invisible to immune surveillance.

Regulatory T cells and specialised uterine-lining (\emph{decidual}) immune
cells, including natural killer cells and macrophages, help maintain this
balance while supporting placental attachment and remodelling of uterine blood
vessels. Hormones, immune-messenger proteins (\emph{cytokines}), and physical
separation contribute as well. Pregnancy therefore changes immune regulation; it
does not simply weaken the entire immune system, and responses vary by
infection and by stage of pregnancy. Tolerance is incomplete and its disruption
may contribute to disorders such as miscarriage, pre-eclampsia, or Rh
alloimmunisation (when maternal antibodies target Rh-positive fetal red blood
cells), although these conditions have multiple causes.

\section{Why do women generally live longer than men?}\index{ageing}\index{autoimmune disease}\index{chromosomes}\index{menopause}\index{oestrogen}

Across most populations, women have a higher average life expectancy than men
(commonly by roughly 4 to 6 years globally, though the gap varies widely by
region and era), but no single cause explains the difference. Before menopause,
oestrogen may provide some cardiovascular protection, and sex-chromosome
differences (such as having two X chromosomes with cellular mosaicism that can
buffer harmful variants) and immune-system or oxidative-stress differences may
influence infection responses and ageing. These biological
effects are probabilistic and change with age; possible advantages are partly
offset by pregnancy-related risks and a greater prevalence of some autoimmune
diseases. They cannot by themselves account for the size of the difference in
every population.

Social and behavioural factors are also important. In many settings, men have
greater exposure to injuries, violence, occupational hazards, smoking, alcohol,
and other risks, and may have less access to or use preventive and medical care.
Women may face different risks, including unequal healthcare access, caregiving
burdens, and violence. The size and even the causes of the gap change with
country, income, culture, public health, discrimination, and historical
period. Life expectancy is an age-specific population average, not a guarantee
of lifespan; the distributions overlap substantially and individual variation
is large.

\section{Why does human fertility decline substantially with age?}\index{chromosomes}\index{DNA}\index{egg cells (oocytes)}\index{embryonic development}\index{menopause}\index{mitochondria}

In people with ovaries, the supply of follicles is finite and declines from
roughly 1--2 million at birth to several hundred thousand at puberty, about
25,000 in the late thirties, and near zero by menopause (averaging around age
51). With age, fewer follicles remain, ovulation becomes less reliable, and
oocyte quality falls. Age-related changes in the cell-division machinery
(\emph{meiotic spindle}), chromosome separation, mitochondria, and the
surrounding follicle increase the chance of embryos with an abnormal chromosome
count (\emph{aneuploid embryos}), miscarriage, and difficulty achieving a live
birth. Fertility often declines before menopause and falls more sharply after
about age 35; menopause ends ovulation, but the decline is gradual rather than
an instantaneous switch.

In people with testes, sperm production often continues into later life because
new sperm are made continuously. However, sperm concentration, movement,
shape, and DNA integrity may decline gradually, and the number of new mutations
in sperm can rise with paternal age. Illness, medications, hormonal changes,
and sexual-function problems can also affect fertility. The rate varies greatly
among individuals. Ovarian age is generally a stronger and earlier predictor of
fertility decline than paternal age, but both partners' health, reproductive
anatomy, genetics, timing, and age can matter; age alone does not determine an
individual's ability to conceive.

\section{What triggers puberty?}\index{puberty}\index{adrenal glands}\index{androgens}\index{brain}\index{hair}\index{hormones}

Puberty begins in the brain when hypothalamic neurons increase the pulsatile
release of gonadotrophin-releasing hormone (GnRH), a shift regulated by
neuropeptides such as kisspeptin alongside metabolic, genetic, and
environmental cues. GnRH stimulates the anterior pituitary gland to secrete
luteinising hormone (LH) and follicle-stimulating hormone (FSH), which act on
the ovaries or testes to drive sex-steroid production (primarily oestradiol and
testosterone) and gamete maturation.

A partly separate maturation of the adrenal glands, called \emph{adrenarche},
also increases adrenal androgen production and contributes to pubic and
axillary hair and body odour. Together, these hormones trigger growth spurts,
body-composition shifts, and secondary sexual characteristics over several
years. Normal puberty spans a wide age range influenced by genetics, nutrition,
and overall health. Unusually early signs of puberty (for example, before age 8
in girls or age 9 in boys) or markedly delayed puberty (no signs by age 13 in
girls or age 14 in boys) should be medically assessed.

\section{Why does the body have dozens of hormones when nerves transmit faster?}\index{hormones}\index{adrenal glands}\index{adrenaline}\index{nervous system!autonomic}\index{brain}\index{metabolism}

Nerve signals are rapid and relatively precisely directed, making them useful
for sensation, movement, and immediate adjustments. A synapse can target a
specific neighbouring cell, while autonomic nerves can regulate organs more
broadly. Hormones are released into the circulation or nearby tissue and act on
cells carrying the appropriate receptors, allowing coordination across distant
organs. Their effects often last longer and help regulate growth, metabolism,
reproduction, stress, and fluid balance. Some hormones, such as adrenaline,
act within seconds, whereas some neural effects can persist, so the categories
are not separated by an absolute speed boundary.

The two systems often work together. For example, the brain can rapidly detect
a challenge, activate sympathetic nerves, and trigger hormone release that
sustains the response. Neuroendocrine cells combine the two modes, and some
signals act locally on neighbouring cells (\emph{paracrine} messengers) or back
on the releasing cell itself (\emph{autocrine} messengers) rather than as
long-range hormones. The important differences are therefore the route,
distance, duration, and breadth of control each signal provides. Having both
systems lets the body respond quickly while also coordinating whole-body
changes over time.

\section{How can tiny quantities of hormones alter the whole body?}\index{hormones}\index{calcium}\index{genes}\index{kidneys}\index{liver}\index{proteins}

Hormones can work at very low concentrations because their target cells carry
specific receptors that bind them selectively. Once a hormone binds, it can
start a signalling cascade in which one activated molecule influences many
others. Peptide and many amine hormones usually act through cell-surface
receptors and internal relay molecules (second messengers such as cyclic AMP,
calcium, or protein-activating \emph{kinase cascades}). Steroid and thyroid
hormones act mainly through intracellular receptors that directly alter which
genes are read (\emph{gene transcription}), although some rapid surface-mediated
(\emph{non-genomic}) effects also occur.

A hormone does not affect every cell it reaches: the response depends on which
cells have the receptor, how many receptors they have, the hormone's free
concentration, and the cell's internal signalling machinery. Pulses, binding
proteins, receptor desensitisation, and feedback further shape the response.
Hormones are eventually inactivated or cleared, especially by the liver and
kidneys. This combination of selectivity, signal amplification, and feedback
allows a small circulating signal to produce a substantial but regulated
effect; abnormal levels or receptor responses can nevertheless cause widespread
disease.

\section{Why does stress affect the heart, intestine, immunity, sleep, and glucose together?}\index{glucose}\index{heart}\index{immune system}\index{sleep}\index{stress}\index{adrenal glands}

Stress is a coordinated brain--body response designed to shift resources toward
immediate survival rather than long-term maintenance. When the brain detects a
physical or psychological challenge, sympathetic nerves and adrenaline released
from the adrenal medulla quickly increase heart rate, blood pressure,
alertness, and the availability of glucose by promoting liver fuel release and
limiting non-essential energy expenditure.

A slower hypothalamus--pituitary--adrenal axis response—via hypothalamic CRH and
pituitary ACTH—releases cortisol from the adrenal cortex, which
helps sustain fuel supply, raises blood glucose, and modulates immune and
inflammatory activity so energy is not diverted unnecessarily during acute
threat. Stress networks also suppress deep sleep to keep the brain vigilant,
while signals from the gut and immune system feed back to the brain. A brief,
proportionate response is useful and resolves through negative feedback;
repeated or severe activation produces cumulative physiological wear and tear
(\emph{allostatic load}), turning those same adaptive shifts into persistent
disturbances of sleep, blood pressure, glucose regulation, digestion, and
inflammation. Effects vary with the type, duration,
predictability, controllability, and meaning of stress, as well as the
individual's health and support.

\section{How can emotions produce genuine physical symptoms?}\index{nervous system!autonomic}\index{brain}\index{breathing}\index{carbon dioxide}\index{diarrhoea}\index{digestion}

Emotions can produce physical symptoms because brain networks---including the
amygdala, insula, and hypothalamus---directly regulate the autonomic nervous
system, stress hormones, skeletal muscle tone, breathing, and immune pathways.
Fear or excitement can raise heart rate, cause sweating or trembling, and speed
breathing; rapid overbreathing lowers blood carbon dioxide, which constricts
cerebral vessels and alters nerve excitability, causing light-headedness,
tingling, or chest tightness. Stress can likewise tighten pericranial or neck
muscles and alter smooth-muscle motility in the gut, leading to headache, chest
discomfort, dizziness, nausea, cramps, diarrhoea, or constipation.

The brain also interprets signals from inside the body through interoceptive
networks. Attention, expectation, memory, and perceived threat can make
sensations feel stronger or weaker, while accompanying changes in heart rate,
breathing, muscles, or digestion may be measurable. Such symptoms are
physically real and are not being consciously faked. Emotion may be one
contributor, however, not a diagnosis by exclusion; new, severe, persistent,
or unexplained symptoms should be medically assessed rather than automatically
attributed to stress.

\section{Why can placebo treatment produce measurable physiological effects?}\index{placebo effect}\index{nervous system!autonomic}\index{bones}\index{brain}\index{cancer}\index{dopamine}

A placebo is an intervention lacking a treatment-specific ingredient for the
condition being treated, but its clinical context can still affect the brain
and body. Expectations, previous experience, conditioning, clinician
communication, and the act of receiving care can change how symptoms such as
pain, nausea, or fatigue are experienced. In some settings, expectancy engages
brain systems involving endogenous opioids, dopamine, and autonomic control,
and can influence measurable responses such as stress or autonomic activity.
The size and biological meaning of the effect vary greatly by condition and
individual; it is often stronger for symptoms than for an underlying disease
process.

An improvement after a placebo is not automatically a placebo effect. Symptoms
can fluctuate naturally, improve through regression to the mean, or change
because of attention and supportive care. Placebos may relieve some symptoms,
but they cannot reliably eradicate an infection, remove a tumour, or repair a
broken bone, and they should not replace effective treatment. Ethical use
requires appropriate informed communication. Negative expectations can produce
a nocebo effect, worsening symptoms or perceived side effects.

\chapter{Pain, Balance, and Whole-Body Coordination}

\section{Why do we feel itch, and how is itch different from pain?}\index{itch (pruritus)}\index{pain}\index{allergies}\index{bones}\index{brain}\index{histamine}

Itch, or pruritus, is an unpleasant sensation that creates an urge to scratch,
often helping remove an irritant from the skin. Pain is a broader warning of
actual or potential tissue threat and usually encourages withdrawal, guarding,
or rest, although neuropathic pain can occur without new tissue damage.

The two sensations also differ in where they arise, how they are carried, and
how they respond to medicines. Itch is largely confined to superficial barrier
tissues---the skin, mucous membranes, and the surface of the eye---whereas pain
can arise both superficially and from deep structures such as muscles, joints,
bones, and internal organs. They use partly different but overlapping sensory
neurons and spinal pathways: itch can be triggered by histamine and by
non-histamine mediators acting on specialised small sensory fibres, and opioid
painkillers suppress pain yet can trigger or worsen itch through distinct spinal
and central circuits. A strong or inflamed stimulus can cause both itch and
pain, and pain from scratching can briefly suppress the itch.

Common triggers of itch include skin irritation, allergy, inflammation,
infection, nerve problems, medications, and some internal diseases such as liver
or kidney disease. Persistent, widespread, or unexplained itch should be
assessed for its cause rather than assumed to be a simple skin irritation.

\section{Why does scratching temporarily relieve an itch?}\index{itch (pruritus)}\index{brain}\index{inflammation}\index{pain}\index{skin}\index{touch}

Scratching activates fast mechanosensory touch and mild pain-sensing fibres,
which stimulate inhibitory interneurons in the dorsal horn of the spinal cord
and temporarily block itch-transmitting neurons from firing. Descending
pathways from the brain can further suppress itch transmission, while the
shift in attention and brief relief can make scratching feel rewarding. Because
the underlying skin irritation, inflammatory mediators, or nerve activity
often remain, the itch returns once inhibition fades.

Hard or repeated scratching damages the skin barrier, increases inflammation
and infection risk, and can make sensory nerves more reactive. Thickening of
the skin may develop, creating a self-reinforcing itch--scratch cycle. Persistent,
severe, widespread, or unexplained itch is better addressed by treating its
cause and should be medically assessed; when appropriate, gentle measures such
as moisturising or cooling may be safer than scratching.

\section{Why can amputees feel pain in a limb that no longer exists?}\index{pain}\index{brain}\index{itch (pruritus)}\index{spinal cord}\index{touch}

The brain keeps a body map of the missing limb even after amputation. Cut
sensory nerves may develop neuromas or generate abnormal impulses, and the
spinal cord can become sensitised. Meanwhile, the brain continues to interpret
signals using its limb map and gradually reorganises connections in response to
the changed sensory input. Mismatch between expected and received visual,
touch, and movement signals may contribute to pain in some people.

A phantom-limb sensation can involve movement, position, touch, temperature,
itch, or pain. Phantom-limb pain is the painful form, whereas residual-limb
pain comes from the remaining stump and may result from a neuroma, poor socket
fit, infection, or another local problem. Both are genuine sensory experiences,
not imagined symptoms. Persistent phantom pain can be managed by clinicians
with an individualised, often multidisciplinary plan that may include
rehabilitation, residual-limb care, medicines, sensory or graded-motor
training, prosthetic optimisation, or mirror-based visual feedback. No single
method works for everyone, and new or worsening stump symptoms need assessment.

\section{Why can chronic pain continue after the original injury heals?}\index{pain}\index{brain}\index{inflammation}\index{sleep}\index{stress}\index{touch}

Pain can continue after an injury appears to have healed for several reasons.
Inflammation, structural problems, recurrent disease, or a damaged nerve may
still be present. Nerve injury can produce burning, electric, or shooting pain
(neuropathic pain). In other cases, prolonged nociceptive input makes the
nervous system more responsive: a normally harmless touch may hurt (allodynia),
or a painful stimulus may feel more intense than expected (hyperalgesia). This
increased responsiveness is called sensitisation and can involve both
peripheral and central mechanisms.

Pain is constructed by the nervous system from incoming signals and their
context, so it is not a precise measure of tissue damage at every moment.
Descending brain control, poor sleep, stress, fear, low mood, illness, and
protective changes in movement can all modify persistence and disability.
Sometimes altered pain processing continues without clear ongoing tissue
damage; this is often called nociplastic pain. These mechanisms can overlap,
and pain without an abnormal scan is still real. However, persistent or
worsening pain should not automatically be attributed to sensitisation: it
needs clinical assessment for treatable causes and an appropriate
multimodal-management plan.

\section{Why do headaches and migraines happen?}\index{headaches}\index{migraine}\index{blood pressure}\index{brain}\index{dehydration}\index{fever}

Although brain tissue itself lacks pain receptors, the head contains many
pain-sensitive structures, including the meninges, blood vessels, cranial and
upper cervical nerves, scalp, and neck muscles. Tension-type headaches are
linked to heightened sensitivity in pericranial muscles and central
pain-processing pathways. Migraines are a distinct neurovascular brain disorder involving altered sensory
processing, activation of the trigeminal nerve system, and release of
vasoactive signalling peptides such as calcitonin gene-related peptide (CGRP)
that sensitise meningeal pain fibres---which is why CGRP-blocking therapies can
prevent or treat migraine attacks. In people who experience visual or sensory
aura, a slow wave of altered cortical activity called cortical spreading
depression typically precedes or accompanies the headache.

Secondary headaches can also result from dehydration, sinus inflammation,
fever, caffeine withdrawal, frequent use of acute painkillers
(medication-overuse headache), or less commonly serious conditions such as
infection, high blood pressure, or intracranial bleeding. A sudden,
``thunderclap'' severe headache, headache with fever, stiff neck, confusion,
rash, weakness, numbness, vision or speech changes, or a new headache after
head injury or in older age requires urgent medical assessment.

\section{Why do humans become dizzy after spinning?}\index{brain}\index{headaches}\index{hearing}\index{inner ear}\index{joints}\index{muscles}

Spinning stimulates the three perpendicular fluid-filled loops of the inner ear
(the \emph{semicircular canals}), which detect head rotation. As the canal
moves, the inner fluid (\emph{endolymph}) lags behind and bends a gelatinous
sensor flap (the \emph{cupula}), changing balance-nerve (\emph{vestibular})
activity. When spinning stops,
the body stops first but the endolymph continues moving briefly. The cupula
then bends in the opposite direction, so the canals signal an apparent turn in
the opposite direction.

This signal conflicts with what the eyes see and with information from muscles,
joints, and the other balance organs. The brain interprets the mismatch as
movement, while vestibular reflexes can make the eyes drift and flick back and
forth involuntarily (\emph{nystagmus}), causing dizziness, a sensation of
rotation, unsteadiness, or nausea. The sensation
usually fades as the fluid settles and the signals agree again. Persistent,
severe, recurrent, or spontaneous vertigo---especially with hearing loss,
weakness, trouble speaking, severe headache, or inability to walk---is not
ordinary post-spinning dizziness and needs prompt medical assessment.

\section{Why cannot the eye see its own blind spot?}\index{brain}\index{optic nerve}\index{retina}\index{vision!visual field}

Each eye has a normal blind spot at the optic disc, where the optic nerve and
retinal blood vessels pass through the retina. This region contains no rods or
cones, so light falling there produces no photoreceptor signal. We usually do
not notice the gap because, with both eyes open, the other eye covers much of
that location. Even with one eye closed, the brain uses nearby visual patterns
and expectations to perceptually fill in the missing information; it does not
receive an image from the blind spot itself.

The blind spot can be demonstrated by closing one eye, fixing the other on a
mark, and moving a small object across the visual field—roughly 15 degrees to
the outer (temporal) side of central fixation—until it disappears. The exact position depends on viewing distance and eye.
A new, enlarging, or otherwise unusual blind spot, especially with flashes,
floaters, or loss of vision, is not the normal physiological finding and should
receive prompt eye assessment.

\section{Why does the pupil constrict in bright light?}\index{pupil}\index{nervous system!autonomic}\index{headaches}\index{muscles}\index{pain}\index{retina}

The pupil acts like an adjustable opening. In bright light, photoreceptors and
intrinsically photosensitive retinal ganglion cells signal the pretectal area of
the midbrain. The pretectal area activates both Edinger--Westphal nuclei, whose parasympathetic
fibres travel with the oculomotor nerves to the ciliary ganglia and then to the
iris sphincter muscles. Their contraction constricts the pupils, limiting light
entry, reducing glare, and often increasing depth of field; it does not make
the retina itself more sensitive.

The reflex normally constricts both pupils, even when light is shone into only
one eye, producing direct and consensual responses. In darkness, sympathetic
activity to the iris dilator muscles generally enlarges the pupils. Pupil size
is also influenced by near focusing, emotion, arousal, drugs, and other
autonomic signals. A new unequal pupil response, especially with eye pain,
drooping eyelid, double vision, severe headache, or neurological symptoms, is
not a normal light reflex and warrants medical assessment.

\section{Why is the retina effectively part of the central nervous system?}\index{retina}\index{brain}\index{cornea}\index{embryonic development}\index{meninges}\index{neurons}

The retina is effectively part of the central nervous system because it develops
as an outgrowth of the embryonic forebrain. It contains photoreceptors, several
classes of interneurons, and retinal ganglion cells arranged in neural circuits,
not merely a sheet of light-sensitive cells. These circuits begin processing
contrast, timing, colour, motion, and spatial information before signals travel
to the brain through the optic nerve.

The optic nerve is a CNS tract: its axons are insulated by CNS myelin-forming
cells (\emph{oligodendrocytes}) and surrounded by protective brain membranes
(\emph{meninges}) and cerebrospinal-fluid spaces, unlike a typical peripheral
nerve. Inside the retina itself, a blood--retinal barrier regulates molecular
exchange much like the blood--brain barrier, and CNS-characteristic support
cells (including star-shaped \emph{astrocytes}, immune \emph{microglia}, and
retinal \emph{M\"uller glia}) support and maintain the neurons. The retina and
optic nerve therefore share important features with brain tissue, including
limited repair after injury—which is why eye doctors can sometimes spot signs
of systemic brain, inflammatory, or vascular disease directly on a retinal
examination. This comparison applies to the neural retina and optic nerve, not to
every part of the eye; the lens, cornea, sclera, and many supporting tissues
are not themselves parts of the CNS.

\section{Why cannot damaged human retinal neurons readily regenerate?}\index{neurons}\index{regeneration}\index{brain}\index{genes}\index{inflammation}\index{optic nerve}

Mature human retinal neurons no longer divide (\emph{post-mitotic}) and normally
cannot replace themselves after they are lost. Retinal ganglion cells, whose
axons form the optic nerve, would also need to regrow a long axon, follow the
correct route, and reconnect with precise targets in the brain. After injury,
neurons may lose intrinsic growth capacity or die, and inflammation, inhibitory
surrounding scaffolding signals (\emph{extracellular-matrix} signals), and
scar-like tissue can discourage regrowth.
Restoring a useful circuit would require reconnecting many cell types in the
correct pattern, not merely producing new cells.

Some animals can use supporting Müller glia to regenerate retinal neurons, but
adult human Müller glia do not normally replace large numbers of lost neurons.
Cell replacement, gene-based treatments, neuroprotection, and optic-nerve
regrowth are active research areas. Some therapies can slow or treat particular
retinal diseases, but they do not yet routinely restore established loss of
human retinal neurons or optic-nerve connections; many regenerative approaches
remain experimental or condition-specific, making early eye checks and prompt
treatment of sudden vision changes essential to protect remaining sight.

\section{How can 30--40 trillion human cells coordinate as one organism?}\index{ageing}\index{blood pressure}\index{cancer}\index{genes}\index{glucose}\index{hormones}

The body's cells coordinate by sending and receiving chemical, electrical, and
mechanical signals. Nearby cells communicate by direct contact, adhesion,
electrical coupling through gap junctions, or short-range chemicals; nerves
provide rapid, targeted messages; hormones travel in the blood to distant
organs; and immune cells use cytokines and contact signals to identify threats.
Local nutrients, oxygen, metabolites, and the extracellular matrix also convey
information. A cell responds only when it has the appropriate receptor,
channel, adhesion molecule, and intracellular machinery.

Most cell types share much of the same genome but activate different gene
programmes, giving them specialised jobs. Tissues and organs are organised so
that local interactions and long-range control cooperate. Feedback systems
continually adjust temperature, blood pressure, oxygen, glucose, fluid balance,
and tissue repair. The nervous, endocrine, immune, and circulatory systems
exchange information at different speeds, allowing the body to act as one
changing, coordinated organism rather than as identical cells following a
single command. Coordination is distributed and includes redundancy, not a
single central controller.

The estimate of 30--40 trillion cells is approximate and varies with body size,
age, sex, tissue composition, and counting method. Coordination is not perfect:
mutation, infection, ageing, cancer, or failed communication can disrupt it and
cause disease.

\chapter{Essential Foundations and Health Literacy}

\section{What is a cell, and how do its major parts work?}\index{ATP (adenosine triphosphate)}\index{calcium}\index{cell membrane}\index{chromosomes}\index{endoplasmic reticulum}\index{genes}

A cell is the smallest unit that can carry out the basic processes of life. Its
plasma membrane controls exchange with the surroundings and carries receptors;
the nucleus stores most of the genetic material within a porous nuclear
envelope, houses the nucleolus where ribosomal subunits are assembled, and
regulates gene expression; ribosomes build proteins; and mitochondria use
nutrients and oxygen to make much of the cell's ATP. The endoplasmic reticulum
synthesises membrane and secreted proteins on its rough regions and makes lipids
and stores calcium in its smooth regions, while the Golgi apparatus modifies,
sorts, and packages those products into vesicles for delivery or secretion.
Lysosomes digest selected materials, peroxisomes break down certain fatty acids
and neutralise toxic peroxides, centrosomes help organise microtubules and
separate chromosomes during cell division, and the cytoskeleton supports shape,
transport, and movement. The cytoplasm and internal compartments provide an
organised setting for chemical reactions.

Different cells use the same broad plan but specialise through their shape,
internal structures, gene activity, and connections. Not every cell has every
organelle: mature red blood cells, for example, lack a nucleus and mitochondria.
Cells also sense signals, communicate, repair damage, divide, or die in a
controlled way. Tissues and organs arise when specialised cells cooperate.

\section{How are DNA, genes, chromosomes, and proteins related?}\index{chromosomes}\index{DNA}\index{genes}\index{proteins}\index{blood!red blood cells}\index{sperm}

DNA is the molecule that stores inherited biological information. A gene is a
functional region of DNA that helps specify an RNA or protein product, or helps
regulate a process; not all genes make proteins. A chromosome is a long,
organised DNA molecule packaged with proteins, and it contains many genes and
regulatory regions. In humans, most cells have 23 pairs of chromosomes, with
exceptions such as mature red blood cells and reproductive cells.

Genes are expressed through transcription of DNA into RNA and, for many genes,
translation of messenger RNA into protein. Alternative RNA processing can let
one gene contribute to more than one product. Cells with the same genome can
behave differently because they activate different genes, regulated by control
sequences, transcription factors, chromatin state, chemical marks on DNA and its
associated proteins, and signals from other cells. Many traits arise from
networks of genes interacting with development and environment rather than from
one gene alone.

Mutations can alter DNA, with effects ranging from harmless to beneficial or
harmful depending on the gene, cell type, and change. Changes in DNA can occur
in body cells without being inherited, whereas changes in egg, sperm, or their
precursor lineages may be passed to descendants.

\section{How do cells produce and use energy?}\index{ATP (adenosine triphosphate)}\index{mitochondria}\index{oxygen}\index{proteins}

Cells obtain energy from nutrients and convert much of it into ATP, a rapidly
recycled molecule that transfers energy to cellular work. Mitochondria usually
produce most ATP by using oxygen to extract energy from carbohydrates, fats,
and other fuels; glycolysis and other substrate-level reactions can produce
some ATP without oxygen. The amount and route vary among tissues: mature red
blood cells, for example, lack mitochondria and rely on glycolysis.

ATP powers processes such as movement, transport across membranes, synthesis of
proteins, and maintenance of electrical signals. Its terminal phosphate can be
transferred during a reaction, after which ADP is recharged to ATP. Energy
conversion is not perfectly efficient, so some energy is released as heat.

With oxygen available, pyruvate from glycolysis enters mitochondria, where it
is converted to a two-carbon fuel unit (acetyl-CoA) and processed through the
citric-acid cycle. Fatty acids can be broken down step by step
(\emph{beta-oxidation}) to supply acetyl-CoA as well. High-energy electrons from
these pathways pass through the mitochondrial respiratory chain, which builds a
proton gradient across the inner membrane that drives the molecular turbine
ATP synthase—yielding roughly 30--32 molecules of ATP per glucose. Without
enough oxygen, cells can make ATP briefly through glycolysis while converting
pyruvate to lactate, regenerating the electron carrier NAD+ needed to keep
glycolysis running; this yields only 2 ATP per glucose and cannot support all
tissues for long.

\section{How does the cell membrane control what enters and leaves?}\index{cell membrane}\index{ATP (adenosine triphosphate)}\index{carbon dioxide}\index{glucose}\index{muscles}\index{oxygen}

The cell membrane is a selectively permeable barrier made mainly of a flexible
lipid bilayer with proteins embedded in it. Small non-polar or lipid-soluble
molecules, including oxygen and carbon dioxide, may cross directly, whereas
channels and transporters help move ions, glucose, water, and other substances.
Water can move by osmosis toward the side with a higher effective solute
concentration, often through aquaporins. Facilitated diffusion moves a
substance down its electrochemical gradient without directly using ATP, whereas
primary active transport uses energy, commonly from ATP---as in the
sodium--potassium pump---to move material against a gradient, and secondary
active transport uses one gradient to drive another substance uphill.

The ion gradients built by membrane pumps help create the electrical signals
used by nerves and muscles, so membrane transport supports communication and
electrical activity as well as nutrition and waste removal. Cells can also bring
larger material in by endocytosis and release material by exocytosis, while
membrane proteins additionally provide receptors, adhesion, and signalling.

\section{How does the body maintain homeostasis?}\index{homeostasis}\index{blood pressure}\index{brain}\index{carbon dioxide}\index{blood!clotting}\index{glucose}

Homeostasis is the continuous regulation of internal conditions within ranges
compatible with life. Sensors detect changes such as altered temperature,
carbon dioxide, blood pressure, glucose, or osmolality. Integrating circuits in
the brain, spinal cord, endocrine organs, or local tissues compare information
with a regulated range and coordinate effectors that adjust organs and blood
flow. Some responses are anticipatory, beginning before a variable changes
substantially.

Most regulation uses negative feedback, which opposes a change. The body does
not keep every value fixed, and regulated ranges can shift with sleep,
exercise, illness, age, hormones, and other circumstances. This flexible
adjustment is sometimes called allostasis. Failure of sensing, integration, or
effector response can contribute to disease.

Control is distributed rather than located in one master switch. The brain,
endocrine glands, kidneys, liver, lungs, immune system, and local tissue
circuits all contribute. Some processes use positive feedback for a limited
time---for example, clotting and labour contractions---so homeostasis means
stable function over time, not that every response always reverses the initial
change.

\section{How do hormones regulate the body?}\index{hormones}\index{circadian rhythm}\index{genes}\index{hypothalamus}\index{metabolism}\index{stress}

Hormones are chemical messengers released by specialised cells, usually into
the bloodstream for endocrine action. They travel widely, but only cells with
the appropriate receptors and signalling machinery respond. Hormones can alter
metabolism, growth, reproduction, water balance, stress responses, and many
other functions.

Hormonal effects are often slower and longer-lasting than many nerve signals,
though some hormones act within seconds. Secretion may be pulsatile or
circadian and is commonly regulated by negative-feedback loops, such as the
hypothalamus--pituitary--thyroid and stress axes. The nervous, endocrine, and
immune systems communicate with one another rather than acting as separate
systems.

Water-soluble hormones usually bind receptors on the cell surface and activate
second-messenger or kinase pathways. Steroid hormones and thyroid hormones
usually act through intracellular receptors that alter gene activity, although
rapid effects are also possible. In both cases, the overall response depends on
hormone concentration, free hormone availability, receptor number and
sensitivity, target-cell state, and the rate of hormone breakdown and clearance.
Nearby cells can also communicate by paracrine signals, and a cell can respond
to its own autocrine signals; these are related forms of chemical communication
rather than circulating endocrine hormones.

\section{What are the components of blood, and what does each do?}\index{antibodies}\index{carbon dioxide}\index{blood!clotting}\index{haemoglobin}\index{hormones}\index{inflammation}

Plasma is the liquid portion of blood and carries water, salts, nutrients,
hormones, proteins, and waste products. Albumin helps maintain oncotic pressure
and transports molecules; immunoglobulins support defence; and fibrinogen and
other clotting proteins help stop bleeding. Red blood cells transport oxygen
using haemoglobin and carry carbon dioxide mainly after it has been converted
to bicarbonate. White blood cells support innate and adaptive immune defence,
while platelets help initiate clotting.

These components work together: blood transports materials, regulates internal
conditions, defends against threats, and repairs vessel injury. The proportion
of blood occupied by red cells is the haematocrit, which affects oxygen-carrying
capacity and viscosity. Abnormal numbers or function of any component can
impair health.

Red cells are flexible, disc-shaped cells without nuclei or mitochondria when
mature, leaving room for haemoglobin but limiting their ability to repair
themselves. Platelets are cell fragments released by megakaryocytes, and
clotting proteins form a reinforcing fibrin network. White cells include
neutrophils, monocytes, lymphocytes, eosinophils, and basophils, with roles in
rapid defence, inflammation, antibody production, and immune memory; a total
white-cell count alone cannot describe immune function.

\section{How does blood pressure work, and what causes hypertension?}\index{blood pressure}\index{hypertension}\index{arteries}\index{brain}\index{cardiac output}\index{heart}

Blood pressure is the force of circulating blood against the vessel walls. Mean
arterial pressure depends largely on cardiac output and systemic vascular
resistance: cardiac output is heart rate multiplied by stroke volume, while
resistance is strongly influenced by small-artery diameter. Arterial elasticity,
blood volume, blood viscosity, and the kidneys and hormones that regulate salt
and water also matter. Pressure normally rises and falls with each heartbeat
and varies with activity, posture, sleep, and emotion; systolic and diastolic
pressure therefore describe different parts of the pulsatile signal.

Hypertension is persistently elevated blood pressure confirmed with appropriate
repeated measurements (typically $\ge$ 130/80~mmHg or $\ge$ 140/90~mmHg
depending on clinical guidelines, compared with a normal resting reading below
120/80~mmHg), not one isolated reading. Most cases are primary
hypertension arising from interacting inherited, vascular, renal, metabolic,
and environmental influences. Secondary causes include kidney or endocrine
disease, sleep apnoea, certain medications or substances, and some vascular
conditions. Hypertension can damage the heart, brain, kidneys, retina, and
arteries while causing no obvious symptoms, so measurement and risk assessment
are important; very high readings accompanied by chest pain, shortness of
breath, severe headache, vision changes, or weakness indicate a hypertensive
emergency requiring immediate medical care. White-coat and masked hypertension
also show why readings need context and, sometimes, home or ambulatory
monitoring.

\section{What causes heart attacks and atherosclerosis?}\index{atherosclerosis}\index{heart}\index{arteries}\index{blood pressure}\index{blood!clotting}\index{diabetes}

Atherosclerosis is a slow, chronic process in which cholesterol-carrying
lipoproteins enter the inner lining of arteries, trigger local inflammation,
and form fatty, fibrous plaques within the vessel wall. Over years, these
plaques can narrow arteries and restrict blood flow, for example causing
exertional chest pain (angina) when the heart muscle's oxygen demand rises.

A heart attack (\emph{myocardial infarction}) most often occurs abruptly when an
unstable plaque ruptures or erodes, exposing material that triggers a blood clot
inside a coronary artery. If the clot blocks blood flow, downstream heart
muscle cells begin dying within minutes from oxygen deprivation and are later
replaced by non-contractile scar tissue. Risk is shaped by LDL cholesterol,
high blood pressure, smoking, diabetes, age, and genetics. Sudden chest pain,
pressure, or tightness---especially spreading to the arm, neck, jaw, or back,
or accompanied by breathlessness, cold sweats, nausea, or faintness---requires
emergency medical care; symptoms can be less typical in women, older adults,
and people with diabetes.

\section{How does haemoglobin transport oxygen?}\index{haemoglobin}\index{oxygen}\index{carbon dioxide}\index{lungs}\index{placenta}\index{proteins}

Haemoglobin is a protein inside red blood cells with four haem groups, each
containing iron that can bind one oxygen molecule reversibly. In the lungs,
where oxygen partial pressure is high, it loads oxygen; in tissues, where oxygen
is being used, it releases more of it. Haemoglobin also carries some carbon
dioxide and helps buffer blood acidity. Low haemoglobin, abnormal haemoglobin,
or too few red blood cells can reduce oxygen delivery even when the lungs are
working normally.

Haemoglobin binds oxygen cooperatively: binding at one site changes the protein
so that other sites bind more readily, producing a sigmoidal saturation curve.
Acidity, carbon dioxide, temperature, and 2,3-BPG encourage unloading in active
tissues (the Bohr effect). Some carbon dioxide binds directly to haemoglobin,
while most is transported as bicarbonate through red-cell reactions. Fetal
haemoglobin has a different affinity that helps transfer oxygen across the
placenta.

Carbon monoxide binds haemoglobin much more strongly than oxygen and also
reduces the remaining sites' ability to release oxygen. Exposure is therefore
an emergency even when skin colour appears normal or a standard pulse
oximeter gives a misleadingly reassuring reading.

\section{What roles do the liver, pancreas, and gallbladder play in digestion?}\index{digestion}\index{gallbladder}\index{liver}\index{pancreas}\index{bile}\index{glucagon}

The liver makes bile acids and other bile components, processes absorbed
nutrients, stores glycogen, some vitamins, and iron, and synthesises important
plasma proteins. It also processes substances absorbed from the intestine
through the portal vein before they reach the general circulation, converts
ammonia to urea, and modifies many drugs, hormones, and toxins for excretion in
bile or urine---though ``detoxify'' does not mean it makes every chemical
harmless. The gallbladder stores and concentrates bile; fat and acid entering
the small intestine stimulate signals including cholecystokinin, which promotes
gallbladder contraction and bile delivery. Rather than digesting fat directly,
bile salts disperse fat droplets so pancreatic lipase can act on them and help
form micelles for absorption.

The exocrine pancreas releases bicarbonate to neutralise gastric acid and
digestive enzymes into the duodenum. These include amylase, lipase, nucleases,
and proteases, many of which are released as inactive precursors to protect the
pancreas. Its endocrine islets release hormones such as insulin and glucagon
into the blood to regulate fuel metabolism. Together, these organs support
both digestion and whole-body metabolism, and injury to them can impair both
functions.

\section{What does the large intestine do?}\index{large intestine}\index{amino acids}\index{diarrhoea}\index{digestion}\index{small intestine}\index{vitamins}

The large intestine---the caecum, colon, rectum, and anal canal---absorbs water
and electrolytes, especially sodium and chloride, from the remaining intestinal
contents while secreting potassium and bicarbonate, helping maintain fluid and
acid--base balance as it compacts waste into stool. Its resident microbes ferment
some undigested fibre and resistant carbohydrate, producing short-chain fatty
acids that nourish colon cells and influence immune and metabolic signalling,
along with some vitamin K and other compounds of varying nutritional
importance. Mucus lubricates stool for passage and, together with the
epithelium and immune tissue, helps contain the microbial community.

The colon stores stool and uses coordinated motility to move it onward; when
stool stretches the rectum, sensory nerves trigger the urge to defaecate, and
the internal and external anal sphincters coordinate continence and voluntary
elimination. The large intestine does not perform all digestion, as most
calories, amino acids, and readily absorbable nutrients are absorbed earlier in
the small intestine. Too little absorption or excessive secretion can contribute
to diarrhoea and potassium and bicarbonate loss, while excessive water removal
or slow transit can contribute to hard stool and constipation.

\section{How does the body maintain salt, water, and acid--base balance?}\index{carbon dioxide}\index{consciousness}\index{heart}\index{hormones}\index{kidneys}\index{lungs}

The kidneys adjust how much water, sodium, chloride, potassium, and other
substances are returned to the blood or excreted in urine. Hypothalamic concentration sensors (\emph{osmoreceptors}) and cardiovascular
volume sensors influence thirst and antidiuretic hormone (vasopressin), which
makes the kidney's final urine tubes (\emph{collecting ducts}) more permeable
to water. Reduced renal perfusion or sodium chloride delivery can activate the
renin--angiotensin--aldosterone hormone cascade (renin, angiotensin II, and
aldosterone), promoting sodium retention and altering potassium and acid
secretion. Water balance and sodium balance are linked but are not controlled by
one identical mechanism.

The bicarbonate, protein, and phosphate buffer systems resist rapid changes in
arterial blood pH (normally kept tightly within 7.35--7.45). The lungs provide relatively rapid regulation by changing how much
carbon dioxide is exhaled: retaining CO$_2$ tends to increase acidity, while
exhaling more tends to reduce it. The kidneys provide slower, more durable
control by reabsorbing or generating bicarbonate, excreting hydrogen ions, and
producing ammonium to carry acid into urine. Severe failure of these systems
can disturb fluid volume, electrolytes, nerves, muscles, consciousness, and
heart rhythm.

\section{Why do humans need water?}\index{blood pressure}\index{breathing}\index{consciousness}\index{digestion}\index{fever}\index{hormones}

Water is the body's main solvent and transport medium. It allows chemical
reactions to occur inside cells, helps maintain cell volume and acid--base
conditions, carries nutrients and hormones through the blood, helps remove
waste in urine, lubricates and cushions tissues, and supports digestion. Its
high heat capacity helps buffer temperature changes, while evaporation from
the skin cools the body. Water also helps maintain circulating volume and
normal blood pressure.

Water is continually lost through urine, faeces, sweat, and breathing, including
insensible losses that are not always noticed. The body replaces it through
drinks, food, and a small amount made during metabolism. Thirst, hypothalamic
sensors, antidiuretic hormone, and the kidneys work together to adjust intake
and conserve or excrete water, although illness, age, medications, or impaired
consciousness can make these controls unreliable. Needs vary with body size,
activity, climate, fever, pregnancy, breastfeeding, and diet; there is no
single amount correct for everyone. Drinking far more water than the kidneys
can excrete can dilute blood sodium and cause dangerous hyponatraemia, so more
is not always better.

\section{What happens when the kidneys fail?}\index{kidneys}\index{anaemia}\index{blood pressure}\index{bones}\index{diabetes}\index{erythropoietin}

When kidney function falls substantially, the body has difficulty removing
waste, regulating water and electrolytes, controlling acid--base balance, and
performing endocrine functions. Fluid, potassium, acids, and nitrogenous waste
can accumulate, causing swelling, breathlessness, confusion, weakness, or
dangerous heart-rhythm disturbances, although early disease may produce few
symptoms. Reduced renin regulation, erythropoietin production, and activation
of vitamin D can also disturb blood pressure, cause anaemia, and impair bone
and mineral balance.

Acute kidney injury can develop over hours or days and may improve when its
cause is treated, although it can leave lasting damage. Chronic kidney disease
develops over months or years and may be stable or progressively reduce
filtering capacity. Testing is important when risk is present because symptoms
often appear late.

Treatment depends on the cause and severity and may include treating an
obstruction or infection, medicines, dietary changes, management of blood
pressure or diabetes, dialysis, or transplantation. Dialysis replaces selected
kidney functions; it is not a complete substitute for every endocrine or
metabolic role. Kidney failure is a medical condition requiring professional
care and is not corrected by simply drinking more water.

\section{Why do allergies occur?}\index{allergies}\index{anaphylaxis}\index{antibodies}\index{breathing}\index{histamine}\index{immune system}

An allergy is an inappropriate or excessive immune response to a usually
harmless substance, such as pollen, food proteins, medicines, or insect venom.
During sensitisation, the immune system develops a response to the allergen;
later exposure can activate antibodies or immune cells and release mediators
such as histamine, causing symptoms in the skin, nose, lungs, or gut. Genetic
susceptibility, barrier function, environment, and the route and dose of
exposure all influence risk.

In many immediate allergies, allergen-specific IgE antibodies attach to mast
cells and basophils; re-exposure cross-links the IgE and triggers release of
histamine and other mediators. These can produce itching, hives, swelling,
bronchoconstriction, vomiting, or widespread vasodilation. Other reactions use
different antibody or T-cell pathways and may be delayed. An allergy is not the
same as an intolerance, which can cause symptoms without the same immune
mechanism.

Reactions range from mild local symptoms to life-threatening anaphylaxis.
Sudden breathing difficulty, throat or tongue swelling, faintness, or symptoms
in more than one body system after an exposure requires emergency treatment;
people with a prescribed epinephrine auto-injector should follow their
emergency plan. Severe reactions can recur after an initial improvement.

\section{How do antibiotics work, and why does resistance develop?}\index{antibiotics}\index{DNA}\index{genes}\index{metabolism}\index{proteins}

Antibiotics target structures or processes used by bacteria, such as cell-wall
construction, protein production, DNA replication, folate metabolism, or
membrane function. Some inhibit growth and others kill bacteria. They do not
treat viruses, and different antibiotics work against different organisms;
choice depends on the likely or confirmed pathogen, site of infection,
resistance patterns, and the patient.

Resistance develops when genetic changes or acquired genes allow some bacteria
to survive an antibiotic. Mechanisms include drug destruction or modification,
altered drug targets, reduced entry, and active efflux. Antibiotic exposure
selects for resistant organisms, which can multiply and spread; bacteria can
also exchange resistance genes by horizontal gene transfer. Unnecessary,
incorrectly chosen, or poorly administered treatment increases this pressure.

Appropriate prescribing, vaccination and infection prevention, diagnostic
testing when useful, and taking a prescribed course as directed help limit
resistance. Patients should not save leftover antibiotics, share them, or
self-treat viral illnesses. The correct drug and duration depend on the
infection, so treatment should not be stopped, extended, or changed without
clinical advice; side effects should be reported. Antibiotics can also disrupt
beneficial microbes and cause adverse reactions, so their benefits and harms
must be weighed.

\section{How do neurons communicate at synapses?}\index{neurons}\index{synapses}\index{calcium}\index{learning}\index{memory}\index{neurotransmitters}

When an action potential reaches a chemical synapse, it opens voltage-gated
calcium channels in the axon terminal. Calcium entry triggers docking proteins
(\emph{SNARE proteins}) to fuse neurotransmitter-filled vesicles with the
membrane, releasing their contents into the synaptic cleft. The transmitter
crosses this narrow space and binds to receptors on the receiving cell.
Fast channel-opening (\emph{ionotropic}) receptors open ion pores quickly;
slower G-protein-coupled (\emph{metabotropic}) receptors act through internal
second messengers, usually more slowly and for longer.

The resulting local change may be excitatory (depolarising, an EPSP),
inhibitory (hyperpolarising, an IPSP), or modulatory, and the same
neurotransmitter can excite one cell while inhibiting another depending on
receptor type, ion gradients, and target-cell state. A neuron combines inputs
arriving at different synapses (spatial summation) and at different times
(temporal summation); the net effect near the trigger zone where impulses launch
(\emph{axon initial segment}) helps determine whether it fires. The transmitter
is then removed by reuptake into neurons or glial cells, enzymatic breakdown, or
diffusion—steps targeted by many medicines (such as antidepressants that slow
serotonin reuptake) and neurotoxins.

By contrast, electrical synapses use
cell-to-cell protein channels (\emph{connexin}-containing gap junctions) to pass
ionic current directly between cells; they are extremely rapid, may conduct in
both directions, and can synchronize activity, though they provide less signal
amplification and flexibility than most chemical synapses.

\section{How does learning change the brain?}\index{brain}\index{learning}\index{genes}\index{hippocampus}\index{hormones}\index{memory}

Learning changes how existing neural circuits respond and communicate. Repeated
co-activation can strengthen synapses (long-term potentiation), whereas other
patterns can weaken them (long-term depression). These changes may involve
altered neurotransmitter release, receptor number or sensitivity, dendritic
spines, gene expression, and protein synthesis. Practice can also recruit new
connections, refine circuit activity, and alter myelination, which can change
the timing and efficiency of communication.

Different kinds of learning use partly different systems. The hippocampus is
important for forming many new explicit memories, while practice can gradually
make skills more dependent on distributed cortical and subcortical circuits.
Attention, motivation, emotion, reward, prior knowledge, and feedback influence
what is encoded. Sleep supports consolidation through coordinated replay and
network changes; retrieval can both strengthen a memory and make it temporarily
modifiable before reconsolidation.

Neuroplasticity is regulated rather than unlimited. Developmental sensitive
periods, age, hormones, health, injury, and the amount and spacing of practice
affect its direction and extent. Plasticity can be beneficial, but it can also
reinforce pain, addiction, fear, or maladaptive habits. Learning does not usually
mean that large numbers of new neurons appear: in adults, limited
neurogenesis has been reported in some regions, but most learning depends on
changes in existing cells and networks. A biological change is not by itself
proof that a particular memory or skill has been stored in one location.

\section{What happens during a stroke or seizure?}\index{seizures}\index{stroke}\index{ATP (adenosine triphosphate)}\index{nervous system!autonomic}\index{brain}\index{breathing}

A stroke is an acute brain injury caused either by blocked blood flow (an
ischaemic stroke) or by bleeding from a ruptured vessel (a haemorrhagic stroke).
In ischaemia, the deprived region loses oxygen and glucose: ATP production
fails, ion gradients collapse, cells swell, and excitotoxic and inflammatory
processes can expand the injury. Some surrounding tissue may remain temporarily
viable, so the extent and reversibility of damage depend on blood flow, location,
collateral vessels, and time. Haemorrhage additionally causes pressure and
irritation from blood within or around the brain. Symptoms may include sudden
facial drooping, arm or leg weakness, speech or vision disturbance, imbalance,
severe headache, or confusion.

A seizure is a period of excessive, abnormally organised or disorganised
electrical activity in a network of brain cells. Excess excitation, reduced
inhibition, or disrupted ion balance can make neurons fire in an abnormal
pattern. Seizures may cause convulsions, staring and impaired awareness,
unusual sensations, involuntary movements, or temporary behavioural or
autonomic changes; not all involve shaking. A temporary trigger such as fever,
low glucose, withdrawal, or an acute brain injury can provoke one seizure, while
epilepsy is a continuing predisposition to recurrent unprovoked seizures.

Stroke and seizure are therefore different processes, although a stroke can
trigger a seizure and a seizure can briefly produce weakness or other
stroke-like postictal symptoms. A transient ischaemic attack causes temporary
stroke symptoms without demonstrable lasting infarction, but it is still an
urgent warning. Sudden neurological symptoms require emergency assessment;
brain imaging, blood tests, and sometimes EEG help distinguish these conditions.
A seizure lasting about five minutes, repeated seizures without recovery, or
breathing difficulty is also an emergency.

\section{How do hearing and speech work?}\index{hearing}\index{speech}\index{bones}\index{brainstem}\index{brain}\index{cochlea}

Sound is collected by the outer ear and makes the eardrum and middle-ear bones
vibrate. The stapes passes this motion through the oval window into cochlear
fluid. A travelling wave bends hair cells in the organ of Corti: inner hair
cells convert the movement into signals in the auditory nerve, while outer hair
cells amplify and sharpen the mechanical response. Different positions along
the cochlea respond preferentially to different frequencies (tonotopy), while
the size and timing of activity help encode loudness and other features.

Signals travel through brainstem and thalamic relays to the auditory cortex,
where the brain analyses patterns as speech, music, or environmental sound.
Inputs from both ears support sound localisation, and the brain uses context
to separate a voice from background noise. Hearing is therefore an active
perceptual process, not simply a microphone-like recording of the ear's output.

Speech production requires breath control, vibration of the vocal folds, and
precise movements of the tongue, lips, jaw, and soft palate. Distributed brain
networks select words and meanings, plan motor sequences, and monitor auditory
and sensory feedback. Hearing helps language and articulation develop, but
speech and language are not identical: hearing loss, aphasia (impaired
language), dysarthria (weak or poorly coordinated speech muscles), and apraxia
of speech can affect communication through different mechanisms.

\section{How do fertilisation, implantation, and placental support occur?}\index{fertilisation}\index{implantation}\index{placenta}\index{amino acids}\index{antibodies}\index{capillaries}

Fertilisation usually occurs in the ampulla of a fallopian tube. After sperm
undergo surface priming in the female tract (\emph{capacitation}), one sperm
penetrates the egg's surrounding cells and protective outer coat (\emph{zona
pellucida}); the egg then changes its surface to block additional sperm. The
sperm and egg nuclei form parental \emph{pronuclei}, which combine to create a
zygote. The zygote undergoes cell divisions as it travels toward the uterus,
becoming a solid ball of cells (\emph{morula}) and then a hollow sphere
(\emph{blastocyst}) with an inner cell mass and an outer placental-precursor
shell (\emph{trophoblast}).

Around six to seven days after fertilisation, the blastocyst begins attaching to
the receptive endometrium. Trophoblast cells invade the lining and later form
major parts of the placenta; the inner cell mass forms the embryo. Implantation
normally occurs in the uterus, whereas implantation elsewhere (most often in a
fallopian tube) is \emph{ectopic} and can be life-threatening—severe one-sided
pelvic pain, shoulder-tip pain, dizziness, or vaginal bleeding in early
pregnancy requires urgent assessment. Embryonic cells then fold into three
foundational layers (\emph{gastrulation}) and undergo coordinated patterning to
form germ layers, organs, and supporting membranes.

The placenta creates a highly vascular exchange organ linked to the fetus by
the umbilical cord. Maternal blood enters spaces around chorionic villi, while
fetal blood remains in capillaries within the villi; their proximity permits
diffusion and carrier-mediated transport without normal free mixing. Oxygen,
water, glucose, amino acids, lipids, and selected minerals cross toward the
fetus, while carbon dioxide and wastes cross toward the mother. Maternal IgG
antibodies can cross, especially later in pregnancy, but transfer is selective
and does not provide complete protection.

The placenta also produces or supports hormones such as hCG, progesterone,
oestrogens, and human placental lactogen, adapts maternal metabolism, and
contributes to immune tolerance. Its barrier and transport functions change
through pregnancy, and placental disease can impair fetal growth or maternal
health.

\section{How can readers distinguish normal symptoms from medical warning signs?}\index{allergies}\index{breathing}\index{headaches}\index{pain}\index{seizures}\index{sleep}

Many symptoms have harmless explanations and vary with age, exertion, sleep,
stress, diet, menstrual or reproductive state, and known conditions. Context,
severity, duration, recurrence, and change from a person's usual pattern matter
more than a symptom label alone. A symptom that is new, unexplained, persistent,
or steadily worsening deserves clinical advice, particularly in an infant,
older adult, pregnant or postpartum person, or someone with major chronic
disease or immune suppression.

Seek emergency help for severe trouble breathing, blue or grey lips, new
confusion or unresponsiveness, fainting with ongoing symptoms, severe chest
pressure, sudden one-sided weakness or speech trouble, a sudden severe
headache, uncontrolled bleeding, a serious allergic reaction with airway or
circulatory symptoms, or a seizure lasting about five minutes or recurring
without recovery. Severe abdominal pain, rapidly worsening infection symptoms,
or sudden loss of vision can also require urgent assessment. The appropriate
local emergency service should be contacted; a person with a possible
life-threatening illness should not delay care while searching for an online
explanation.

Red flags do not identify one diagnosis, and their absence cannot guarantee
safety. Conversely, a symptom can be important even when it is mild if it is
recurrent or represents a major change. This is general safety guidance, not a
diagnosis: a clinician may need a history, examination, and targeted testing.

\section{How does breathing supply oxygen and remove carbon dioxide?}\index{breathing}\index{carbon dioxide}\index{oxygen}\index{alveoli}\index{brain}\index{capillaries}

When the diaphragm contracts and descends, supported by the external
intercostal muscles, chest volume increases and air flows into the lungs. When
these muscles relax, elastic recoil usually drives quiet exhalation; forced
breathing recruits additional muscles. Surfactant lowers surface tension and
helps keep alveoli open. Air moves through conducting airways before reaching
the respiratory zone, where alveoli lie beside a dense capillary network.

Because of differences in partial pressure, oxygen diffuses across the thin
air--blood barrier into pulmonary capillaries and binds mainly to haemoglobin.
Carbon dioxide travels in blood mostly as bicarbonate, with smaller amounts
dissolved or bound to proteins; it is converted back to carbon dioxide in the
lungs and diffuses into alveoli for exhalation. The circulation then delivers
oxygenated blood to tissues, where the gradients are reversed. Oxygen content
therefore depends not only on lung oxygen transfer but also on haemoglobin and
cardiac output.

Breathing rhythm is generated and adjusted by networks in the medulla and pons.
Central chemoreceptors respond chiefly to carbon-dioxide-driven changes in
brain-fluid acidity, while peripheral carotid and aortic bodies detect low
arterial oxygen as well as changes in carbon dioxide and pH. Stretch and other
sensory signals refine the response. Ventilation, pulmonary blood flow,
diffusion, and matching of air to blood flow (the ventilation--perfusion ratio)
must all be adequate. Airflow obstruction, fluid or thickened alveolar walls,
poor perfusion, or inadequate respiratory drive can therefore impair oxygen
uptake, carbon-dioxide removal, or both.

\section{How does digestion turn food into usable nutrients?}\index{digestion}\index{amino acids}\index{bile}\index{gallbladder}\index{hormones}\index{large intestine}

Digestion combines movement, chemical breakdown, absorption, and coordinated
transport. Chewing and saliva begin mechanical processing and starch digestion;
swallowing and peristalsis move food through the gut. In the stomach, acid
denatures proteins, pepsin begins their digestion, and muscular mixing forms
chyme. The stomach also limits many swallowed microbes, although it is not a
sterile barrier.

In the small intestine, bicarbonate from the pancreas neutralises much of the
acid. Pancreatic and brush-border enzymes break carbohydrates into absorbable
sugars, proteins into amino acids and small peptides, and nucleic acids into
nucleotides and related components. Bile acids made by the liver and released
from the gallbladder emulsify fat and help form micelles, allowing pancreatic
lipase products and fat-soluble vitamins to enter intestinal cells. Hormones
such as secretin and cholecystokinin, the enteric nervous system, and feedback
from nutrients coordinate secretion and motility.

Villi and microvilli enlarge the absorptive surface. Water-soluble nutrients,
many minerals, and short-chain fatty acids generally enter intestinal blood and
travel first through the portal vein to the liver, where they are processed or
stored. Most long-chain fat is packaged into chylomicrons and initially enters
lymph before reaching the bloodstream. Transporters, digestion-resistant
components, and the gut microbiota influence what is absorbed.

The large intestine absorbs additional water and electrolytes, houses microbes
that produce metabolites such as short-chain fatty acids, and compacts and
stores stool. Absorbed molecules can then be oxidised for energy, used to build
and repair tissues, or stored. Disease affecting motility, enzymes, bile,
intestinal surface, pancreas, liver, or microbiota can therefore cause
malabsorption even when food intake is adequate.

\section{How does bone marrow make new blood cells?}\index{bone marrow}\index{bones}\index{ageing}\index{amino acids}\index{anaemia}\index{cancer}

Bone marrow performs haematopoiesis: haematopoietic stem cells both self-renew
and produce progressively specialised progenitors. Broad myeloid and lymphoid
pathways give rise to red cells, platelets, granulocytes, monocytes, and
lymphocytes, although some lineage relationships are more flexible than a
simple tree suggests. Stem cells interact with stromal cells, blood vessels,
chemical signals, and the marrow niche, which help regulate when they divide,
remain quiescent, or mature.

Low oxygen stabilises hypoxia-inducible signals in the kidney, increasing
erythropoietin and stimulating erythroid precursors. They make haemoglobin,
shrink, and normally expel their nucleus before becoming red blood cells;
adequate iron, vitamin B12, folate, amino acids, and other nutrients are
needed. Thrombopoietin, produced mainly by the liver, promotes megakaryocyte
development. Megakaryocytes extend processes into marrow vessels and release
many platelets. Colony-stimulating factors and inflammatory signals adjust
production of innate immune cells, while lymphocyte development also involves
the thymus and other lymphoid tissues.

Mature cells enter the circulation and have different lifespans and functions:
red cells transport gases, white cells coordinate innate and adaptive defence,
and platelets support clot formation and repair. The spleen and other tissues
also help remove or recycle ageing blood cells. Marrow production responds to
demand, but it cannot instantly replace cells. Marrow failure, infiltration by
cancer, infection, chemotherapy, severe inflammation, or nutrient deficiency
can reduce multiple cell lines, causing combinations of anaemia, infection
risk, and bleeding. Excessive or abnormal production can also cause disease.

\section{What does the lymphatic system do?}\index{lymphatic system}\index{capillaries}\index{fever}\index{heart}\index{muscles}\index{proteins}

The lymphatic system is a network of vessels, lymph nodes, and lymphoid organs
(such as the spleen, thymus, and tonsils) that supports both fluid balance and
immune defence. As blood circulates through capillaries, a small fraction of
plasma filters into surrounding tissues. Lymphatic vessels collect this excess
interstitial fluid---now called \emph{lymph}---along with proteins and cellular
debris, and return it via larger ducts to the venous circulation near the
heart. In the small intestine, specialised lymphatic vessels called
\emph{lacteals} also absorb dietary fats and fat-soluble vitamins.

Along the way, lymph passes through bean-shaped lymph nodes packed with
macrophages, T cells, and B cells. These nodes filter particles and microbes
and serve as hubs where immune responses are organised, which is why nodes in
the neck, armpits, or groin often swell and feel tender during an infection.
Unlike blood, lymph is moved mainly by vessel smooth muscle, one-way valves,
and surrounding muscle movement rather than a central pump. Persistent, firm,
painless, or enlarging lymph nodes---especially with unexplained fever, night
sweats, or weight loss---or progressive limb swelling should be medically
assessed.

\section{How do bones stay strong, and why does bone density decline with age?}\index{bones}\index{calcium}\index{collagen}\index{hormones}\index{menopause}\index{oestrogen}

Bone is living, metabolically active tissue rather than an inert scaffold. It
combines a collagen-rich protein matrix, which provides tensile strength and
slight flexibility, with calcium-phosphate mineral crystals, which provide
compressive stiffness. Throughout life, specialised cells continuously remodel
bone: \emph{osteoclasts} break down old or micro-damaged bone, while
\emph{osteoblasts} lay down new matrix that gradually mineralises. This turnover
repairs microscopic wear and helps regulate blood calcium levels.

Remodelling responds to mechanical loading, vitamin D, calcium intake,
parathyroid hormone, and sex hormones. Peak bone mass is typically reached in
young adulthood; afterward, bone resorption slowly outpaces formation. In women,
the sharp fall in oestrogen around menopause accelerates bone loss, increasing
the risk of osteopenia and osteoporosis, though men also lose bone density
with age. Weight-bearing and resistance exercise, adequate calcium and vitamin
D, and avoiding smoking and excess alcohol help preserve bone strength. A
fracture from a minor fall, progressive height loss, or sudden back pain in
later life should be medically assessed.

\chapter*{Scientific Sources and Medical Note}
\phantomsection
\addcontentsline{toc}{chapter}{Scientific Sources and Medical Note}

This book is intended for general education. Numerical values are typical
estimates, not exact limits, and individual bodies vary. Symptoms described in
these pages are not a diagnosis; persistent, severe, sudden, or worrying
symptoms should be discussed with a qualified health professional.

The answers summarise the best-supported findings from anatomy, physiology,
neuroscience, genetics, and medicine. Where evidence is limited, variable, or
still debated, the text says so. Evolutionary explanations and numerical
estimates are especially likely to be approximate. The sources below are
starting points for checking medical information; they do not replace current
clinical guidance or advice from a qualified professional.

Useful, regularly updated sources include:

\begin{itemize}
  \item National Institutes of Health, \emph{MedlinePlus}:\newline
  \url{https://medlineplus.gov/}
  \item National Heart, Lung, and Blood Institute, \emph{How Sleep Works}:\newline
  \url{https://www.nhlbi.nih.gov/health/sleep}
  \item National Institute of Neurological Disorders and Stroke, \emph{Headache}:\newline
  \url{https://www.ninds.nih.gov/health-information/disorders/headache}
  \item National Institute of Diabetes and Digestive and Kidney Diseases:\newline
  \url{https://www.niddk.nih.gov/health-information}
  \item National Cancer Institute, \emph{Cancer Causes and Prevention}:\newline
  \url{https://www.cancer.gov/about-cancer/causes-prevention}
  \item Centers for Disease Control and Prevention, \emph{Vaccines}:\newline
  \url{https://www.cdc.gov/vaccines/}
  \item National Library of Medicine, \emph{Bookshelf}:\newline
  \url{https://www.ncbi.nlm.nih.gov/books/}
  \item Eunice Kennedy Shriver National Institute of Child Health and Human
  Development, \emph{Pregnancy and sleep position}:\newline
  \url{https://www.nichd.nih.gov/newsroom/news/092019-pregnancy-sleep-position}
  \item American College of Obstetricians and Gynecologists, \emph{Female fertility and ageing}:\newline
  \url{https://www.acog.org/}
\end{itemize}

\clearpage
\phantomsection
\addcontentsline{toc}{chapter}{Index}
\printindex

\end{document}
